\chapter{稀疏条件与鲁棒重建}

\paragraph*{本章阅读锚点}
原始 3DGS 默认有较可靠的相机位姿、较密集的多视角图像、相对清晰的训练帧和稳定静态场景。一旦视角变少、位姿缺失、输入模糊、视频压缩严重或场景中有行人等干扰物，3DGS 的显式高斯会很快过拟合训练视角：高斯可能长成针状、椭球过大、颜色贴在错误深度上，或者把瞬态物体写进静态场景。本章按“缺什么先验”来读：稀疏视角方法补几何、共视、扩散或视频先验；无 SfM 方法补位姿和初始化；退化图像方法补成像模型、事件流和低光照校正；主动重建方法补下一视角选择；干扰物方法则补语义或不一致性掩码。

\section{稀疏视角新视角合成}

稀疏视角是 3DGS 最早暴露鲁棒性问题的设置。只有 3 到 5 张图像时，原始 3DGS 的梯度密度控制会把有限训练像素解释得过好，却无法约束未观察区域。早期方法如 FSGS、CoherentGS、CoR-GS 从深度、虚拟视角和双模型一致性入手；随后 DropGaussian/DropoutGS 发现随机丢弃高斯本身就是强正则；再往后，视频扩散、点图扩散、流匹配和潜空间扩散被用来补外推与遮挡内容。



% METHOD: 2312.00451
\paragraph{FSGS（少样本高斯反池化）（2023.12）}
FSGS~\cite{2312_00451} 是 few-shot 3DGS 的早期锚点。原始 3DGS 依赖 SfM 点云初始化，再通过梯度阈值 clone/split 做 densification；在只有 3 张左右输入图像时，SfM 点云本身就稀疏，未观测区域没有足够残差触发分裂，高斯容易只贴着训练视角过拟合。FSGS 的核心改变是把 densification 从“哪里误差大就复制哪里”改成“哪里几何覆盖不足就补哪里”，也就是 proximity-guided Gaussian unpooling。

具体流程上，FSGS 在已有高斯之间建立近邻关系，寻找相邻高斯之间的代表性空旷区域，并在这些区域插入新的高斯；新高斯的颜色、尺度和位置不是随机给定，而是从邻近高斯及其观测中初始化。这样新增 primitive 的来源不再只是高梯度像素，而是稀疏点云中缺失的空间覆盖。随后方法在训练相机附近采样虚拟视角，用当前 3DGS 渲染伪图像，并引入 DPT 等单目深度先验，对真实视角和增强视角同时施加深度一致性约束。新增高斯因此既要解释已有照片，又要在附近视角下保持合理几何。相对 DropGaussian 这类正则化方法，FSGS 更像是给少样本 3DGS 补“初始骨架”，后续许多稀疏视角方法都在它的补点、伪视角或深度约束上继续扩展。
\textbf{实现细节：} 初始化质量对稀疏视角 3DGS 的最终性能具有决定性影响。稀疏点云会导致高斯无法覆盖场景空缺，进而引发严重的过拟合。本节首先详细介绍里程碑式工作 FSGS，然后按稠密初始化、致密化策略两条线梳理后续进展。FSGS（2023.12） 是该方向的里程碑式工作，其核心创新在于系统性地解决了标准 3DGS 在极稀疏点云初始化下的两大失败模式：（1）启发式位置梯度致密化无法覆盖整个场景空缺，以及（2）部分高斯趋向生长为极大体积导致过拟合训练视角。FSGS 提出三管齐下的解决方案： （1）邻近引导的高斯反池化（Proximity-guided Gaussian Unpooling）：受计算机图形学中网格细分算法的启发，该方法在优化过程中动态构建邻近图（proximity graph），连接每个高斯与其 $K$ 个最近邻（实践中 $K{=}3$），并计算邻近分数 $P_i = \frac{1}{K}\sum_{j=1}^{K} D_i^K$，即高斯 $G_i$ 到其 $K$ 近邻的平均欧氏距离。当 $P_i$ 超过阈值 $t_p$ 时，在源高斯与目标高斯连线的中点生成新高斯，其尺度与不透明度属性继承目标高斯，旋转和球谐系数则初始化为零。这一策略鼓励新致密化的高斯分布在最具代表性的位置周围，在优化过程中逐步填充观测空缺，而非仅在已有高斯附近分裂或复制。（2）伪视角增强：为缓解稀疏视角下视差不足导致的过拟合，FSGS 在训练视角间通过插值相机朝向和添加高斯位置噪声 $\mathcal{N}(0, \sigma)$ 来合成伪相机视角。这些在线采样的伪视角在 2000 次迭代后启用（确保高斯已大致表达场景结构），使 3D 高斯在优化过程中持续接收新视角的几何约束，从而降低对训练视角的过拟合风险。（3）单目深度引导的几何正则化：FSGS 利用预训练的 Dense Prediction Transformer（DPT）估计训练视角和伪视角的单目深度图，并通过 Pearson 相关系数 $\rho(\mathbf{D}_{ras}, \mathbf{D}_{est}) = \frac{\mathrm{Cov}(\mathbf{D}_{ras}, \mathbf{D}_{est})}{\sqrt{\mathrm{Var}(\mathbf{D}_{ras}) \cdot \mathrm{Var}(\mathbf{D}_{est})}}$ 衡量渲染深度与估计深度的分布差异。这种放松的约束允许对齐深度结构而不受绝对深度值尺度不一致的阻碍。
\textbf{方法信息：} 引用很多，Semantic Scholar 约 371 次。



% METHOD: 2403.19495
\paragraph{CoherentGS（相干高斯与深度初始化）（2024.03）}
CoherentGS~\cite{2403_19495} 是早期极稀疏输入的重要方法。原始 3DGS 在 3 张图像下容易生成针状、杂乱的高斯；CoherentGS 的改变是把高斯运动约束搬到 2D 图像空间，让来自同一图像局部区域的高斯保持相干。

CoherentGS 先用单目深度把每个输入像素反投影成一个初始高斯，因此高斯和 2D 像素保持可追踪对应。训练早期不让这些高斯在 3D 空间里完全自由漂移，而是在图像平面上施加相干约束：隐式卷积解码器以像素坐标和深度分组为输入，预测局部共享的 3D 偏移，使相似深度邻域的高斯一起移动；多视角约束则把所有高斯渲染出的深度图送入 total variation 损失，抑制针状深度突变；光流损失进一步把不同输入图像中互相对应的像素高斯拉到一致位置。实现上，论文先做粗对齐，再做约束优化，前段固定或限制旋转和尺度，后段才放开更多高斯属性。相对 FSGS 的反池化扩展，CoherentGS 更强调“稀疏视角下的局部结构不要散”。
\textbf{实现细节：} 在稠密初始化方面，Dust-GS（2024.09） 引入自适应深度掩码技术来增强稀疏视角下的点云初始化质量。LoopSparseGS（2024.08） 提出基于循环的渐进高斯初始化策略（PGI），利用训练过程中渲染的伪图像迭代致密化初始点云，同时通过深度对齐正则化（DAR）提供精确的几何监督。CoherentGS（2024.03） 引入结构化高斯表示，通过隐式卷积解码器和全变差损失约束高斯位置，防止其在优化过程中独立移动。TWINGS（2026.05） 利用薄板样条（Thin Plate Splines）进行非刚性变形建模，将反投影点与三角化控制点对齐以获得几何精确的初始化。
\textbf{设计动机与边界：} 相对原始 3DGS 以独立高斯和图像重建误差为主的设定，CoherentGS 把问题落在以下可实现的改动上：CoherentGS针对3DGS在极稀疏输入（3张图像）下产生「针状杂乱」表示的问题，提出正则化优化和深度初始化两大策略。（1）结构化高斯表示：为每个输入图像的每个像素分配一个高斯，在2D图像空间施加约束，避免在3D空间中直接约束困难的问题。（2）单视角约束：通过隐式卷积解码器，强制具有相似深度的像素对应的高斯在优化过程中保持相干运动。解码器将2D坐标映射到3D位置偏移量，确保局部区域的高斯运动一致。（3）多视角约束：使用全变差（Total Variation）损失确保所有高斯渲染的深度图保持平滑，增强跨视角一致性。（4）基于光流的损失：在引入相干性后，进一步约束两幅图像中对应像素的高斯位置相似性。（5）深度初始化：使用现成的单目深度估计模型获取每个视角的深度，通过尺度对齐将相对深度转换到世界坐标系，为高斯提供合理的初始位置。其设计目的不是单独增加一个模块，而是用上述表示、监督或渲染约束消除原流程在目标任务中的具体失配。提出将结构化高斯表示引入3DGS，通过在2D图像空间控制高斯运动来增强相干性。引入隐式卷积解码器的单视角约束和全变差损失的多视角约束，以及基于光流的损失函数。提出使用单目深度估计进行高斯初始化，显著改善稀疏视角重建质量。适用时还应保留论文给出的边界：依赖单目深度估计的质量；在纹理极度匮乏区域仍可能出现伪影；对相机姿态精度有一定要求；计算开销因隐式解码器而略有增加。
\textbf{方法信息：} Meta 参与了这篇论文；引用很多，Semantic Scholar 约 80 次。



% METHOD: 2405.12110
\paragraph{CoR-GS（双辐射场共正则）（2024.05）}
CoR-GS~\cite{2405_12110} 沿着另一条路线处理稀疏视角：不额外引入深度或生成先验，而是让两个 3DGS 模型互相暴露不可靠区域。它的观察是，同样的稀疏输入、不同随机初始化或 densification 过程会产生不同高斯点云；如果某个区域真实约束充分，两个模型应当收敛到相近结构，反之则会出现点云位置和渲染结果分歧。

方法同时训练两个 3DGS。第一类约束是点级 co-pruning：把两组高斯中心看成点云，用近邻匹配找对应关系，若某个高斯在另一组中找不到足够近的对应点，就把它视为由稀疏监督诱发的异常 primitive 并周期性剪掉。第二类约束是渲染级共正则：从相邻训练相机之间或附近采样伪视角，让两个模型在该视角的渲染图上通过 $L_1$ 与 D-SSIM 保持一致，同时真实训练视角仍使用常规 photometric loss。推理时不需要保留双模型，只使用其中一个高斯场。相对 FSGS 的补点思路，CoR-GS 的重点是把“不确定性”变成可优化信号；它不是假设某个外部先验正确，而是利用两个独立 3DGS 的分歧来清理浮点和抑制过拟合。
\textbf{方法信息：} 引用很多，Semantic Scholar 约 143 次。
\textbf{设计动机与边界：} 相对原始 3DGS 以独立高斯和图像重建误差为主的设定，CoR-GS 把问题落在以下可实现的改动上：CoR-GS基于一个关键发现：同时训练两个3DGS辐射场时，由于densification实现中的随机性，两者会表现出点分歧（高斯位置差异）和渲染分歧（渲染像素差异），且这两种分歧与重建质量呈负相关——分歧越大，重建质量越差。基于此，方法包含两个模块：(1) 共同剪枝(Co-pruning)：将两个3DGS辐射场视为两个点云，进行点对匹配，将没有近邻匹配点的高斯视为异常点并剪枝。设置最大距离阈值5，计算非匹配掩码，周期性执行（每5次优化/密度控制循环）；(2) 伪视角共正则化：从两个最近训练视角之间采样伪视角，分别用两个高斯集合渲染图像，最小化两者之间的颜色重建损失（L1+D-SSIM），抑制渲染分歧。最终损失函数结合训练视角的真值监督损失和伪视角的共正则化项（权重=1.0）。训练完成后保留一个辐射场用于推理。其设计目的不是单独增加一个模块，而是用上述表示、监督或渲染约束消除原流程在目标任务中的具体失配。首次提出共正则化视角改善稀疏视角3DGS，发现并量化了两个3DGS辐射场之间的点分歧和渲染分歧与重建质量的负相关关系，基于此实现了无需真值信息的重建质量评估和改善适用时还应保留论文给出的边界：训练两个辐射场增加了计算和内存开销；共同剪枝的周期和阈值需要手动设置；点匹配在极度稀疏视角下可能不稳定
% METHOD: 2408.00254
\paragraph{LoopSparseGS（循环式稀疏友好初始化）（2024.08）}
LoopSparseGS~\cite{2408_00254} 可以看作把 FSGS 的“补初始化”推进成循环式重建。FSGS 在已有高斯之间插入新 primitive，但初始点云仍主要受少量输入视角限制；LoopSparseGS 的判断是，只要当前 3DGS 已经在训练视角附近形成可用渲染，就可以把这些渲染反过来喂给 SfM，重新产生更密的初始化点云。

它的管线是“训练 $\rightarrow$ 生成伪视角 $\rightarrow$ SfM 补点 $\rightarrow$ 再训练”的多轮闭环。每轮先训练一个当前 3DGS，然后围绕每个训练相机采样多个小幅移动的伪相机，渲染伪视角图像；真实图像和伪图像一起进入 SfM 或点投影流程，得到新的 3D 点，再用扩展后的点集重新初始化或继续训练高斯场。为了避免伪深度和单目深度尺度不一致，LoopSparseGS 引入 depth-alignment regularization，把稀疏但相对可信的 SfM/projected depth 与窗口化 monocular depth 对齐。它还加入 sparse-friendly sampling，专门处理稀疏场景中超大椭球高斯覆盖大量像素、造成局部误差被单个 primitive 支配的问题。公开代码中也能看到训练、looping 和 sparse sampling 参数，这说明该方法的贡献不只是一个正则项，而是一套可重复执行的循环补点流程。
\textbf{实现细节：} 在稠密初始化方面，Dust-GS（2024.09） 引入自适应深度掩码技术来增强稀疏视角下的点云初始化质量。LoopSparseGS（2024.08） 提出基于循环的渐进高斯初始化策略（PGI），利用训练过程中渲染的伪图像迭代致密化初始点云，同时通过深度对齐正则化（DAR）提供精确的几何监督。CoherentGS（2024.03） 引入结构化高斯表示，通过隐式卷积解码器和全变差损失约束高斯位置，防止其在优化过程中独立移动。TWINGS（2026.05） 利用薄板样条（Thin Plate Splines）进行非刚性变形建模，将反投影点与三角化控制点对齐以获得几何精确的初始化。
\textbf{方法信息：} 代码：\url{https://github.com/pcl3dv/loopsparsegs}。
\textbf{设计动机与边界：} 相对原始 3DGS 以独立高斯和图像重建误差为主的设定，LoopSparseGS 把问题落在以下可实现的改动上：三个核心组件：(1)渐进式高斯初始化(PGI)：利用训练过程中渲染的伪图像迭代密集化初始点云，在每轮循环中为每个训练视角生成4个新伪视角，通过SfM获取新点云并合并到初始化点集中；(2)深度对齐正则化(DAR)：结合SfM的稀疏但可靠的深度和窗口化密集单目深度估计，对齐后提供精确的几何监督；(3)稀疏友好采样(SFS)：处理稀疏输入下超大高斯椭球体问题，通过采样策略减少像素误差。整体框架基于循环机制逐步优化。其设计目的不是单独增加一个模块，而是用上述表示、监督或渲染约束消除原流程在目标任务中的具体失配。提出基于循环机制的3DGS稀疏视角框架，通过渐进式高斯初始化策略迭代密集化初始点云，结合SfM稀疏深度和窗口化密集单目深度的对齐正则化，以及处理超大高斯椭球体的稀疏友好采样策略，在四个数据集上超越现有SOTA方法。适用时还应保留论文给出的边界：循环机制增加了训练时间；单目深度估计的尺度不一致性仍有影响；对极端稀疏视角（如2-view）仍有挑战
% METHOD: 2408.11413
\paragraph{Pano2Room（单张全景图到室内 3DGS）（2024.08）}
Pano2Room~\cite{2408_11413} 是单全景输入的特殊稀疏视角方法。公开 README 直接说明流程是把输入 panorama 转成 3DGS。它面对的问题比普通 few-shot 更极端：只有一个室内全景视图，大量遮挡背面和房间角落没有观测。

方法先做 Pano2Mesh：输入是单张 panorama 及其深度图，系统在全景图像网格上连接相邻像素，再按球面投影反投影到 3D，形成保留原始纹理的初始网格。普通边长阈值会误删远处表面，所以 Pano2Room 改用深度边缘过滤器：从深度图中检测物体轮廓，在深度突变处断开三角面，而不是按绝对边长切面。随后，系统迭代选择最需要补全的新视角，用 panoramic RGBD inpainter 生成遮挡区域的颜色和深度；新内容进入网格前要经过 geometry conflict avoidance，避免与已有表面穿插或覆盖可信区域。每轮迭代都会同时收集 3D 一致伪视角。最后的 Mesh2GS 阶段把精化网格转换为 3DGS，并用这些伪视角优化高斯场。相对一般扩散补全，Pano2Room 始终让生成内容先通过网格几何检查，再进入高斯训练。
\textbf{实现细节：} 在特定场景方面，PanoPlane（2026.05） 针对室内 360 度全景场景，提出布局锚定注意力引导机制，将扩散模型的注意力导向检测到的平面表面，仅用 3 个输入视角即可实现高质量新视角合成。SparseGS-W（2025.03） 和 MS-GS（2025.09） 分别处理不受约束的真实世界照片集中的稀疏视角问题，前者利用约束扩散先验处理遮挡和外观变化，后者引入多外观能力以适应不同时间和季节的拍摄条件。Free360（2025.03） 提出分层高斯表示来解决无界 360 度场景的极稀疏无位姿重建。Pano2Room（2024.08） 从单张全景图重建室内 3D 场景，通过迭代全景 RGBD 修复逐步收集伪新视角来训练 3DGS。
\textbf{设计动机与边界：} 相对原始 3DGS 以独立高斯和图像重建误差为主的设定，Pano2Room 把问题落在以下可实现的改动上：方法流程：(1) Pano2Mesh：将输入全景图和深度图通过三角化和反投影生成初始网格，提出深度边缘过滤器替代传统边长阈值，准确断开不同物体；(2) 迭代网格精化：搜索最优相机视角，用全景RGBD修复器生成遮挡区域的纹理和几何，通过几何冲突避免机制将新内容整合到网格中，同时收集3D一致的伪新视角；(3) 网格转3DGS：将修复后的网格转换为3DGS，用收集的伪新视角训练，最终3DGS支持照片级真实感的新视角合成。其设计目的不是单独增加一个模块，而是用上述表示、监督或渲染约束消除原流程在目标任务中的具体失配。提出Pano2Room从单张全景图自动重建高质量3D室内场景。核心流程：全景转网格，迭代RGBD修复精化处理大遮挡，收集3D一致伪新视角，最终转换为3DGS训练。创新包括深度边缘过滤器替代传统边长阈值、最优相机视角搜索策略、几何冲突避免机制。适用时还应保留论文给出的边界：全景深度图的质量影响初始重建；迭代修复的计算成本较高；对极端遮挡场景可能需要更多迭代。
\textbf{方法信息：} 代码：\url{https://github.com/trickygo/pano2room}；GitHub 关注度较高，约 151 stars。



% METHOD: 2409.14316
\paragraph{MVPGS（MVS 先验稀疏视角 3DGS）（2024.09）}
MVPGS~\cite{2409_14316} 是 MVS 先验路线的直接代表。它用 MVSFormer 估计密集、视图一致的深度图，替代 COLMAP 稀疏点初始化；这样初始高斯就覆盖更多真实表面。

随后，方法利用 MVS 几何把训练视角图像 forward warp 到未见视角，作为额外外观监督。具体来说，MVSFormer 以每个训练视角为 reference，把其他视角特征投影到 3D cost volume 中回归深度；几何一致过滤后，多视角深度被融合成初始高斯中心，比 COLMAP 稀疏点更密。给定一个源图像、源视角 MVS 深度和目标相机，MVPGS 把源图像颜色沿深度前向投到目标视角，并用反向双线性采样处理非整数像素，得到“应该在未见视角看到什么”的伪监督。训练时，这个外观约束和 3DGS 的渲染损失并行使用；同时，view-consistent geometry loss 约束高斯深度不要偏离 MVS 可信区域，单目深度正则补充 MVS 不可靠处。相对 NexusGS 的对极深度，MVPGS 更依赖学习式 MVS 网络，目标是 few-shot NVS 而不是纯初始化研究。
\textbf{方法信息：} 代码：\url{https://github.com/zezeaaa/mvpgs}。
\textbf{设计动机与边界：} 相对原始 3DGS 以独立高斯和图像重建误差为主的设定，MVPGS 把问题落在以下可实现的改动上：四项关键技术：(1) MVS先验初始化：使用MVSFormer估计密集视图一致深度图作为3DGS初始点云，比COLMAP稀疏输出提供更丰富的初始信息；(2) 前向变形外观约束：基于MVS计算的几何，将训练视角图像前向变形到未见视角，提供额外外观监督防止过拟合；(3) 视图一致几何约束：对高斯参数施加几何正则化，保持准确结构不偏离；(4) 单目深度正则化：用单目深度估计作为补充监督信号。所有约束无缝集成到3DGS优化中。其设计目的不是单独增加一个模块，而是用上述表示、监督或渲染约束消除原流程在目标任务中的具体失配。提出基于MVS先验的稀疏视角3DGS方法MVPGS，通过学习型MVS初始化、前向变形外观约束和几何正则化，在few-shot NVS上达到SOTA且保持实时渲染速度适用时还应保留论文给出的边界：MVS质量受视角数量和重叠度影响；训练时间因MVS前处理增加；极端稀疏（2-3视角）下性能下降
% METHOD: 2410.16266
\paragraph{3DGS-Enhancer（视频扩散一致性增强）（2024.10）}
3DGS-Enhancer~\cite{2410_16266} 把 3D 多视角一致性问题重写为视频时间一致性问题。给定初始稀疏 3DGS，它渲染相邻参考视角之间的低质量中间视角，再用稳定视频扩散模型增强这些帧。

它先用稀疏输入训练一个初始 3DGS，再沿两张相邻参考视角之间的相机轨迹渲染低质量中间帧；这些帧和两端真实参考视角一起组成一段“视频”。视频扩散部分以 SVD 式 latent diffusion 为骨干，但条件不是单张首帧：两端参考视角经过 CLIP 编码提供语义条件，VAE latent 提供颜色和纹理条件，扩散 U-Net 在时空注意力中恢复中间帧的高质量 latent。原始 VAE decoder 容易带来颜色漂移和时序不稳，因此方法再训练 spatial-temporal decoder，用 temporal convolution、controllable feature warping 和颜色归一化把扩散 latent、低质量渲染和参考图像融合成增强图像。最后，增强帧不会无差别加入训练，而是按视角距离和置信度选择性微调 3DGS，避免扩散幻觉反过来污染真实输入。相对 GSFixer，3DGS-Enhancer 是更早的“视频扩散修正稀疏渲染”范式。
\textbf{实现细节：} ArtifactWorld（2026.04） 建立了 3DGS 伪影的精细分类学，构建了包含 107.5K 多样化配对视频片段的训练集，通过同构预测器定位结构缺陷并引导修复。3DGS-Enhancer（2024.10） 利用 2D 视频扩散先验解决 3D 视角一致性问题，通过时空解码器整合增强特征，并微调初始 3DGS 模型。初始化改进与综合优化策略。本节讨论两条技术路线：（1）初始化改进与致密化策略；（2）综合框架与特定场景方法。
\textbf{方法信息：} 引用很多，Semantic Scholar 约 108 次。
\textbf{设计动机与边界：} 相对原始 3DGS 以独立高斯和图像重建误差为主的设定，3DGS-Enhancer 把问题落在以下可实现的改动上：方法核心是将3D视角一致性挑战重新表述为视频生成的时间一致性任务：(1)给定稀疏视角训练的初始3DGS模型，渲染低质量新视角图像；(2)利用预训练的稳定视频扩散(SVD)模型作为先验，将相邻两个参考视角和中间的低质量渲染视角组成视频序列输入视频扩散模型；(3)使用CLIP编码器提取参考视角序列的条件输入，VAE编码器提取丰富的颜色信息通过分类器自由引导策略注入；(4)通过时空解码器恢复视角一致的高质量潜在特征并解码为增强图像；(5)设计选择性微调策略将增强视角集成到3DGS微调过程。训练目标包含扩散损失和渲染损失。其设计目的不是单独增加一个模块，而是用上述表示、监督或渲染约束消除原流程在目标任务中的具体失配。将3D一致性问题重构为视频时间一致性问题，利用视频扩散模型恢复视角一致的潜在特征并微调3DGS，大幅提升稀疏视角下的渲染质量适用时还应保留论文给出的边界：依赖视频扩散模型的质量和计算开销；增强过程中可能引入扩散模型的生成伪影；对极端大视角间隔效果有限
% METHOD: 2411.00144
\paragraph{SE-GS（自集成少样本 3DGS）（2024.11）}
Self-Ensembling Gaussian Splatting~\cite{2411_00144} 借鉴教师-学生和模型集成思想。少样本下训练多个独立 3DGS 再平均确实可能更稳，但成本高，而且不同模型往往仍收敛到相似错误。SE-GS 的改变是只维护两个模型，并从其中一个模型动态扰动出多个“临时成员”，让主模型学习这些成员的平均行为。

训练中有一个 $\Delta$-model 和一个 $\Sigma$-model。$\Delta$-model 按普通 3DGS 方式拟合训练图像，但系统会在训练视角之间插值采样 pseudo views，并把 $\Delta$-model 在这些伪视角上的多次历史渲染存入 buffer。若某个像素在不同训练步渲染差异大，就被视为高不确定区域；SE-GS 只对与这些像素相关的高斯属性加入随机扰动，而不是对所有高斯盲目加噪。扰动后的 $\Delta$-model 产生多样化伪视角渲染，$\Sigma$-model 则通过 photometric self-ensembling loss 与这些扰动渲染保持一致，同时仍在真实训练视角上受监督。推理时只保留 $\Sigma$-model，因此没有集成推理成本。相对 CoR-GS 的双模型分歧剪枝，SE-GS 不把分歧直接当作异常点删除，而是把不确定区域变成有控制的扰动源，让主模型学到更平滑的泛化解。
\textbf{实现细节：} 稀疏视角下抑制过拟合可沿优化层面正则、扩散先验注入与几何约束三条主线展开，三者并非互斥而是互补。沿优化层面正则方向，正则化策略针对训练动态而非显式几何下手：自集成通过对权重或预测做时间平均、提供平滑监督 ；与之互补，深度-密度联合引导 从几何稳定性角度、频率自适应锐度正则 从优化景观角度分别抑制过拟合，二者关注点不同但目标一致。在表示层面，针对球谐系数易编码视角信息诱发过拟合的现象，部分工作将球谐替换为低阶锚点表示 ，从源头削弱表示容量；几何-查询扩散 则进一步将几何约束直接嵌入扩散采样过程，把正则与生成合二为一。
\textbf{实现补充：} 实现上还要核对输入、输出、参数、特征、初始化、训练损失、正则、渲染或推理接口，以及优化中的剪枝、分裂、采样、编码或解码步骤；这些抓手决定该方法到底改了 3DGS 管线的哪一环。
\textbf{方法信息：} Alibaba 参与了这篇论文。
\textbf{设计动机与边界：} 相对原始 3DGS 以独立高斯和图像重建误差为主的设定，Self-Ensembling Gaussian Splatting for Few-Shot Novel View Synthesis 把问题落在以下可实现的改动上：SE-GS的核心方法：(1) 双模型训练：同时训练Sigma-model（主模型）和Delta-model（扰动模型），在训练图像上监督两个模型；(2) 不确定性感知扰动：在伪视角上计算像素级渲染不确定性（基于不同训练迭代的渲染缓冲），据此对Delta-model的高斯属性添加随机噪声生成多样化的扰动模型；(3) 自集成损失：最小化Sigma-model与扰动模型渲染结果的差异，通过教师-学生范式使Sigma-model学习到多样化模型的平均表示；(4) 推理时仅使用Sigma-model，无需额外推理开销。避免了训练多个独立3DGS模型的高计算成本和缺乏多样性的问题。其设计目的不是单独增加一个模块，而是用上述表示、监督或渲染约束消除原流程在目标任务中的具体失配。提出SE-GS，首次将自集成策略引入3DGS框架解决少样本过拟合问题，通过不确定性感知扰动策略以可忽略的计算开销训练多样化的3DGS集成，在LLFF、Mip-NeRF360、DTU、MVImgNet四个数据集上超越SOTA。适用时还应保留论文给出的边界：不确定性估计的准确性影响扰动质量；伪视角的选择策略需要调优；对于极稀疏视角（如仅2-3张）效果仍有限。
% METHOD: 2411.16779
\paragraph{NovelGS（新视角去噪生成 3DGS）（2024.11）}
NovelGS~\cite{2411_16779} 面向输入视角未覆盖区域的前馈生成，位置更接近 image-to-3D / sparse-view object reconstruction，而不是每个场景从零优化的 3DGS。LGM、GS-LRM 等前馈方法通常从输入图像预测 pixel-aligned Gaussians，但像素只覆盖相机看见的表面，背面或侧面容易缺失。NovelGS 的改变是把“未见目标视角”也纳入扩散去噪过程，让模型在生成 3DGS 前先学会对这些新视角做一致去噪。

训练时，输入由两类视图组成：少量 clean condition views 和若干 noisy target views，每个视图都带 RGB 与 Plucker ray embedding。Transformer denoiser 接收这些 patch tokens，并为每个像素预测对应的高斯属性图，包括深度、旋转、尺度、不透明度和 SH 颜色；所有 condition/target 视图预测出的 pixel-aligned Gaussians 合并成一个 3DGS。监督不是只看输入视角，而是把该 3DGS 渲染到目标视角和额外视角，用 $L_2$ 与 LPIPS 约束。推理时，目标视图从纯噪声开始，denoiser 每一步输出高斯并渲染目标视图，再把渲染结果按扩散时间步更新为下一步输入，直到 timestep 0 得到最终高斯场。相对 RI3D 这种 per-scene diffusion prior，NovelGS 是训练一个大模型，让“新视角去噪”和“3DGS 参数预测”成为同一个前馈生成过程。
\textbf{方法信息：} Tencent 参与了这篇论文。
\textbf{设计动机与边界：} 相对原始 3DGS 以独立高斯和图像重建误差为主的设定，NovelGS 把问题落在以下可实现的改动上：核心创新在于将扩散去噪与3DGS生成结合：(1) 模型输入包含m张干净条件视图和n张噪声目标视图，目标视图预设为未见区域。(2) Transformer-based去噪网络接收条件视图和噪声视图的图像+Plucker射线嵌入，预测像素对齐的3DGS参数。(3) 训练时：对目标视图加噪声，网络预测3DGS，渲染多视图并用L2+LPIPS损失监督。(4) 推理时：目标视图初始化为纯噪声，迭代去噪，逐步从3DGS渲染和加噪交替进行直到timestep 0获得最终3DGS。(5) 条件机制引入timestep和clean image的嵌入，增强可控性。其设计目的不是单独增加一个模块，而是用上述表示、监督或渲染约束消除原流程在目标任务中的具体失配。提出novel-view denoising范式，通过扩散模型为未见区域生成一致的新视角内容，再从去噪后的多视角生成3DGS。该方法在生成未见区域时保持3D一致性和锐利纹理，显著超越现有image-to-3D框架。适用时还应保留论文给出的边界：推理需要多步迭代去噪，速度慢于单步前馈方法；扩散过程可能引入不确定性
% METHOD: 2412.10051
\paragraph{TSGaussian（目标特定稀疏重建）（2024.12）}
TSGaussian~\cite{2412_10051} 把稀疏视角重建限定到指定目标。普通 3DGS 会把高斯资源分给整幅图，复杂背景和遮挡会消耗大量 primitive，并把目标区域的几何监督稀释掉；在农业植株、花盆、玩偶等细长或复杂目标上，这种冗余尤其明显。TSGaussian 的改变是先用 2D 语义把目标身份稳定下来，再把语义作为可微属性写进每个高斯，最后让 clone/split/prune 围绕目标而不是整幅图工作。

流程先用 YOLOv9 检测目标框，把框作为 SAM prompt 得到目标 2D mask；多视角中同类目标可能 ID 不一致，因此再用 DEVA 等 zero-shot tracker 做跨视角身份关联，得到一致的 target identity。进入 3DGS 后，每个高斯除了位置、协方差、颜色和 opacity，还带一个低维 identity encoding。语义渲染时，这些 identity code 像颜色一样按 splatting 权重聚合到像素，再经分类头预测 2D identity mask，并与 SAM/track 得到的 mask 做 cross-entropy loss。为了避免语义只在 2D 投影上成立，方法还对 $K$ 近邻高斯施加 3D semantic regularization，使空间邻近的同一目标高斯身份一致。最后，semantic-driven control 根据 identity 属性决定哪些高斯应 clone、split 或 prune，背景高斯和错误身份高斯会被压缩，目标区域则保留更多容量；深度正则再缓解少视角目标几何退化。相对 Gaussian Grouping 的全场景分组，TSGaussian 更偏“先指定目标，再重建目标”。
\textbf{实现细节：} PointmapDiff（2025.01） 利用点图作为扩散模型的条件信号。TSGaussian（2024.12） 结合语义约束和深度先验来重建特定目标。SPC-GS（2025.03） 利用场景布局初始化和语义提示一致性约束。HBSplat（2025.09） 提出混合损失深度估计和双向翘曲虚拟视角合成。FrameTwin（2026.05） 将高斯锚定到参数化曲线上进行线框结构的稀疏视角捕获。S2D（2026.03） 通过稀疏到稠密的提升管道，利用一步扩散模型修复图像伪影。Sparse-View 3DGS in the Wild（2026.04） 处理含有干扰物的真实世界稀疏视角场景。Gaussian Scenes（2024.11） 利用深度增强的扩散先验进行无位姿稀疏视角重建。Init-to-Generalize（2025.10） 证明初始化是稀疏视角 3DGS 的决定性因素。
\textbf{方法信息：} 代码：\url{https://github.com/leon2000-ai/tsgaussian}；GitHub 约 23 stars。
\textbf{设计动机与边界：} 相对原始 3DGS 以独立高斯和图像重建误差为主的设定，TSGaussian 把问题落在以下可实现的改动上：TSGaussian包含三个核心模块：(1) 一致目标语义分割：使用YOLOv9检测目标生成边界框，作为SAM的prompt生成精确2D掩码，通过零样本跟踪器确保多视角语义身份一致性；(2) 语义约束2D到3D：为每个高斯引入Identity Encoding（身份编码）分配语义属性，通过语义渲染将3D高斯投影回2D语义掩码计算身份损失，加入3D语义正则化损失增强一致性；(3) 语义驱动高斯控制和剪枝：根据语义属性自适应克隆、分裂和剪枝高斯，将计算资源集中在指定目标区域，减少背景冗余。使用深度先验约束避免稀疏视角下的几何退化。其设计目的不是单独增加一个模块，而是用上述表示、监督或渲染约束消除原流程在目标任务中的具体失配。提出TSGaussian，首个针对特定目标的稀疏视角3DGS重建框架。结合YOLOv9检测和SAM分割生成一致的2D语义掩码，通过Identity Encoding为每个高斯分配语义属性，3D空间一致性正则化实现高斯聚类，语义驱动剪枝减少冗余。在三个标准数据集和一个新数据集上优于SOTA。适用时还应保留论文给出的边界：依赖2D检测和分割模型的准确性；对严重遮挡目标效果有限；语义粒度受YOLOv9类别限制。
% METHOD: 2501.02913
\paragraph{PointmapDiff（点图条件扩散新视角）（2025.01）}
Pointmap-Conditioned Diffusion~\cite{2501_02913} 用 point map 作为扩散条件。RGB 条件跨视角对应弱，而点图是 3D 场景坐标投影到目标视图，天然携带几何对应；因此它比普通图像条件更适合一致新视角生成。

模型采用参考 U-Net 和目标 U-Net 双流结构：参考分支读取输入 RGB 图像并提取可传递的外观特征，目标分支从噪声 latent 开始生成新视角。关键条件不是普通深度图，而是 point map，即每个像素对应的 3D 场景坐标；参考和目标 point map 经过坐标变换后告诉扩散模型哪些像素在 3D 中对应同一位置。ControlNet 把这些 point map 特征注入目标 U-Net，参考交叉视角注意力再把输入视角的纹理和语义传给目标视角。用于 3DGS 时，先训练或初始化一个街景 3DGS，再在外推相机位置渲染低质量图像；PointmapDiff 以最近参考图像和 LiDAR/深度点图为条件修复这些外推视角，生成结果加入增强训练集继续微调 3DGS，并用 LiDAR 深度损失保持几何对齐。相对 FlowR，它仍是扩散生成；相对 Splatent，它的核心条件是显式 3D 坐标图。
\textbf{实现细节：} PointmapDiff（2025.01） 利用点图作为扩散模型的条件信号。TSGaussian（2024.12） 结合语义约束和深度先验来重建特定目标。SPC-GS（2025.03） 利用场景布局初始化和语义提示一致性约束。HBSplat（2025.09） 提出混合损失深度估计和双向翘曲虚拟视角合成。FrameTwin（2026.05） 将高斯锚定到参数化曲线上进行线框结构的稀疏视角捕获。S2D（2026.03） 通过稀疏到稠密的提升管道，利用一步扩散模型修复图像伪影。Sparse-View 3DGS in the Wild（2026.04） 处理含有干扰物的真实世界稀疏视角场景。Gaussian Scenes（2024.11） 利用深度增强的扩散先验进行无位姿稀疏视角重建。Init-to-Generalize（2025.10） 证明初始化是稀疏视角 3DGS 的决定性因素。
\textbf{方法信息：} Huawei 参与了这篇论文。
\textbf{设计动机与边界：} 相对原始 3DGS 以独立高斯和图像重建误差为主的设定，Pointmap-Conditioned Diffusion for Consistent Novel View Synthesis 把问题落在以下可实现的改动上：PointmapDiff采用双流架构：参考U-Net处理输入视角图像产生特征，目标U-Net接收噪声隐变量生成新视角图像。核心创新在于使用点图（point maps，即3D场景坐标的2D投影）作为几何条件信号，通过ControlNet注入扩散模型。引入参考交叉视角注意力模块确保隐式信息从参考视角传递到生成目标视角。点图建立了比RGB更好的跨视角对应关系，且可适应稀疏LiDAR数据。生成的视角可对齐LiDAR扫描，作为额外监督信号精炼3DGS等场景表示。其设计目的不是单独增加一个模块，而是用上述表示、监督或渲染约束消除原流程在目标任务中的具体失配。提出PointmapDiff框架，利用预训练2D扩散模型进行新视角合成，使用点图（栅格化3D场景坐标）作为条件信号捕获几何和光度先验。引入参考注意力层和ControlNet实现跨视角一致性。可蒸馏到3DGS表示以改善视角外推。适用时还应保留论文给出的边界：依赖预训练扩散模型质量；点图条件需要深度或LiDAR数据；生成速度受扩散采样限制。
% METHOD: 2503.05082
\paragraph{Scene-Grounded Video Diffusion（场景接地视频先验）（2025.03）}
Taming Video Diffusion Prior with Scene-Grounding Guidance~\cite{2503_05082} 针对稀疏输入中的两个难点：视野外外推和遮挡补全。普通视频扩散模型能沿相机轨迹生成连续视角，但如果只把少量输入图像作为条件，模型很容易在未观测区域生成不存在的物体或帧间不一致的纹理；这些图像再用于训练 3DGS，会把幻觉固化成黑洞、阴影或漂浮结构。该方法的改变是：不微调视频扩散模型，而是在采样过程中用当前 3DGS 的渲染序列对扩散 latent 施加 scene-grounding guidance。

具体做法是先用稀疏输入训练一个强基线 3DGS，并把这个模型作为 grounding model。系统初始化一组相机轨迹，重点覆盖训练视角之外和容易出现遮挡的区域；对每条轨迹，3DGS 渲染 RGB 序列以及由 opacity blending 得到的 transmittance/mask，用来标记哪些像素属于已有场景的可信投影，哪些像素需要扩散模型补全。视频扩散采样时，每一步把 noisy latent 预测回 clean space，与 3DGS 渲染序列在 mask 区域计算 $L_1$ 和 perceptual guidance，其梯度被加入 denoising 过程。这样扩散模型仍保留生成未观测内容的能力，但被当前场景的几何和外观约束住，不容易生成与房间结构冲突的元素。生成序列最后和原始稀疏图像一起进入 3DGS 优化；论文实验中还使用 DUSt3R 点云与 FSGS unpooling 形成较强初始模型。相对 GSFixer 的伪影修复，它更偏“先沿轨迹补未观测区域，再用补全序列重训场景”。
\textbf{实现细节：} 在视频扩散先验的利用方面，Taming Video Diffusion Prior（2025.03） 解决了两个被忽视的实际问题：外推和遮挡。它提出了基于渲染序列的场景接地引导来驯服扩散模型生成一致的帧序列。FlowR（2025.04） 提出多视角流匹配模型，直接学习从稀疏重建渲染到密集重建渲染的映射，在 360 万图像对上训练后可一次前向传播处理多达 45 个视角。BetterScene（2026.02） 在 Stable Video Diffusion 框架中引入时间等变正则化和视觉基础模型对齐的表示，以增强 VAE 模块对 3D 一致性的建模。
\textbf{方法信息：} Huawei 参与了这篇论文。



% METHOD: 2503.10860
\paragraph{RI3D（修复与补全扩散先验）（2025.03）}
RI3D~\cite{2503_10860} 把极稀疏 NVS 明确拆成两个任务：可见区域修复和缺失区域补全。普通 3DGS 在 3-view 或 6-view 360 度场景中会把“已看见但渲染有伪影的区域”和“完全未看见的区域”混在同一个 photometric loss 中处理；生成式方法若直接补全整张图，又容易把不一致内容灌入 3D 优化。RI3D 的改变是分别个性化两个扩散模型：repair model 负责把当前 3DGS 渲染修成高质量可见区域图像，inpainting model 专门负责未观测区域。

初始化上，RI3D 先组合两类深度：DUSt3R 提供跨视角一致但偏平滑的几何，单目相对深度提供局部细节；二者通过 Poisson blending 融合成每个输入图像的深度图，然后给每个像素放置一个高斯并投影到 3D。优化分两阶段。第一阶段只处理可见区域：当前 3DGS 渲染被送入 repair diffusion，输出更干净、更锐利的图像，作为 pseudo ground truth 反过来监督高斯颜色和几何。第二阶段打开缺失区域：inpainting diffusion 在 mask 区域 hallucinate 内容，repair model 继续帮助新内容与周围可见区域融合，3DGS 则通过多视角优化把补全内容固定到一致的 3D 位置。相对 Scene-Grounded Video Diffusion 的“用 3DGS 渲染约束视频扩散采样”，RI3D 更强调任务拆分：先修复可见，再补全不可见。
\textbf{方法信息：} Meta 参与了这篇论文。
\textbf{设计动机与边界：} 相对原始 3DGS 以独立高斯和图像重建误差为主的设定，RI3D 把问题落在以下可实现的改动上：RI3D的核心贡献：(1) 将视角合成分离为两个任务，引入两个个性化扩散模型：修复模型（repair）以渲染图像为输入预测高质量图像作为伪GT约束优化，修复模型（inpainting）专门负责幻觉缺失区域细节；(2) 两阶段优化策略：第一阶段用修复模型重建可见区域，第二阶段用修复模型重建缺失区域同时通过进一步优化确保一致性；(3) 新的高斯初始化方法，结合3D一致且平滑的深度与高细节的相对深度获取逐图像深度估计。两个扩散模型都在输入图像上个性化调优。其设计目的不是单独增加一个模块，而是用上述表示、监督或渲染约束消除原流程在目标任务中的具体失配。首次将稀疏视角NVS明确分离为可见区域重建和缺失区域幻觉两个任务，分别使用个性化修复和修复扩散模型，结合新深度初始化方法，在极稀疏输入下超越SOTA。适用时还应保留论文给出的边界：两阶段优化增加训练时间；对扩散模型个性化质量有依赖；360度场景的极端稀疏输入仍具挑战。
% METHOD: 2503.18794
\paragraph{NexusGS（对极深度纽带）（2025.03）}
NexusGS~\cite{2503_18794} 通过对极几何直接补深度。它与 FSGS 的差别在于，FSGS 在已有高斯邻域之间插点，仍会受原始点云分布影响；NexusGS 则试图在训练开始前就生成更密、更准的点云，减少 3DGS adaptive density control 在稀疏监督下随机 split 的空间。它也不同于 DNGaussian 这类单目深度正则，因为它把相机位姿和跨视角光流放回到三角化几何里，得到带 metric grounding 的像素深度。

具体步骤是三段。Epipolar Depth Nexus 先用 FlowFormer++ 等光流模型估计源视角像素到目标视角的匹配点；由于光流点未必正好落在对极线上，方法把它投影到目标视角对极线的最近点，再用两相机内外参显式求该源像素深度。一个源像素会从多个目标视角得到多个深度候选，Flow-Resilient Depth Blending 不是简单平均，而是估计匹配点沿对极方向扰动时深度变化率，优先选择或加权那些对光流误差不敏感的深度。最后，Flow-Filtered Depth Pruning 用跨视角流一致性过滤不可靠深度点。保留下来的深度反投影成密集点云，用来初始化 3DGS，使后续优化更多是在可靠几何上细化颜色和透明度，而不是从极稀疏点云中猜位置。相对 HBSplat，NexusGS 的管线更集中：核心就是把对极深度做准，再把深度嵌入初始点云。
\textbf{实现细节：} NexusGS（2025.03） 利用核线几何（epipolar geometry）将深度信息直接嵌入点云，通过核线深度连接、流鲁棒深度混合和流过滤深度剪枝三步构建稠密初始点云。GBR（2024.12） 集成神经束调整模块和生成式深度精修模块，仅用 4--6 张输入视角即可实现高保真的高斯泼溅和网格重建。CuriGS（2025.11） 引入课程学习框架，在训练视角周围采样不同扰动级别的学生视角，渐进解锁更高扰动级别以稳定增强训练。S2Gaussian（2025.03） 同时处理稀疏视角和低分辨率的复合退化问题，通过两阶段策略先优化低分辨率高斯表示再精修高分辨率版本。
\textbf{方法信息：} 代码：\url{https://github.com/usmizuki/nexusgs}；GitHub 约 73 stars。
\textbf{设计动机与边界：} 相对原始 3DGS 以独立高斯和图像重建误差为主的设定，NexusGS 把问题落在以下可实现的改动上：NexusGS包含三个关键步骤：(1)对极深度纽带(Epipolar Depth Nexus)：利用光流估计和相机位姿计算像素级深度值，通过将对极线上的流估计点投影为最近点来缓解光流误差；(2)流鲁棒深度混合(Flow-Resilient Depth Blending)：对来自不同视角的深度值进行鲁棒融合，使用光流置信度作为权重，有效减轻光流估计不准确性带来的深度误差；(3)流过滤深度剪枝(Flow-Filtered Depth Pruning)：基于光流一致性过滤不可靠的深度点，确保初始点云的可靠性。整体策略避免了复杂的正则化设计，通过直接将深度信息嵌入点云来减少高斯分裂操作的随机性。其设计目的不是单独增加一个模块，而是用上述表示、监督或渲染约束消除原流程在目标任务中的具体失配。提出基于对极几何深度先验的密集点云初始化策略，通过光流和相机位姿计算精确深度图，有效解决稀疏视角下3DGS的过平滑和过拟合问题。适用时还应保留论文给出的边界：依赖光流估计的质量；在纹理极度缺失的区域深度估计困难；计算开销主要来自光流估计步骤。
% METHOD: 2503.20998
\paragraph{CoMapGS（共可见性图监督）（2025.03）}
CoMapGS~\cite{2503_20998} 认为稀疏视角的监督不应均匀对待每个像素。某些区域被多个视角共同观察，几何更可信；某些区域几乎单视角可见，最容易歧义。它的改变是构建 per-pixel covisibility map。

共可见性图来自密集对应模型而不是人工标注。CoMapGS 用 MASt3R 等模型在训练图之间建立 dense correspondences，对每个像素累计它能在多少其它视角中找到可靠对应；高计数区域说明几何被多视角约束，低计数或 mono-view 区域说明该处容易被 3DGS 误解释。这个分数随后控制两件事。第一是增强初始点云：方法利用已优化的 COLMAP 相机，把单目深度按现有点云尺度对齐，再在高共视区域和低共视区域都补充 COLMAP 缺失的点，避免高斯只能围绕稀疏关键点生长。第二是自适应监督：渲染损失和 proximity loss 按 covisibility score 加权，低共视区域权重更高；一个 proximity MLP classifier 预测高斯是否接近原始/增强几何，用来约束 mono-view 区域不要飘离合理表面。相对 MVPGS 的“用 MVS 直接提供深度和外观伪监督”，CoMapGS 更关注监督分配：哪些像素可信、哪些像素欠约束，应当在 loss 中被区别对待。
\textbf{方法信息：} Huawei 参与了这篇论文。
\textbf{设计动机与边界：} 相对原始 3DGS 以独立高斯和图像重建误差为主的设定，CoMapGS 把问题落在以下可实现的改动上：CoMapGS的核心贡献包括三个部分。（1）共可见性图作为核心组件：重新定义新视角合成问题，利用每图像每像素的共可见性计数（共可见性图）来识别区域特定的不确定性，取代传统的均匀监督方式。（2）增强初始点云：针对低和高不确定性区域均增强COLMAP衍生的稀疏点云，补偿特征匹配生成的稀疏点云缺乏几何细节的问题。这种增强也使其他少样本3DGS方法受益。（3）自适应监督：通过基于共可见性分数的加权和近邻分类器实现自适应监督。近邻分类器用于正则化高不确定性（单视角）区域的合理几何。共可见性分数衍生出的稀疏度分数在不同稀疏度场景中实现一致的性能提升。其设计目的不是单独增加一个模块，而是用上述表示、监督或渲染约束消除原流程在目标任务中的具体失配。提出了基于共可见性图的CoMapGS方法，通过构建共可见性图解决稀疏视角下区域特定的不确定性问题，增强初始点云并应用不确定性感知的加权监督，在Mip-NeRF 360和LLFF数据集上超越SOTA。适用时还应保留论文给出的边界：共可见性图计算增加预处理开销；对极端稀疏（仅2-3张图像）的效果可能有限。
% METHOD: 2504.00773
\paragraph{DropGaussian（随机丢弃高斯的结构正则）（2025.04）}
DropGaussian~\cite{2504_00773} 是少视角 3DGS 中“只改训练机制、不引入外部先验”的锚点。原始 3DGS 在 3-view 这类设置下会把有限训练图像拟合得很紧：靠近相机或前景的高斯遮住后方高斯，后方高斯的透射率 $T_i$ 很低，渲染公式中的权重也低，于是几乎拿不到足够梯度；训练后期这种不均衡会进一步变成过拟合，前景或训练视角附近长出许多只服务当前视角的高斯。DropGaussian 的改变，是把深度学习里的 dropout 思想放到高斯集合上，让每次训练迭代只使用一个随机子集渲染。

具体流程很简单，但作用点明确。训练时先给定当前丢弃率 $r(t)$，随机把一部分高斯的贡献置零；剩余高斯的不透明度贡献用 $1/(1-r)$ 做补偿，使整体亮度不因随机删除而系统性变暗。这样一来，原来被前方高斯遮挡的远处高斯会在某些迭代里变得可见，获得更大的颜色和几何梯度，参与 densification 和属性更新。$r(t)$ 不是固定常数，而是随迭代逐步增加：早期让大部分高斯参与，先建立场景粗结构；后期过拟合更严重时再加强随机扰动。论文和代码都强调这是一种 plug-and-play 修改：不需要单目深度、MVS、扩散模型，也不改变测试阶段，推理时仍渲染全部高斯。

相对 FSGS 的“补点和虚拟视角”路线，DropGaussian 不补外部几何，而是让已有高斯获得更公平的优化机会；相对 Co-Adaptation 和 DropoutGS，后者可以看作在它的随机扰动思想上继续做诊断指标或边缘恢复。它之所以适合作为锚点，是因为它把少视角过拟合从“缺先验”解释为“显式高斯之间可见性和梯度分配不均”，后续许多轻量正则方法都沿着这个角度展开。
\textbf{实现细节：} DropGaussian（2025.04） 的核心思想受启发于神经网络中的 dropout（一种通过随机丢弃网络单元来防止过拟合的正则化技术）：在每次训练迭代中，随机丢弃一部分高斯基元，使其暂时退出渲染过程。直觉上来说，当前方遮挡基元被丢弃后，原本被遮挡的远处基元突然变得"可见"，从而获得更丰富的梯度信号，有机会更新自身属性。这一操作完全不需要引入外部先验模块，仅通过训练阶段的随机扰动即可显著缓解过拟合。具体实现上，DropGaussian 在每次迭代时以丢弃率 $r$ 随机选择高斯基元予以移除。例如 $r = 0.1$ 表示丢弃 10\% 的基元。为了维持渲染亮度的一致性，被丢弃基元的贡献通过补偿因子分配给保留基元：保留基元的不透明度乘以 $M(i) = 1/(1-r)$，被丢弃基元的 $M(i) = 0$。这样即使在丢弃过程中，每个像素的总颜色贡献量仍然保持不变，渲染结果不会因丢弃操作而变暗或变亮。更进一步，DropGaussian 观察到一个重要现象：随着训练推进，过拟合倾向会逐步加剧。因此丢弃率不应该是固定值，而应随训练迭代线性增长。其公式为： 其中 $\lambda \in (0,1)$ 是缩放因子（默认 $\lambda = 0.2$），$t$ 是当前迭代，$t_{\mathrm{total}}$ 是总迭代数。这一渐进策略在训练初期施加轻微正则化以避免破坏几何建立，后期则加大正则化力度以对抗日益严重的过拟合。此外，随机丢弃优于基于梯度或距离的选择性丢弃，因为后者存在反复丢弃关键重建基元的风险。DropGaussian 仅临时停用基元而非永久移除，被停用的基元在后续迭代中可以重新参与优化。
\textbf{实现补充：} 实现上还要核对输入、输出、参数、特征、初始化、训练损失、正则、渲染或推理接口，以及优化中的剪枝、分裂、采样、编码或解码步骤；这些抓手决定该方法到底改了 3DGS 管线的哪一环。
\textbf{锚点细化：} DropGaussian：针对 3DGS 远处基元因前方遮挡、透射率极低而收不到有效梯度、最终只在近处区域堆叠参数的过拟合温床，用``随机丢弃高斯基元 + 渐进式丢弃率''替代固定优化，核心思想是丢弃前方遮挡基元后原本被遮挡的远处基元突然``可见''、获得梯度信号更新属性，丢弃率随训练迭代线性增长以匹配日益严重的过拟合倾向，被丢弃基元贡献通过 $1/(1-r)$ 补偿因子分配给保留基元、维持亮度一致。
\textbf{设计动机与边界：} 相对原始 3DGS 以独立高斯和图像重建误差为主的设定，DropGaussian 把问题落在以下可实现的改动上：DropGaussian的核心思想来自深度学习中的Dropout正则化。在稀疏视角条件下，远离相机的高斯的透射率T\_i很低，且由于输入视角有限，这些高斯很少被观察到，导致其属性（缩放、颜色、不透明度等）无法被有效更新，进而降低整体渲染质量并加剧过拟合。DropGaussian在每次训练迭代中，以概率p随机移除（丢弃）一部分高斯，使剩余高斯的透射率提高，获得更大的梯度反馈，从而更积极地参与渲染优化。关键设计：丢弃比率p随训练进程逐步增加（progressive dropout ratio），从低比率开始逐步提高。这是因为训练早期需要所有高斯参与以建立基本场景结构，后期过拟合风险增大时需要更强的正则化。DropGaussian可无缝集成到原始3DGS框架的优化循环中，不引入额外计算开销、外部深度估计或生成模型。其设计目的不是单独增加一个模块，而是用上述表示、监督或渲染约束消除原流程在目标任务中的具体失配。提出DropGaussian，一种无需外部先验的稀疏视角3DGS正则化方法。受dropout启发，在训练过程中随机移除部分高斯，使未被移除的高斯获得更大梯度并增加可见性，从而更积极地参与优化过程。丢弃比率随训练进行逐步增加，更好缓解后期过拟合问题。仅需对原始3DGS框架做简单修改，无额外计算开销即可达到与先验方法竞争的性能。适用时还应保留论文给出的边界：丢弃比率需要调参；对极稀疏视角（如仅两视角）的改善有限；丢弃操作可能影响细节恢复。
\textbf{方法信息：} 代码：\url{https://github.com/dcvl-3d/dropgaussian_release}；GitHub 关注度较高，约 144 stars；引用很多，Semantic Scholar 约 58 次。



% METHOD: 2504.01647
\paragraph{FlowR（从稀疏渲染到真实图像的流匹配）（2025.04）}
FlowR~\cite{2504_01647} 也是用生成模型补稀疏视角，但它与视频扩散路线的关键差异在于 source distribution。扩散式方法通常从噪声出发，再用输入图像或相机条件约束生成；FlowR 直接把初始 3DGS 在目标视角的低质量渲染当作源分布，学习一条从“错误或不完整渲染”流向“同视角真实图像”的 velocity field。这样如果某个视角已经被初始重建解释得很好，模型可以学会接近恒等映射；如果某个视角有 floaters、模糊或几何缺口，模型才进行较大修正。

完整流程分两层。第一层是鲁棒初始重建：系统用 MASt3R 在稀疏共可见图的边上提取匹配，利用匹配传递性建立 tracks，再用 COLMAP 三角化 3D 点；为了得到度量尺度，还把三角化深度与 Metric3D 等单目深度估计做尺度对齐。基于这批点训练一个短 schedule 的 3DGS，作为低质量但几何对齐的初始渲染器。第二层是数据密集化：FlowR 构造数百万“初始渲染--真实图像”对，训练多视角 Transformer 流匹配模型，一次前向可以同时增强多个新视角。推理时，初始 3DGS 先渲染候选新视角，FlowR 把这些渲染增强为更接近密集采集质量的图像，再把增强图像作为额外训练视图重新优化 3DGS。相对 scene-grounded diffusion，FlowR 的锚点不是“扩散采样时加场景约束”，而是“直接学习稀疏渲染到真实图像的分布映射”。
\textbf{实现细节：} 在视频扩散先验的利用方面，Taming Video Diffusion Prior（2025.03） 解决了两个被忽视的实际问题：外推和遮挡。它提出了基于渲染序列的场景接地引导来驯服扩散模型生成一致的帧序列。FlowR（2025.04） 提出多视角流匹配模型，直接学习从稀疏重建渲染到密集重建渲染的映射，在 360 万图像对上训练后可一次前向传播处理多达 45 个视角。BetterScene（2026.02） 在 Stable Video Diffusion 框架中引入时间等变正则化和视觉基础模型对齐的表示，以增强 VAE 模块对 3D 一致性的建模。
\textbf{实现补充：} 实现上还要核对输入、输出、参数、特征、初始化、训练损失、正则、渲染或推理接口，以及优化中的剪枝、分裂、采样、编码或解码步骤；这些抓手决定该方法到底改了 3DGS 管线的哪一环。
\textbf{方法信息：} Meta 参与了这篇论文。



% METHOD: 2504.09491
\paragraph{DropoutGS（随机丢弃与边缘引导分裂）（2025.04）}
DropoutGS~\cite{2504_09491} 与 DropGaussian 同属 dropout 正则路线，但它从过参数化角度重新解释稀疏视角失败：训练图像越少，过多小尺度高斯越容易记住训练视角细节，而不是形成稳定 3D 结构。论文的 pilot study 显示，在 3-view 条件下，较少 primitive 的模型虽然细节更糊，却更不容易出现空洞和严重过拟合；因此 DropoutGS 先用随机丢弃降低有效模型复杂度，再用边缘引导分裂补回高频结构。

训练时，Random Dropout Regularization 以概率 $p$ 随机停用一部分高斯，只用剩余 primitive 渲染当前视角；但它不是简单拿 dropout 子模型去直接拟合图像，而是让 dropout 渲染与完整模型渲染保持一致，从而把完整模型当作 teacher，迫使任意子集都能形成合理解释。推理时所有高斯重新启用。这个正则会抑制少数高斯对训练像素的过度记忆，但也可能牺牲边缘纹理，于是 Edge-guided Splitting Strategy 用 Sobel 边缘图计算 edge score，在边缘响应强且高斯尺度满足条件的位置更积极地 split。相对 DropGaussian 的 progressive dropout ratio，DropoutGS 更像“复杂度正则 + 边缘细化”的组合插件；它可以接到原始 3DGS、FSGS 或 CoR-GS 后面，用较小改动改善稀疏视角泛化。
\textbf{实现细节：} 在 DropGaussian 的基础上，一系列后续工作从不同角度改进了 dropout 策略。DropoutGS（2025.04） 发现模型中高斯基元数量越少，过拟合反而越轻，因此提出随机 dropout 正则化（RDR，Random Dropout Regularization）结合边缘引导的分裂策略（ESS，Edge-guided Splitting Strategy），在减少基元数量的同时保留高频细节。
\textbf{方法信息：} 引用很多，Semantic Scholar 约 35 次。
\textbf{设计动机与边界：} 相对原始 3DGS 以独立高斯和图像重建误差为主的设定，DropoutGS 把问题落在以下可实现的改动上：DropoutGS包含两个核心模块：(1) 随机丢弃正则化（RDR）：受观察到较少高斯基元反而减轻过拟合的启发，在训练过程中随机丢弃部分高斯基元，降低模型复杂度以缓解过拟合。该操作类似于神经网络中的Dropout机制，但对3DGS的离散基元进行操作。(2) 边缘引导分裂策略（ESS）：为补偿Dropout导致的高频细节损失，ESS利用边缘检测引导高斯基元在边缘区域更集中地分裂，从而恢复纹理细节。两个模块构成即插即用方案，可直接集成到现有3DGS方法中。损失函数结合渲染损失和正则化项进行端到端优化。其设计目的不是单独增加一个模块，而是用上述表示、监督或渲染约束消除原流程在目标任务中的具体失配。提出随机丢弃正则化（RDR）和边缘引导分裂策略（ESS），作为即插即用模块缓解3DGS在稀疏视角下的过拟合问题。RDR通过随机丢弃高斯基元降低模型复杂度，ESS通过边缘引导恢复高频细节。在Blender、LLFF和DTU数据集上达到SOTA性能。适用时还应保留论文给出的边界：依赖Dropout比例超参数调优；对于极端稀疏（如2-view）场景仍有提升空间；仅作为后处理模块，未从根本上改进3DGS的表征能力。
% METHOD: 2504.14548
\paragraph{VGNC（验证集引导高斯数量控制）（2025.04）}
VGNC~\cite{2504_14548} 从过拟合监控角度处理稀疏视角。原始 3DGS 的 densification 倾向于持续增加高斯数，训练视角质量会随之提高，但少视角下测试视角可能在某个点后下降；也就是说，高斯数本身成了过拟合开关。VGNC 的改变是人为构造一个“验证集”，用它在训练中判断什么时候该增长、什么时候该丢弃冗余高斯。

具体流程先用 ViewCrafter 这类生成式 NVS 模型，在输入视角之间生成候选验证图；例如每对输入图之间生成一段中间视角序列。由于生成图可能有形变和幻觉，VGNC 先通过几何一致性过滤掉明显不可信的图像。为了避免生成图反过来污染相机估计，方法先只用原始输入估计相机，然后固定这些相机参数，把过滤后的验证图加入 COLMAP/初始化流程以补全初始点云。训练过程中，validation-guided monitor 定期在验证图上评估渲染质量：如果训练视角继续变好而验证视角不再受益，就说明高斯数已越过最佳点。增长-丢弃机制据此先允许高斯生长补结构，再在指定阶段删除冗余或导致验证质量下降的高斯。相对 DropGaussian 的随机 dropout，VGNC 更像给 3DGS 加上早停和模型容量选择机制，只是容量单位从神经网络参数变成了高斯数量。
\textbf{实现细节：} 在致密化策略方面，AD-GS（2025.09） 发现不受控的致密化是导致稀疏视角伪影的关键因素，提出交替致密化框架，在高致密化阶段（激进添加基元）和低致密化阶段（激进剪枝+几何正则化）之间交替。SDI-GS（2025.09） 利用分割驱动初始化，仅保留结构显著区域的点来减少高斯数量。VGNC（2025.04） 首次利用生成式新视角合成模型产生验证图像，据此确定最优高斯数量以减轻过拟合。
\textbf{方法信息：} Intel 参与了这篇论文。
\textbf{设计动机与边界：} 相对原始 3DGS 以独立高斯和图像重建误差为主的设定，VGNC 把问题落在以下可实现的改动上：VGNC的核心思路是利用验证集指导3DGS自动确定最优高斯数量。具体流程：(1)使用生成式NVS模型从稀疏输入视图生成新视角图像作为验证集；(2)设计过滤策略，基于几何一致性消除扭曲严重的生成图像；(3)联合初始化：使用过滤后的验证图像和原始输入图像进行初始化，获得更好的全局几何结构；(4)验证引导监控器：在训练过程中使用验证集监控过拟合；(5)增长-丢弃机制：动态调整高斯数量，先增长后逐步丢弃冗余高斯。这种方法不仅减少过拟合，还降低高斯数量从而减少存储和加速训练渲染。其设计目的不是单独增加一个模块，而是用上述表示、监督或渲染约束消除原流程在目标任务中的具体失配。首次将验证集引入3DGS训练过程，提出基于生成式NVS模型的验证图像生成和过滤策略，以及增长-丢弃机制动态调整高斯数量，有效减少稀疏视角3DGS的过拟合问题，同时降低存储需求并加速训练渲染。适用时还应保留论文给出的边界：依赖生成式NVS模型的质量；过滤策略可能误删有用的验证图像；额外的验证步骤增加了训练时间。
% METHOD: 2504.20378
\paragraph{Sparse2DGS（稀疏视角表面重建）（2025.04）}
Sparse2DGS~\cite{2504_20378} 把稀疏视角问题从“新视角图像是否好看”推进到“表面是否准确”。原始 3DGS 或 2DGS 在稀疏输入下通常从 COLMAP 的少量点出发，后续渲染损失又容易被颜色拟合主导，结果是图像可能局部合理，但面片位置、法向和尺度并不稳定。Sparse2DGS 的改变是把 2DGS 的平面高斯作为表面 primitive，同时把 MVS 产生的密集几何作为初始化和训练监督，让几何优先于外观过拟合。

具体流程先用 CLMVSNet 预测每个输入视角的密集深度，把深度反投影并融合成点云，用它初始化 2D Gaussian disks，而不是依赖 3-view COLMAP 点。由于 MVS 网络还提供特征图，Sparse2DGS 将每个 MVS 点对应的 RGB 和几何特征采样到高斯上，并在渲染时同时渲染颜色和 feature map；训练中用 feature cosine loss 约束渲染特征与 MVS 特征一致，同时固定颜色/特征值，避免高斯通过改外观掩盖几何错误。然后，direct Gaussian primitive regularization 用重参数化的 disk sampling 把高斯的位置、方向和尺度转成可做跨视图一致性的点约束。最后，selective Gaussian update 替换普通 2DGS 的密度控制，只在渲染几何线索可信的位置更新或增强高斯。最终表面通过 TSDF fusion 从渲染深度中提取。相对 NexusGS 这类点云初始化方法，Sparse2DGS 不只补初始点，而是让 MVS 几何持续参与 2DGS 表面优化。
\textbf{方法信息：} 代码：\url{https://github.com/wuuu3511/sparse2dgs}；GitHub 约 67 stars。



% METHOD: 2508.09667
\paragraph{GSFixer（参考引导视频扩散修复）（2025.08）}
GSFixer~\cite{2508_09667} 是“先修复 3DGS 渲染伪影，再蒸馏回 3DGS”的生成式锚点。它的基线不是在 3DGS 内部再加一个深度或稀疏正则，而是承认少视角 3DGS 的新视角渲染会出现浮动结构、几何扭曲和纹理破碎，然后把这些渲染帧当成待修复视频。与普通视频扩散不同，GSFixer 的修复必须服从输入参考视角，否则生成内容虽然好看，却不能作为 3DGS 的监督。

训练数据先从 DL3DV 等多视角数据中构造：用 3、6 或 9 张输入训练一个低质量 3DGS，沿原始相机轨迹渲染 artifact frames，并用对应的真实帧作为 clean target，形成 DL3DV-Res 这类“3DGS 伪影视频修复”数据。修复网络基于 CogVideoX/DiT 式视频扩散骨干。输入不是单张图，而是一段由 3DGS 渲染出的连续伪影帧；两端或邻近参考视角提供条件。DINOv2 从参考图像提取 2D semantic tokens，用来保持物体身份、材质和语义细节；VGGT 从参考视角提取 3D geometric tokens，用来约束相机间几何一致性。这些 token 经过投影和融合后通过 cross-attention 注入 DiT，使去噪过程知道“修成什么”和“在 3D 上应当怎样对齐”。

推理和重建时，GSFixer 先用当前 3DGS 沿 reference-guided trajectory 渲染一串待修复新视角。轨迹不是单纯插值，也不是大圆环随机绕一圈，而是在参考视角附近保证修复质量，同时沿球面路径增加角度覆盖。视频扩散模型输出修复后的帧后，这些帧被加入监督集，和原始稀疏输入一起继续优化 3DGS；优化损失由真实输入视角的重建项和修复新视角的生成监督项组成。相对 3DGS-Enhancer，GSFixer 的关键增量是把 2D 语义和 3D 几何两类参考 token 都放进视频修复；相对 GenFusion/Difix3D+，它更强调修复帧与原始参考视角之间的身份和空间一致性。
\textbf{实现细节：} 在扩散模型的高效利用方面，RealisticDreamer（2025.11） 提出引导分数蒸馏（GSD，Guided Score Distillation），结合深度翘曲引导和语义特征引导来校正视频扩散模型的噪声预测方向，确保分数更新与正确的相机姿态和几何一致。FixingGS（2025.09） 提出一种无需训练的方法，直接利用现有扩散模型为稀疏视角 3DGS 提供更准确的跨视角一致先验。GSFixer（2025.08） 构建了基于 DiT（Diffusion Transformer）的视频扩散修复模型，以参考视角的 2D 语义特征和 3D 几何特征为条件，修复 3DGS 渲染的伪影帧。
\textbf{设计动机与边界：} 相对原始 3DGS 以独立高斯和图像重建误差为主的设定，GSFixer 把问题落在以下可实现的改动上：GSFixer框架包含：(1) 参考引导视频修复模型：基于DiT架构的视频扩散模型，在有伪影的3DGS渲染帧和干净帧的配对数据上训练。以输入稀疏视角作为参考，提取2D语义标记（来自视觉编码器Ev）和3D几何标记（来自空间编码器Es），注入扩散过程引导修复。(2) 2D语义特征：确保修复内容与参考视图的语义一致；(3) 3D几何特征：确保修复内容在3D空间中的几何一致性；(4) 蒸馏回3DGS：将修复后的图像蒸馏回3DGS表示以提升其质量。构建DL3DV-Res基准数据集，包含用低质量3DGS渲染的伪影帧。其设计目的不是单独增加一个模块，而是用上述表示、监督或渲染约束消除原流程在目标任务中的具体失配。提出GSFixer框架，基于参考引导的视频修复模型改善稀疏视图3DGS质量。利用DiT视频扩散模型，提取参考视角的2D语义特征和3D几何特征引导修复，保持语义一致性和3D一致性。构建DL3DV-Res基准数据集用于3DGS伪影修复评估。适用时还应保留论文给出的边界：修复质量受视频扩散模型能力限制；处理大范围缺失区域可能产生幻觉内容；推理速度受扩散模型迭代限制。
\textbf{方法信息：} 代码：\url{https://github.com/gvclab/gsfixer}；GitHub 关注度较高，约 167 stars；引用很多，Semantic Scholar 约 24 次。



% METHOD: 2508.12720
\paragraph{Co-Adaptation 3DGS（共适应量化与缓解）（2025.08）}
Quantifying and Alleviating Co-Adaptation in Sparse-View 3D Gaussian Splatting~\cite{2508_12720} 把稀疏视角伪影解释为高斯之间的共适应。多个高斯在有限训练像素上共同拟合一个局部外观，但这种耦合未必对应真实 3D 结构；到了新视角，颜色和不透明度的协同关系失效，就出现漂浮和纹理错位。

它先定义 co-adaptation score：对同一个目标视角反复随机丢弃不同高斯子集并渲染，再计算像素级方差；如果少数高斯组合一变，颜色就大幅波动，说明模型不是学到了稳定外观，而是在训练视角上用高斯之间的偶然协同“拼”出颜色。基于这个诊断，方法给出两类训练期扰动。随机 Gaussian dropout 直接切断一部分高斯贡献，迫使剩余高斯独立解释像素；opacity 乘性噪声只扰动每个高斯在 alpha blending 中的贡献强度，比扰动位置或颜色更不容易破坏几何收敛。代码 README 也把这两项作为可插入 3DGS 渲染器和训练循环的改动：dropout mask 作用在 means、SH、opacity、scale、rotation 等参与渲染的张量上，opacity noise 则作用在透明度分支。相对 DropGaussian，它多了一个可诊断指标；相对 CoR-GS 的双模型分歧，它从同一模型内部的高斯耦合解释过拟合。
\textbf{设计动机与边界：} 相对原始 3DGS 以独立高斯和图像重建误差为主的设定，Quantifying and Alleviating Co-Adaptation in Sparse-View 3D Gaussian Splatting 把问题落在以下可实现的改动上：该方法系统研究了稀疏视角3DGS中的外观伪影问题。核心发现是：多个高斯通过联合优化过度纠缠以拟合有限的训练像素，忽略了场景的真实外观分布。为此提出：(1) 共适应评分(CA)：通过随机dropout子集高斯并多次渲染同一视角，计算像素级方差来量化纠缠程度。理论推导表明CA评分直接反映高斯颜色和不透明度属性间的耦合关系。(2) 两种缓解策略：随机高斯dropout（训练时随机丢弃部分高斯）和乘性噪声注入（在不透明度上添加乘性噪声）。两者均为即插即用设计，可应用于各种稀疏视角3DGS方法。其设计目的不是单独增加一个模块，而是用上述表示、监督或渲染约束消除原流程在目标任务中的具体失配。首次揭示稀疏视角3DGS中高斯间的共适应(co-adaptation)现象是导致外观伪影的根本原因，提出CA评分量化工具和两种轻量级缓解策略，为理解稀疏视角3DGS提供新视角。适用时还应保留论文给出的边界：未解决稀疏视角下的几何伪影问题；dropout率等超参数需调整；主要针对外观而非几何。
\textbf{方法信息：} Huawei 参与了这篇论文；代码：\url{https://github.com/chenkangjie1123/Co-Adaptation-of-3DGS}。



% METHOD: 2509.24893
\paragraph{HBSplat（混合损失深度与双向扭曲）（2025.09）}
HBSplat~\cite{2509_24893} 是稀疏视角中较完整的几何约束组合。它的出发点是，少量输入图像会同时造成三个问题：初始化点云覆盖不足、优化时高斯中心在 3D 空间乱漂、遮挡区域没有监督。HBSplat 因此不只加一个深度正则，而是把 dense matching、ray-constrained optimization、hybrid depth、virtual view synthesis 和 occlusion-aware reconstruction 串成一个管线。

方法先用 GIM-DKM 等 dense matcher 替代普通稀疏特征匹配，获得更密的初始对应和点云。随后高斯中心不再自由优化三维坐标，而是被参数化到相机射线 $r=o+zd$ 上，只优化深度 $z$，这样可以减少稀疏监督下的漂浮自由度。Hybrid-Loss Depth Estimation 进一步估计深度：reprojection loss 用匹配点的跨视角投影误差约束几何，point propagation loss 利用“a 与 b 匹配、b 与 c 匹配”形成的间接对应传播深度一致性，TV smoothness 抑制局部噪声。基于这些深度，bidirectional warping 从训练视角向虚拟视角和反方向同时扭曲图像/深度，并融合多视角结果，生成比单向 warp 更稳定的伪训练图。对遮挡区域，方法用深度差 mask 找出 disocclusion，再调用 inpainting 模型补背景或缺失纹理。相对 NexusGS 的对极深度初始化，HBSplat 更重：它覆盖初始化、深度估计、虚拟视角监督和遮挡补全，代价是管线更复杂。
\textbf{实现细节：} PointmapDiff（2025.01） 利用点图作为扩散模型的条件信号。TSGaussian（2024.12） 结合语义约束和深度先验来重建特定目标。SPC-GS（2025.03） 利用场景布局初始化和语义提示一致性约束。HBSplat（2025.09） 提出混合损失深度估计和双向翘曲虚拟视角合成。FrameTwin（2026.05） 将高斯锚定到参数化曲线上进行线框结构的稀疏视角捕获。S2D（2026.03） 通过稀疏到稠密的提升管道，利用一步扩散模型修复图像伪影。Sparse-View 3DGS in the Wild（2026.04） 处理含有干扰物的真实世界稀疏视角场景。Gaussian Scenes（2024.11） 利用深度增强的扩散先验进行无位姿稀疏视角重建。Init-to-Generalize（2025.10） 证明初始化是稀疏视角 3DGS 的决定性因素。
\textbf{方法信息：} 代码：\url{https://github.com/eternalland/hbsplat}；GitHub 约 13 stars。
\textbf{设计动机与边界：} 相对原始 3DGS 以独立高斯和图像重建误差为主的设定，HBSplat 把问题落在以下可实现的改动上：HBSplat包含五个紧密集成的组件：(1) 密集匹配初始化：用密集匹配网络替代COLMAP的稀疏特征匹配，产生更密集可靠的初始3D点云。(2) 射线约束优化：将高斯中心约束在相机射线上，将3D位置优化简化为1D深度调整，减少过拟合。(3) 混合损失深度估计：融合重投影损失（多视图一致性）、点传播损失（深度传播约束）和平滑损失（局部几何平滑），提供鲁棒的几何先验。(4) 双向扭曲虚拟视图合成：通过双向深度-图像扭曲和多视图融合生成高保真虚拟视图，为新视角提供强约束。(5) 遮挡感知重建：使用深度差掩码识别遮挡区域，结合学习型修复模型恢复遮挡内容。所有组件协同工作，保持实时推理能力。其设计目的不是单独增加一个模块，而是用上述表示、监督或渲染约束消除原流程在目标任务中的具体失配。提出统一框架HBSplat，三大核心贡献：(1) 混合损失深度估计模块，融合重投影、点传播和平滑约束确保多视图一致性；(2) 双向扭曲虚拟视图合成方法，通过双向深度-图像扭曲和多视图融合生成高保真虚拟视图提供更强约束；(3) 遮挡感知重建组件，利用深度差掩码和学习型修复模型恢复遮挡区域。在LLFF、Blender和DTU上达到SOTA（PSNR达21.13dB）。适用时还应保留论文给出的边界：多阶段流程增加系统复杂度。密集匹配和修复模型引入额外计算开销。在极端稀疏（2视图）场景下性能退化。
% METHOD: 2511.14633
\paragraph{SparseSurf（立体几何-纹理对齐表面重建）（2025.11）}
SparseSurf~\cite{2511_14633} 同样面向稀疏视角表面重建，但它指出了 Sparse2DGS/FatesGS 一类 flattened Gaussian 的新风险：为了贴合表面，高斯被压得更扁，表达更各向异性；在视角少时，这种强表面假设反而更容易贴错训练视角并在新视角暴露几何错误。SparseSurf 的核心是用自生成的立体监督和伪视角 feature consistency，把“渲染质量”和“表面几何”绑在一起优化。

具体做法有两层。第一层是 Stereo Geometry-Texture Alignment：对每个训练相机，沿最近邻训练视角方向生成一个带水平基线的 pseudo-stereo camera，当前高斯场渲染原视角和伪立体视角，送入预训练 stereo matching 网络得到 disparity，再转换为深度和法向；这些 stereo prior 经过一致性 mask 过滤后，用于约束 3DGS 渲染深度、渲染法向和由深度计算的法向。训练过程中系统周期性重渲染 stereo pair，因而当高斯渲染变好时，stereo depth 也会变得更可靠。第二层是 Pseudo-Feature Enhanced Geometry Consistency：方法从多视角特征网络提取 feature representation，通过 feature distillation 让高斯能够渲染伪视角特征，并同时约束训练视角和 pseudo-unseen views 的 feature consistency，缓解只在 3 个训练视角上贴表面的过拟合。相对 Sparse2DGS 的 MVS 初始化和特征监督，SparseSurf 更强调训练中不断从当前渲染生成新的立体几何约束。
\textbf{方法信息：} 代码：\url{https://github.com/miya-oi/SparseSurf}；GitHub 约 32 stars。
\textbf{设计动机与边界：} 相对原始 3DGS 以独立高斯和图像重建误差为主的设定，SparseSurf 把问题落在以下可实现的改动上：SparseSurf针对稀疏视角下3DGS表面重建的过拟合问题：(1)立体几何-纹理对齐：利用立体匹配框架桥接渲染质量和几何估计。对每个训练视角生成伪立体视角（在其最近邻方向以水平基线b偏移），渲染立体对图像输入预训练立体匹配网络（如StereoAnywhere）获取视差图，转换为深度先验。随渲染质量提升，深度先验更准确——形成良性循环。先验通过几何一致性检查过滤噪声。(2)伪特征增强几何一致性：不仅利用训练视角，还整合未视角的几何约束。从多视角一致性出发约束高斯的几何属性。使用压缩的高斯椭球（沿最小缩放因子方向）更好地贴合表面。(3)将3DGS表面表示为2D高斯面元，通过法线和距离渲染，利用射线-平面交叉确定无偏深度图。其设计目的不是单独增加一个模块，而是用上述表示、监督或渲染约束消除原流程在目标任务中的具体失配。提出SparseSurf方法，引入立体几何-纹理对齐提供度量监督桥接渲染质量和几何估计，以及伪特征增强几何一致性整合训练和未视角缓解过拟合。在DTU、BlendedMVS和Mip-NeRF360上达到SOTA。适用时还应保留论文给出的边界：依赖预训练立体匹配网络的质量；伪立体视角的基线选择需要适配场景；对极度稀疏（<3视角）的效果可能有限。
% METHOD: 2512.09923
\paragraph{Splatent（VAE 潜空间高斯与扩散细化）（2025.12）}
Splatent~\cite{2512_09923} 将 3DGS 从 RGB 空间推到扩散模型常用的 VAE latent space。这样做的吸引力是显然的：latent 分辨率低、渲染快，也能和 Stable Diffusion 类模型自然衔接。但问题在于 VAE latent 的高频分量并不天然满足多视角一致性；同一纹理从不同视角编码到 latent 后，高频通道可能相互冲突。直接优化 latent 3DGS 时，这些冲突会在跨视角监督中抵消，得到的场景只剩较低频结构，解码后细节发糊。

Splatent 的两阶段设计正是围绕这个问题。第一阶段把 3DGS 扩展成 feature-space 3DGS：输入图像先经过冻结 VAE encoder 得到 latent，每个高斯不只携带颜色，还携带 latent feature，训练目标是在 novel view 中渲染 latent 并通过冻结 decoder 还原图像。这个阶段主要负责低频几何和跨视角一致内容。第二阶段不再试图在 3D latent field 里恢复所有高频，而是把当前 3DGS 渲染出的 novel-view latent 与若干最近参考视角 latent 排成输入网格，送入带 multi-view attention 的单步扩散细化模型。扩散模型在 2D latent 空间从参考视角搬运纹理细节，同时保留 3DGS 渲染 latent 提供的几何布局。相对 PointmapDiff 的显式 3D 坐标条件，Splatent 的锚点是“冻结 VAE + latent 3DGS 负责结构 + 2D diffusion 负责高频细节”。
\textbf{实现细节：} RI3D（2025.03） 将视角合成过程分解为可见区域的"修复"（repair）和缺失区域的"补全"（inpainting）两个任务，分别训练两个个性化扩散模型，通过两阶段优化策略在可见和不可见区域均产生精细纹理。Splatent（2025.12） 在 VAE（Variational Autoencoder，变分自编码器）潜空间中操作 3DGS，通过多视角注意力机制从输入视角的 2D 恢复细粒度细节，而非在 3D 空间中重建。
\textbf{实现补充：} 实现上还要核对输入、输出、参数、特征、初始化、训练损失、正则、渲染或推理接口，以及优化中的剪枝、分裂、采样、编码或解码步骤；这些抓手决定该方法到底改了 3DGS 管线的哪一环。
\textbf{方法信息：} Amazon 参与了这篇论文。
\textbf{设计动机与边界：} 相对原始 3DGS 以独立高斯和图像重建误差为主的设定，Splatent 把问题落在以下可实现的改动上：Splatent包含两个阶段：(1) 潜空间3DGS：参照Feature-3DGS，将3DGS扩展到特征空间（使用VAE潜空间作为特征空间）。每个高斯额外参数化f\_z\^{}d潜特征值。将输入图像编码到潜空间，然后在这些潜变量上进行3DGS优化重建潜辐射场。也可通过前馈网络（如MVSplat360）直接预测潜3DGS。分析发现VAE潜空间的高频分量是视角依赖的，跨视角严重不一致，导致优化过程中高频信号冲突抵消，低频信号保留。(2) 扩散增强细化：核心创新——不在3D空间修复细粒度细节，而是在2D空间从输入视角恢复。利用带多视角注意力的潜扩散模型（参考Make-A-Video、Imagen Video的时空注意力），在去噪过程中将3DGS渲染的潜特征与源输入视角的潜变量融合。冻结VAE保持解码器质量，扩散模型通过注意力机制从参考视角补充高频信息。可与现有前馈框架（MVSplat360等）集成，持续改善细节保持。其设计目的不是单独增加一个模块，而是用上述表示、监督或渲染约束消除原流程在目标任务中的具体失配。提出Splatent，突破传统3D视角修复思路，在2D空间通过多视角注意力从输入视角恢复高频细节。在冻结VAE设置下保持预训练VAE的重建质量，同时实现忠实的高频细节恢复，在多个基准上建立VAE潜空间辐射场重建的新SOTA。适用时还应保留论文给出的边界：扩散模型推理增加渲染时间；依赖参考视角的质量和覆盖；冻结VAE可能限制某些场景的表示能力。
% METHOD: 2604.12251
\paragraph{ArtifactWorld（3DGS 伪影数据与视频生成修复）（2026.04）}
ArtifactWorld~\cite{2604_12251} 更像 GSFixer 后的系统化数据扩展。GSFixer 类方法证明了视频/扩散模型可以修复 3DGS 渲染伪影，但训练数据通常覆盖不足，模型也缺少明确的“哪里坏了”的空间提示。ArtifactWorld 的改变是先建立细粒度 3DGS artifact taxonomy，再构造 107.5K 规模的成对视频训练集，使修复模型见过更多几何破碎、漂浮、颜色退化和稀疏视角伪影。

模型结构是 homogeneous dual-model paradigm：在冻结 LTX-Video backbone 上接两个任务 LoRA。第一阶段预测器接收退化 3DGS 视频，输出 artifact heatmap/noisemap，定位结构缺陷；第二阶段修复器用同类视频生成骨干做 restoration。Artifact-Aware Triplet Fusion 在去噪的不同时间步组合三种 latent/mask：完整输入保留全局风格，空掩码释放生成能力，heatmap 掩码把注意力压到伪影区域。论文实现中早期保留较多 full latent 以吸收整体场景，中期提高 heatmap 权重定位结构缺陷，后期用 null 分支让模型专注高频纹理和边界融合。公开 README 还说明仓库提供两阶段推理、benchmark preprocessing、LoRA 权重和 AATF 辅助 latent；但大规模训练数据和训练代码仍标注为后续发布。相对 GSFixer，ArtifactWorld 的重点不是单个修复模块，而是“伪影 taxonomy + 大规模 paired videos + heatmap-guided video restoration”。
\textbf{实现细节：} ArtifactWorld（2026.04） 建立了 3DGS 伪影的精细分类学，构建了包含 107.5K 多样化配对视频片段的训练集，通过同构预测器定位结构缺陷并引导修复。3DGS-Enhancer（2024.10） 利用 2D 视频扩散先验解决 3D 视角一致性问题，通过时空解码器整合增强特征，并微调初始 3DGS 模型。初始化改进与综合优化策略。本节讨论两条技术路线：（1）初始化改进与致密化策略；（2）综合框架与特定场景方法。
\textbf{技术实现补充：} ArtifactWorld（2026.04） 建立了 3DGS 伪影的精细分类学，构建了包含 107.5K 多样化配对视频片段的训练集，通过同构预测器定位结构缺陷并引导修复。3DGS-Enhancer（2024.10） 利用 2D 视频扩散先验解决 3D 视角一致性问题，通过时空解码器整合增强特征，并微调初始 3DGS 模型。初始化改进与综合优化策略。本节讨论两条技术路线：（1）初始化改进与致密化策略；（2）综合框架与特定场景方法。
\textbf{方法信息：} 代码：\url{https://github.com/fyting/ArtifactWorld}；GitHub 约 10 stars。
\textbf{设计动机与边界：} 相对原始 3DGS 以独立高斯和图像重建误差为主的设定，ArtifactWorld 把问题落在以下可实现的改动上：ArtifactWorld包含三个核心设计：（1）伪影分类与数据构建：建立3DGS伪影的细粒度现象学分类法，构建107.5K多样化配对视频剪辑训练集增强模型鲁棒性；（2）同构双模型范式：在冻结的LTX-Video骨干上部署两个LoRA适配器（LoRA1用于伪影预测生成伪影热力图，LoRA2用于修复），使用AdamW优化器在8块A100 GPU上训练14K步，LoRA秩128；（3）Artifact-Aware Triplet Fusion机制：通过动态时间步调度在去噪过程中组合三种掩码（z\_full、z\_null、z\_heatmap），实现精确的强度引导时空修复。修复过程分为三个阶段：早期全局风格吸收（0.8 z\_full + 0.2 z\_heatmap）、中期伪影热力图引导、后期纯null掩码聚焦高频纹理。其设计目的不是单独增加一个模块，而是用上述表示、监督或渲染约束消除原流程在目标任务中的具体失配。提出ArtifactWorld框架，通过系统化数据扩展和同构双模型范式解决3DGS伪影修复问题。建立了3DGS伪影的细粒度现象学分类法，构建107.5K多样化配对视频训练集。利用同构预测器定位结构缺陷，通过Artifact-Aware Triplet Fusion机制在视频扩散骨干中实现精确的时空修复。适用时还应保留论文给出的边界：依赖大规模训练数据；视频扩散模型的推理速度可能限制实时应用；主要针对稀疏视角伪影，其他类型伪影待验证。

\section{位姿估计与无 SfM 重建}

3DGS 对相机位姿和初始化非常敏感。传统管线依赖 COLMAP/SfM，但低纹理、手术场景、长视频、模糊或稀疏视角都可能让 SfM 失败。此路线分两类：一类用 DUSt3R/MASt3R/VGGT/FlowMap/InstantSfM 这类几何基础模型或 GPU SfM 替代经典 SfM；另一类把位姿和高斯参数放进同一个训练目标中联合优化。
% METHOD: 2403.20309
\paragraph{InstantSplat（几何基础模型初始化与 GBA）（2024.03）}
InstantSplat~\cite{2403_20309} 是“几何基础模型 + 3DGS”的早期锚点。原始 3DGS 依赖 COLMAP 提供相机位姿和 SfM 点云；少视角时，COLMAP 特征匹配不稳定，点云又稀又偏，后续高斯优化只能在错误几何上补救。InstantSplat 的改变，是先用 DUSt3R/MASt3R 这类预训练多视图几何模型，从 2 到 3 张图像直接预测 pixel-aligned point maps、相机和置信度，再把这些 dense priors 转成 3DGS 初始化。

它的初始化不是简单把所有预测点拼起来。首先根据视图间 co-visibility 过滤重复或低质量区域，避免同一表面被多个视角重复初始化，也避免低置信点把高斯放到错误深度。随后建立高斯集合：点图给出中心，几何模型的置信度影响哪些点进入初始化以及后续优化步长。核心优化步骤是 Gaussian Bundle Adjustment。与原始 3DGS 固定相机不同，GBA 把相机位姿和高斯位置、尺度、旋转、不透明度、颜色放进同一个可微渲染目标，用 photometric error 同时修正场景和相机。置信度感知优化器让高置信区域小步精调，低置信区域允许更大更新，以便从粗几何里恢复可渲染结构。

相对 FlowMap 的可微视频 SfM，InstantSplat 更贴近 few-shot 静态场景：它把几何基础模型当作快速前端，然后用 3DGS 的渲染误差做自监督 bundle adjustment。相对 LongSplat，它处理的是短输入下快速重建，而不是长视频的持续建图。论文报告相比传统 SfM+3DGS 明显缩短重建时间，并在稀疏输入下大幅提升初始化质量，因此后续许多无 SfM 或 sparse-view 方法都可以把它视为“用基础模型替代 SfM”的参照点。
\textbf{锚点细化：} InstantSplat：针对 COLMAP 分钟级 SfM 速度瓶颈的局限，用``MASt3R 前馈 + 自监督端到端优化''替代 COLMAP 管道，核心思想是利用 MASt3R（DUSt3R 的增强版，额外输出像素级匹配以提供更稠密的对应）几何基础模型成对推理、回归像素对齐点图和内外参，构建完全连通图对齐到全局坐标系，配合共可见冗余消除（按置信度排序剪枝冗余点）和置信度感知高斯束调整（纯光度损失 + 动态学习率），将整个流程从分钟级缩短至秒级、实现 30 倍加速。
\textbf{方法信息：} NVIDIA 参与了这篇论文；引用很多，Semantic Scholar 约 59 次。



% METHOD: 2404.12547
\paragraph{Alternatives to SfM Initialization（初始化策略评估）（2024.04）}
Evaluating Alternatives to SFM Point Cloud Initialization~\cite{2404_12547} 是初始化问题的系统研究。它不是提出一个新的渲染器，而是把原始 3DGS 的“必须用 COLMAP 点云初始化”拆开检查：高斯是局部有界 primitive，训练中又依赖 split/clone/prune 这类离散操作，所以如果初始中心离真实表面太远，光度梯度往往不足以把粗结构从空场景中拉出来。论文的核心问题是：如果相机位姿已经有了，点云是否一定要来自 SfM。

论文先重新评估随机初始化。原始 3DGS 中“在相机包围盒附近随机撒点”的设置并不是随机初始化的上限；只要扩大覆盖范围、增加初始高斯数并调整分布，随机初始化就可以比朴素版本好很多。更重要的方案是 NeRF-guided initialization：先用 NerfAcc/INGP 类低成本 NeRF 训练很短时间，利用体渲染终止概率沿训练射线采样表面点，把这些采样点的位置和颜色作为 3DGS 初始点云；训练 3DGS 时还可以用 NeRF 预测深度做深度蒸馏，减少错误剪枝和局部空洞。这样 NeRF 不负责最终实时渲染，只提供粗几何和深度先验，最终质量仍由 3DGS 优化得到。相对原始 3DGS，这篇工作的价值在于把“点云来源”从固定 SfM 前提变成可替换模块：随机初始化、低成本体渲染采样、深度监督都可以单独评估，也为后续无外部点云、SLAM 位姿输入或大规模视频初始化方法提供了基线。\textbf{方法信息：} Google DeepMind 参与了这篇论文；项目页：\url{https://theialab.github.io/nerf-3dgs/}；代码入口对应的网页仓库为 \url{https://github.com/theialab/nerf-3dgs}，约 13 stars，该仓库主要是项目页文件，不按完整训练实现计。
\textbf{设计动机与边界：} 相对原始 3DGS 以独立高斯和图像重建误差为主的设定，Evaluating Alternatives to SFM Point Cloud Initialization for Gaussian Splatting 把问题落在以下可实现的改动上：该方法深入分析了3DGS初始化阶段的问题，指出3DGS对SfM的依赖主要源于其离散的局部化剪枝和增长操作难以有效捕获场景粗结构。对于随机初始化，发现原始论文的均匀分布策略性能不佳，而改进的随机初始化（如在相机包围盒内均匀分布更多高斯、使用更好的分布策略）可以显著提升性能。对于NeRF引导初始化，利用NeRF体渲染权重的概率解释，从短期训练的NeRF模型中提取两个信息：(1) 通过体渲染权重分布采样点云作为3DGS的初始化点；(2) 使用NeRF的深度估计作为深度监督信号。NeRF训练仅需几分钟，远快于COLMAP的SfM处理。两种策略的组合使3DGS摆脱了对昂贵SfM的依赖。其设计目的不是单独增加一个模块，而是用上述表示、监督或渲染约束消除原流程在目标任务中的具体失配。系统性地研究了3DGS的各种初始化策略，发现精心设计的随机初始化可以显著优于朴素随机初始化。提出利用低成本NeRF模型的结构蒸馏（点云初始化和深度监督）来替代SfM初始化，证明改进的随机初始化与NeRF结构引导的组合可以达到甚至超越SfM初始化的质量。适用时还应保留论文给出的边界：NeRF引导引入了额外的训练开销（约5分钟）；在某些极端复杂场景下仍不如SfM初始化。
% METHOD: 2404.15259
\paragraph{FlowMap（可微视频 SfM）（2024.04）}
FlowMap~\cite{2404_15259} 本身不是 3DGS 表示方法，而是给 3DGS 提供可微的相机、内参和密集深度初始化。传统 COLMAP 先做特征匹配、外点剔除和 bundle adjustment，输出稀疏点云；FlowMap 则把视频 SfM 改写成一组可用梯度下降求解的光流一致性问题，使它能嵌入后续学习系统，也能直接把密集深度反投影成 3DGS 初始点云。

输入视频先用现成 optical flow 和 point tracking 得到短期、长期对应。FlowMap 不把深度、位姿和内参都当自由变量直接优化，而是设计了三种重参数化：深度由一个 CNN/深度网络预测并按场景优化；给定深度和对应关系后，相邻帧位姿通过线性最小二乘模块求出；内参则通过光流一致性做可微选择。训练目标比较两种光流：一种是由当前深度、位姿、内参把像素投到另一帧得到的 camera-induced flow，另一种是预计算对应给出的 flow/track。长期 point tracks 约束同一 3D 点在多帧中的一致性，避免只靠相邻光流造成漂移。输出的逐帧 dense depth 可以和相机参数一起反投影成稠密点云，再交给 3DGS 训练；论文显示这个下游 3DGS 结果接近 COLMAP SfM，且明显优于早期基于可微渲染的 pose optimization 方法。相对 InstantSfM 的 GPU 全局 SfM，FlowMap 的锚点价值在于“用可微视频几何替代传统 SfM 预处理”。\textbf{方法信息：} 代码：\url{https://github.com/dcharatan/flowmap}；GitHub 关注度很高，约 980 stars；引用很多，Semantic Scholar 约 66 次。
\textbf{设计动机与边界：} 相对原始 3DGS 以独立高斯和图像重建误差为主的设定，FlowMap 把问题落在以下可实现的改动上：FlowMap的核心思想是将SfM问题转化为端到端可微的梯度下降优化。输入为视频序列，预计算光流和点追踪作为对应关系。目标函数为简单的最小二乘：比较由深度、内参和位姿诱导的光流与预计算对应关系之间的差异。三个关键的可微分重参数化设计：(1) 深度通过神经网络参数化，而非作为自由变量直接优化；(2) 位姿通过深度和光流的线性最小二乘解来参数化，无需自由参数；(3) 内参通过基于光流一致性的可微分选择来确定。此外，利用点追踪（而非仅光流）鼓励长期几何一致性。整个方法不依赖RANSAC等不可微操作，完全可微，可嵌入端到端深度学习流程。输出为密集的逐帧深度估计（而非SfM的稀疏点云）。其设计目的不是单独增加一个模块，而是用上述表示、监督或渲染约束消除原流程在目标任务中的具体失配。提出FlowMap，端到端可微方法通过梯度下降优化精确估计视频序列的相机位姿、内参和逐帧密集深度。方法通过最小化光流诱导与预计算对应关系的简单最小二乘目标实现。引入深度、内参和位姿的可微分重参数化。恢复的相机参数和深度使Gaussian Splatting在360度轨迹上实现照片级新视角合成，性能与COLMAP持平。适用时还应保留论文给出的边界：对极端长视频序列的计算效率有待优化；假设静态场景，不处理动态物体；对初始深度网络参数有一定敏感性。
% METHOD: 2407.02918
\paragraph{Free-SurGS（手术场景 SfM-free 重建）（2024.07）}
Free-SurGS~\cite{2407_02918} 解决的是手术内窥镜视频中的 SfM 失效问题。手术场景纹理少、镜面反射多、光照变化明显，COLMAP 依赖的局部特征匹配和光度一致性很容易崩掉；如果直接把错误位姿和稀疏点云交给 3DGS，高斯会长成漂浮物或错误几何。Free-SurGS 的做法是完全绕开 SfM，把位姿估计改写成“投影流对齐光流”的问题。

流程从单目视频开始。第一帧用 Depth Anything 估计深度，并反投影成点云来初始化高斯中心，颜色和其他高斯参数沿用 3DGS 的初始化规则。随后每来一帧，方法先用 RAFT 估计相邻帧光流，把当前 3D 高斯在候选相机位姿下投影到下一帧，由高斯中心的二维位移得到 projection flow。优化相机位姿时，不再主要依赖渲染图和输入图之间的 photometric loss，而是让 projection flow 尽量匹配 RAFT 光流；极线几何一致性检查会剔除不满足刚性运动的光流点，避免软组织形变、反光和遮挡把位姿带偏。

完成当前帧位姿后，Free-SurGS 再随机采样已有帧来优化高斯的位置、尺度、旋转、不透明度和颜色，并用 progressive growing 逐步扩展高斯覆盖区域。这样，位姿估计依赖当前高斯提供的三维投影，高斯优化又依赖已估计位姿提供多视角监督，二者交替推进。相对 InstantSplat 的通用几何基础模型路线，Free-SurGS 是面向内窥镜/手术视频的专门 SfM-free 设计。\textbf{方法信息：} Tencent AI Lab 参与了这篇论文；代码：\url{https://github.com/wrld/free-surgs}。
\textbf{设计动机与边界：} 相对原始 3DGS 以独立高斯和图像重建误差为主的设定，Free-SurGS 把问题落在以下可实现的改动上：Free-SurGS的核心是将位姿估计问题转化为最小化投影流和光流之间的差异。具体流程：(1) 初始化：利用单目深度（Depth Anything）将首帧像素反投影为3D点云，初始化高斯位置；(2) 位姿估计：利用RAFT获取帧间光流作为伪真值，将3D高斯投影到下一帧得到投影流，通过最小化投影流和光流之间的损失优化相机位姿。引入极线几何一致性检查过滤不可靠的光流点，保留刚性和可靠的对应关系；(3) 场景优化：随机采样帧优化3D高斯参数，采用渐进式增长策略逐步扩展3D高斯覆盖范围。整个过程不依赖COLMAP/SfM。其设计目的不是单独增加一个模块，而是用上述表示、监督或渲染约束消除原流程在目标任务中的具体失配。提出首个SfM-free的3DGS手术场景重建方法Free-SurGS，利用光流先验替代光度损失引导投影流优化相机位姿，并通过一致性检查过滤异常值，在SCARED数据集上实现优于现有方法的位姿估计和新视角合成。适用时还应保留论文给出的边界：依赖光流估计的准确性；对快速运动和大变形场景效果可能受限；深度先验的单目深度可能不准确
% METHOD: 2410.15392
\paragraph{EF-3DGS（事件辅助自由轨迹 3DGS）（2024.10）}
EF-3DGS~\cite{2410_15392} 面向高速或低帧率自由轨迹视频。CF-3DGS、LocalRF 这类帧相机方法依赖相邻帧之间有足够重叠；一旦相机移动太快，帧间位移变大，光流、特征匹配和 photometric pose refinement 都会不稳定。EF-3DGS 引入事件相机，因为事件流在两帧之间连续记录亮度变化，可以补上普通帧没有采到的运动和结构线索。

它把事件信息接到 3DGS 优化的三个位置。第一是 EGM driven optimization：相邻视频帧之间的时间区间被划分为多个子区间，事件在每个子区间内累积成事件帧，再和最近的普通帧积分，恢复中间时刻的潜在强度图像。3DGS 在这些中间时刻渲染图像，并用 L1 与 D-SSIM 形式的事件重渲染损失监督，从而把离散帧之间的盲区变成可训练视角。第二是 CMax 与 LEGM：方法用 3DGS 当前渲染深度和相对位姿推导事件运动场，将事件帧扭曲到参考时刻；如果位姿和深度正确，warped events 会聚成更锐利的 image of warped events，因此可以最大化其对比度来校正相机轨迹。线性事件生成模型进一步把亮度变化写成图像梯度和光流的关系，让事件也在梯度域约束高斯表示。第三是 photometric bundle adjustment：由于事件本身没有颜色，方法仍需普通 RGB 帧来约束颜色和跨视角外观一致性，并通过固定结构/优化颜色的训练策略减小结构和颜色互相拉扯。

相对 E2GS 更偏向用事件解决模糊输入重建，EF-3DGS 的锚点意义在于把事件相机放进自由轨迹位姿估计和 3DGS 场景优化闭环。代码材料需要区分：NVLabs 的 CF-3DGS 仓库是相关基线，不是 EF-3DGS 官方实现；EF-3DGS 项目页目前列出论文与结果，未确认官方代码链接。\textbf{方法信息：} 项目页：\url{https://lbh666.github.io/ef-3dgs/}。

\textbf{实现细节：} 最后，边缘加权损失在重建过程中根据边缘置信度空间加权光度误差：$w(x) = 1 + \beta \cdot M(x)$，其中 $M(x)$ 为边缘掩码、$\beta$ 为权重系数，使优化优先聚焦于几何显著性边界，实现更快的收敛和更准确的深度估计。EF-3DGS（2024.10） 首次将事件相机引入随意拍摄视频的场景构建，通过事件生成模型（将事件流与亮度变化关联）、对比度最大化框架（从事件流中提取边缘信息）和光度束调整（联合优化相机位姿和场景结构）三个组件融合事件和帧信息。
\textbf{实现补充：} 实现上还要核对输入、输出、参数、特征、初始化、训练损失、正则、渲染或推理接口，以及优化中的剪枝、分裂、采样、编码或解码步骤；这些抓手决定该方法到底改了 3DGS 管线的哪一环。
\textbf{设计动机与边界：} 相对原始 3DGS 以独立高斯和图像重建误差为主的设定，EF-3DGS 把问题落在以下可实现的改动上：方法包含三个核心组件：(1)EGM驱动优化：将相邻视频帧间的时间区间划分为N个子区间，利用事件流累积形成高时间分辨率的事件帧，重建中间时刻的潜在强度图像作为渲染监督信号，解决高速场景下视角稀疏问题；(2)统一CMax和LEGM优化：通过对比度最大化框架和线性事件生成模型，利用3DGS渲染深度和相对相机运动构建运动场，通过最大化扭曲事件图像(IWE)的锐度来校正估计姿态，同时利用IWE中的亮度信息在梯度域约束3DGS；(3)光度束调整(PBA)：引入光度约束确保事件和帧之间的跨视角一致性，采用分段方式进行CMax计算。其设计目的不是单独增加一个模块，而是用上述表示、监督或渲染约束消除原流程在目标任务中的具体失配。首次将事件相机与3DGS结合用于自由轨迹场景重建，通过EGM融合事件和帧、CMax校正姿态、PBA确保跨视角一致性适用时还应保留论文给出的边界：需要事件相机硬件；事件数据缺少颜色信息；对事件相机标定精度敏感
% METHOD: 2412.01553
\paragraph{SFGS（分层训练的无 SfM 3DGS）（2024.12）}
SfM-Free 3D Gaussian Splatting via Hierarchical Training~\cite{2412_01553} 面向只有视频、没有 SfM 位姿和点云的输入。它继承 CF-3DGS 的无 SfM 思路，但把问题拆得更细：先估计相邻帧相对位姿，再用分层合并解决单帧初始化覆盖不足。原始 3DGS 的密度控制依赖已有高斯附近的梯度来 clone/split；如果第一帧反投影只能覆盖局部场景，后续未覆盖区域梯度很弱，高斯很难“长”过去。

流程先把第一帧设为参考坐标，对每对相邻帧训练一个单图 3DGS，并固定这个高斯模型，只优化作用在高斯中心上的仿射变换；渲染出的下一帧若能和真实下一帧对齐，这个仿射变换就被解释成相邻帧相对位姿。相邻帧运动过大时，方法用视频插值模型生成中间帧，把一次大位姿估计拆成两次小位姿估计，从而减少单图 3DGS 看不到区域带来的伪影。得到整段视频位姿后，SFGS 不尝试一次性训练完整场景，而是把视频切成重叠片段，每个片段训练一个基础 3DGS；随后按层级把相邻片段的高斯集合合并，保留重要高斯，再用原始帧、基础模型渲染出的伪视图和 VFI 插值帧共同再训练。这样，片段模型负责填补彼此不可见区域，合并本身相当于一次跨片段 densification。相对 LongSplat 的长序列增量优化，SFGS 的关键差异是把无 SfM 初始化写成“相邻位姿堆叠 + 局部高斯模型合并”的分层训练问题。\textbf{方法信息：} 代码：\url{https://github.com/jibo27/3DGS_Hierarchical_Training}；GitHub 约 75 stars。
\textbf{设计动机与边界：} 相对原始 3DGS 以独立高斯和图像重建误差为主的设定，SfM-Free 3D Gaussian Splatting via Hierarchical Training 把问题落在以下可实现的改动上：SFGS方法包含三个核心模块。第一，相机位姿估计：通过堆叠相邻帧间的相对位姿获取全局相机位姿，利用单图3DGS模型和仿射变换估计帧间相对位姿。第二，分层训练策略：将视频序列分割为重叠片段，每个片段训练一个基础3DGS模型，然后从相邻片段迭代合并，每次合并后用原始训练帧、基础模型生成的伪视图帧和VFI插值帧进行再训练，直到获得整个场景的统一3DGS模型。第三，视频帧插值（VFI）增强：使用现成的深度模型将输入视频长度翻倍，桥接较大相机运动间隔，同时作为额外监督信号。该方法有效解决了无SfM场景下高斯初始化困难和自适应密度控制在稀疏区域失效的问题。其设计目的不是单独增加一个模块，而是用上述表示、监督或渲染约束消除原流程在目标任务中的具体失配。提出首个面向视频输入的无SfM预处理3DGS方法SFGS。引入分层训练策略，将视频分段训练多个子3DGS模型后逐步合并为统一模型，解决了无初始点云情况下高斯基元初始化和增长的难题。利用视频帧插值模型弥补大相机运动，并提供多源监督减少过拟合。在Tanks and Temples上PSNR平均提升2.25dB，CO3D-V2上提升1.74dB，大幅超越已有无SfM方法。适用时还应保留论文给出的边界：依赖视频输入且假设小相机运动；相对位姿堆叠可能累积误差；处理大规模场景时计算成本较高；VFI生成的插值帧缺乏完美几何一致性。
% METHOD: 2505.13440
\paragraph{Pensieve（未标定视频的隐式到显式两阶段）（2025.05）}
Recollection from Pensieve~\cite{2505_13440} 从未标定视频学习新视角合成和位姿，输入只有原始视频帧，不提供相机内外参或额外几何先验。它和 SPFSplat 同属自监督 pose-free 路线，但问题判断不同：直接让网络一开始就预测显式高斯，渲染梯度只能在高斯投影范围内建立对应，容易被错误初始化锁住；Pensieve 先让模型在隐空间学会跨帧重建，再把隐式表示对齐到物理 3D 高斯。

第一阶段由共享 Per-Frame Transformer 提取每帧特征，Camera Transformer 预测每帧焦距、旋转和平移，Context Transformer 从严格少于全部帧的上下文中编码场景；随后一个 LVSM 风格的视图合成模块根据预测相机和上下文特征重渲染输入帧。这个阶段不显式输出高斯，等价于让网络先学“哪些帧应该对应同一场景、相机参数如何组织视图变化”。第二阶段在同一框架中额外预测像素对齐的 2DGS/3DGS primitive，并加入 Gaussian Splatting 渲染损失和 depth reprojection loss：前者要求显式高斯能从给定相机渲染目标图，后者用跨视图深度投影把隐式相机和显式几何对齐。论文还区分 target-aware evaluation 和 target-aligned evaluation，说明训练时 target 图像可用于估计 target pose，但真正评估 NVS 时仍要避免把目标图像信息当作渲染输入。相对 SPFSplat 的一步显式预测，Pensieve 的增量是用“隐式预训练缓冲层”降低自监督对应学习难度，再用显式高斯和深度投影把潜空间拉回真实几何。\textbf{方法信息：} 代码：\url{https://github.com/dwawayu/pensieve}；GitHub 约 50 stars。
\textbf{设计动机与边界：} 相对原始 3DGS 以独立高斯和图像重建误差为主的设定，Recollection from Pensieve 把问题落在以下可实现的改动上：方法包含两个训练阶段：（1）第一阶段（隐式预训练）：给定输入视频，网络预测每帧的隐式相机参数和场景上下文特征。采用视图合成模型（受LVSM启发）作为代理渲染器，通过重渲染输入帧实现自监督。网络架构包括共享的Per-Frame Transformer提取特征，Camera Transformer预测相机参数，Context Transformer重建场景。此阶段完全端到端，避免了显式3D表示带来的优化偏差。（2）第二阶段（显式对齐）：额外预测像素对齐的3D高斯基元（采用2DGS的几何公式），计算3DGS渲染损失和深度重投影损失。此阶段将隐式重建与真实3D几何对齐，使模型能够从给定真实相机参数合成新视角。使用Gumbel-Softmax重参数化等技巧确保梯度可传播。其设计目的不是单独增加一个模块，而是用上述表示、监督或渲染约束消除原流程在目标任务中的具体失配。提出从原始视频帧无需相机参数即可训练视角合成模型的两阶段策略。第一阶段在隐空间进行隐式场景重建，避免显式3D表示的优化偏差；第二阶段引入3DGS渲染损失和深度投影损失，将隐式表示与物理3D几何对齐。两阶段互补协同，实现了高质量新视角合成和精确相机估计。适用时还应保留论文给出的边界：两阶段训练增加了整体训练复杂度；对于高度动态或大基线场景可能存在挑战；相机模型假设理想针孔模型和固定内参。
% METHOD: 2508.01171
\paragraph{SPFSplat（完全无位姿前馈 3DGS）（2025.08）}
No Pose at All~\cite{2508_01171} 提出 SPFSplat，目标是在训练和推理阶段都不使用真实相机位姿。它和 NoPoSplat、Splatt3R 一样从稀疏多视图图像直接预测 canonical space 中的像素对齐高斯，但更激进的一点是：训练时目标视角位姿也由模型自己估计，而不是由数据集标注提供。因此它要同时避免两个反馈问题：位姿错会让渲染监督错，渲染错又会反过来污染位姿学习。

网络由共享 ViT 编码器、跨视图 ViT 解码器、pose head 和 Gaussian heads 组成。每个输入视图被 patchify 成 token，可选地加入内参 token；解码器让不同视图 token 做交叉注意力，随后一个 DPT 风格 head 输出每个像素对应的高斯中心，另一个 head 输出旋转、尺度、不透明度和 SH 颜色，所有高斯都放到第一视图定义的 canonical space。pose head 复用同一套解码特征，通过全局池化和轻量 MLP 预测各视图到参考视图的相对位姿。训练时，context views 生成高斯，target image 只用于估计 target pose 和计算渲染损失；模型用估计 target pose 渲染目标图，再用 L2/LPIPS 类图像损失训练外观。同时，reprojection loss 把预测的像素对齐高斯中心投回其他视图，要求它们和图像对应关系一致，防止高斯只靠颜色拟合而失去几何对齐。相对 InstantSplat 的“先用几何基础模型给初始化，再做 3DGS”路线，SPFSplat 把位姿估计和高斯预测合进同一个前馈模型，并用重投影约束压住自监督训练的不稳定性。\textbf{方法信息：} 代码：\url{https://github.com/ranrhuang/SPFSplat}；GitHub 关注度较高，约 153 stars；引用很多，Semantic Scholar 约 33 次。
\textbf{设计动机与边界：} 相对原始 3DGS 以独立高斯和图像重建误差为主的设定，No Pose at All 把问题落在以下可实现的改动上：SPFSplat架构包含四个组件：(1)ViT编码器：将输入图像patch化并编码为token序列，可选地注入相机内参；(2)ViT解码器：通过交叉注意力聚合多视图信息；(3)位姿头：估计输入视图相对于参考视图的相对变换（旋转和平移）；(4)高斯预测头：预测像素对齐的3DGS参数（中心、旋转、缩放、不透明度、球谐颜色）。训练时使用估计的位姿进行新视角渲染（自监督范式），通过渲染损失和重投影损失联合优化3DGS和位姿估计。重投影损失强制像素对齐高斯基元满足跨视角几何约束。关键创新在于共享骨干促进场景重建和位姿学习的一致特征表示，解决自监督位姿-重建耦合的训练不稳定性。其设计目的不是单独增加一个模块，而是用上述表示、监督或渲染约束消除原流程在目标任务中的具体失配。提出SPFSplat，首个在训练和推理阶段均无需真实位姿的3DGS框架。采用共享特征提取骨干，单次前馈同时预测3DGS和相机位姿。引入重投影损失增强像素对齐高斯基元的几何约束。在显著视角变化和有限图像重叠下仍达到SOTA的新视角合成和位姿估计性能。适用时还应保留论文给出的边界：假设已知相机内参；在极端稀疏（2视图）场景中性能受限；估计位姿的精度仍低于SfM。
% METHOD: 2508.14041
\paragraph{LongSplat（随意长视频的无位姿 3DGS）（2025.08）}
LongSplat~\cite{2508_14041} 是无位姿长视频 3DGS 的重要锚点。InstantSplat、MASt3R+3DGS 这类方法已经说明几何基础模型可以替代部分 SfM，但它们通常处理少量图像或较短序列；随手拍长视频的问题更复杂：COLMAP 可能失败，MASt3R 的尺度和相机误差会随序列累积，直接把每帧点云都转成高斯又会很快爆显存。LongSplat 的改变，是把“先估全局位姿再训练 3DGS”改成流式的 pose tracking、octree anchor 表示和局部/全局联合优化。

它的启动阶段用 MASt3R 对初始若干帧做全局对齐，得到深度、对应关系和粗相机；这些点云不直接全部变成自由高斯，而是通过 Octree Anchor Formation 变成空间锚点。具体做法是把点云体素化，再按局部密度递归 split 或 prune：密度高、结构复杂的区域使用更细 anchor，平坦或稀疏区域保留较粗 anchor。每个 anchor 类似 Scaffold-GS 的局部支架，保存特征并解码多个 neural Gaussians。这样 LongSplat 继承了 Scaffold-GS 的 anchor 思想，但用途不同：Scaffold-GS 主要为紧凑表示和视角自适应渲染，LongSplat 则用 anchor 控制长序列增量建图的内存增长。

每插入新帧时，系统先利用已有 3DGS/anchor 状态和新图像之间的对应关系做 PnP+RANSAC，得到粗位姿；然后用可微渲染做 photometric pose refinement，并用深度、关键点重投影和 RGB 光度损失共同约束。如果 PnP 失败或误差过大，系统触发全局重优化以恢复轨迹。位姿确定后，新帧中由遮挡 mask 识别出的新观察区域被反投影成点，再合并进 octree anchor。优化不是每来一帧就全场景更新，而是先选一个 visibility-adapted local window，只更新与当前帧共享可见 anchor 的局部相机和高斯；周期性 global refinement 再传播长期一致性。相对 LocalRF 的分段局部辐射场，LongSplat 的局部窗口由 anchor 可见性决定，并且相机和高斯始终在同一目标中被联合修正。
\textbf{实现细节：} LongSplat（2025.08） 针对随意拍摄的长视频，引入增量联合优化（同时优化位姿和高斯）、鲁棒位姿估计模块和八叉树（octree，一种将三维空间递归划分为八个卦限的层次化数据结构）锚点形成机制。KeyGS（2024.12） 以关键帧为中心，先用 SfM 快速获取粗位姿（秒级），再利用 3DGS 的稠密表示精修位姿，将训练时间从小时级压缩到分钟级。
\textbf{方法信息：} NVIDIA 参与了这篇论文；代码：\url{https://github.com/nvlabs/longsplat}；GitHub 关注度高，约 794 stars；引用很多，Semantic Scholar 约 27 次。



% METHOD: 2509.17246
\paragraph{SPFSplatV2（位姿 token 与掩码注意力）（2025.09）}
SPFSplatV2~\cite{2509_17246} 延续 SPFSplat 的目标：训练和推理时都不给真实相机位姿，只从稀疏多视图图像直接预测 3D 高斯和相机位姿。SPFSplatV2 的问题意识更具体：前一代用全局平均池化从图像 token 中读出位姿，容易把局部对应、边缘和遮挡等关键几何线索平均掉；同时，自监督渲染损失把“高斯是否正确”和“位姿是否正确”耦合在一起，训练容易形成反馈不稳定。

架构上，输入的 context views 先经过共享 ViT 编码器，特征在 MASt3R 风格或 VGGT 风格的解码器中做跨视图交互。每个视图都附加 learnable pose token，SPFSplatV2-L 还区分参考视图 token 和非参考视图 token；这些 token 在解码时通过注意力选择与位姿最相关的图像特征，再由 pose head 输出相对参考视图的旋转和平移。高斯 head 则为 context views 预测像素对齐的高斯中心、旋转、尺度、不透明度和 SH 颜色，并统一放在第一视图定义的 canonical space 中。

训练时还会输入 target views，但掩码注意力禁止 context tokens 从 target 图像泄漏信息，保证高斯重建只由 context views 产生；target tokens 可以读取 context 特征来估计目标位姿。随后用估计目标位姿渲染目标图像，通过渲染损失训练高斯和位姿，并加入 reprojection loss，使像素对齐高斯在估计位姿下与多视图几何一致。相对 SPFSplat，它真正的增量不是“再做一个前馈网络”，而是用 pose token、masked attention 和重投影约束把自监督 pose-free 训练变稳。公开 README 给出 MASt3R/VGGT 两种变体、Hugging Face 预训练权重和 `src.main +experiment=...` 的训练/测试入口。\textbf{方法信息：} 代码：\url{https://github.com/ranrhuang/SPFSplatV2}。
\textbf{设计动机与边界：} 相对原始 3DGS 以独立高斯和图像重建误差为主的设定，SPFSplatV2 把问题落在以下可实现的改动上：SPFSplatV2采用共享特征提取骨干，同时从无位姿输入预测3D高斯和相机位姿。架构基于MASt3R/VGGT风格的ViT编码器-解码器。关键设计：(1) 可学习位姿token：替代SPFSplat的全局平均池化，避免稀释关键几何线索，通过交叉注意力在解码过程中与图像token交互提取位姿信息；(2) 掩码注意力机制：训练时高效估计目标视角位姿，避免不稳定耦合；(3) 重投影损失：强制像素对齐的高斯原语与估计位姿一致，提供更强几何约束。两个模型变体：SPFSplatV2（MASt3R架构）和SPFSplatV2-L（VGGT设计）。训练仅使用渲染损失（无需位姿真值监督），具有更强数据集扩展性。其设计目的不是单独增加一个模块，而是用上述表示、监督或渲染约束消除原流程在目标任务中的具体失配。提出SPFSplatV2，高效的feed-forward框架，从无位姿稀疏多视角图像同时预测3D高斯和相机位姿。引入可学习位姿token替代全局平均池化，掩码注意力机制高效估计训练目标位姿，重投影损失提供像素对齐几何约束。无需位姿监督即达到SOTA性能，且泛化到域外场景。适用时还应保留论文给出的边界：自监督训练可能面临位姿-重建耦合不稳定的风险；对极端大基线视角的泛化仍有提升空间。
% METHOD: 2510.13310
\paragraph{InstantSfM（GPU 原生 SfM）（2025.10）}
InstantSfM~\cite{2510_13310} 不是新的高斯表示，而是把 3DGS 上游最耗时的 SfM 前端改成 GPU 原生模块。它接收多视图图像，先做特征提取和匹配，再在 PyTorch 里完成旋转平均、全局定位和 BA，输出相机位姿与稀疏点云，供后续 3DGS 初始化使用。相对 COLMAP 这类 CPU/C++ 管线，它的重点不是改变渲染方程，而是让位姿和点云估计能进入现代 GPU 学习管线。

它的关键改动有两处。第一，depth-constrained Jacobian structure 把 RGB-D 或单目深度给出的度量深度直接接入全局定位和 BA，而不是等重建完成后再做尺度对齐。对于有有效深度的观测，方法将对应的逆深度尺度变量固定为已知值，并从优化变量中移除；这些“尺度锚点”通过共享稀疏雅可比结构传播到相机和没有深度的点。BA 阶段再加入深度残差，使重投影误差和深度一致性共同约束相机、点和内参。第二，dynamic parameter extraction 在每次 LM 迭代中检查离群观测和欠约束变量，把会导致秩亏的变量动态移出当前正规方程，从而避免大规模全局 SfM 因少量坏匹配而整体发散。

对 3DGS 读者来说，InstantSfM 可以理解为“更快、更可组合的 COLMAP 替代前端”：它仍然服务于相机和点云初始化，但深度先验、稀疏雅可比和动态变量过滤让它更适合 RGB-D、机器人和大规模场景。公开仓库 README 还给出 `ins-feat` 和 `ins-sfm` 的命令行流程，以及 Docker demo，说明它已不是只有论文描述的概念系统。\textbf{方法信息：} 代码：\url{https://github.com/cre185/instantsfm}；GitHub 关注度较高，约 486 stars。
\textbf{设计动机与边界：} 相对原始 3DGS 以独立高斯和图像重建误差为主的设定，InstantSfM 把问题落在以下可实现的改动上：InstantSfM的核心包含两个技术贡献。第一，提出深度约束雅可比结构（Depth-Constrained Jacobian Structure），将度量深度先验作为实时约束直接嵌入全局定位（GP）和光束法平差（BA）优化中。具体而言，当度量深度测量可用时，将每观测尺度变量固定为已知值（1/深度），从优化参数集中移除，使深度先验通过共享稀疏雅可比结构传播至未约束的相机和点。第二，引入动态参数提取（Dynamic Parameter Extraction）机制，在每次LM迭代中实时识别并移除离群值，维持全秩参数空间，这是现有CPU和GPU求解器所不具备的能力。整个系统完全基于PyTorch编程原语实现，使用PyTorch原生稀疏数据结构和GPU并行化算子，无需外部C++优化工具链。系统流程包括：特征提取与匹配、旋转平均、平移平均（全局定位）、光束法平差，所有阶段均在GPU上执行。其设计目的不是单独增加一个模块，而是用上述表示、监督或渲染约束消除原流程在目标任务中的具体失配。提出首个完全基于GPU和PyTorch的全局SfM系统InstantSfM，通过深度约束雅可比结构将度量深度先验嵌入全局定位和BA优化中解决尺度歧义，并采用动态参数提取机制确保数值鲁棒性。在大规模场景上相比COLMAP实现约40倍加速，同时保持可比较的重建精度。适用时还应保留论文给出的边界：系统对初始特征匹配质量仍有依赖；在极端稀疏视图场景下的鲁棒性有待验证；全局SfM方法在处理大规模闭环场景时仍存在挑战。
\textbf{尾项实现澄清：} InstantSfM 要解决的是 CPU 导向 SfM 与 GPU/学习式 3DGS 管线之间的上游瓶颈。其输入是多视图特征匹配和可选度量深度，依次在 GPU 上执行旋转平均、平移平均和 BA；深度约束雅可比将已知尺度直接固定并经稀疏优化传播，动态参数提取在每轮 LM 优化移除离群或欠约束变量。它输出相机位姿与点云，作为 3DGS 初始化而非替代 3DGS 渲染；优化目标仍是几何重投影一致性和深度约束。

\section{去模糊与图像退化处理}

真实采集不只缺视角，还常有运动模糊、散焦、滚动快门、低光照、压缩视频和颜色不一致。原始 3DGS 假设训练图像是清晰、曝光一致的针孔相机观测；退化处理方法的共同改变是把成像过程也纳入优化，而不是让高斯外观硬拟合退化图像。
% METHOD: 2404.04211
\paragraph{Robust Gaussian Splatting（位姿、模糊与颜色鲁棒）（2024.04）}
Robust Gaussian Splatting~\cite{2404_04211} 提供一个统一鲁棒框架，目标不是只处理某一种模糊，而是同时放松原始 3DGS 的三个输入假设：位姿完全准确、图像清晰、颜色一致。它的核心是把运动模糊建模为相机位姿上的 Gaussian 分布。也就是说，一张模糊训练图像不再对应一个固定相机，而对应曝光期间一簇局部扰动相机的期望观测；优化时每张图像只需少量额外位姿噪声参数，就能同时解释相机 pose error 和 motion blur。

由于 3DGS 本身就是 Gaussian primitive，论文可以用一阶近似把相机扰动传播到每个 3D 高斯的投影均值和协方差上，并给出 opacity scaling，保持高斯质量不因协方差放大而改变。这样训练时不必显式渲染大量曝光子帧，也能得到“模糊后的高斯投影”。散焦模糊则通过投影平面上的额外协方差补偿建模：每张图像学习与焦平面偏差相关的 defocus 参数，使离焦区域的 projected Gaussian 更宽。颜色不一致性由每图像 affine color transform 处理，用来吸收白平衡、阴影、环境光或自动曝光差异。

相对 Deblurring 3DGS 的训练时 MLP 协方差偏移，Robust Gaussian Splatting 更强调概率公式内的解析传播和少参数鲁棒校正；相对只做 pose refinement 的方法，它同时处理模糊和颜色漂移。因此它适合作为真实手机/手持采集场景中的鲁棒 3DGS 锚点。\textbf{方法信息：} Meta 参与了这篇论文。
\textbf{设计动机与边界：} 相对原始 3DGS 以独立高斯和图像重建误差为主的设定，Robust Gaussian Splatting 把问题落在以下可实现的改动上：方法从三个维度提升3DGS鲁棒性。核心贡献是将运动模糊建模为相机位姿上的Gaussian分布：模糊图像可理解为曝光时间内多个局部扰动相机捕获的积分。仅需每张图像6个额外参数即可建模运动模糊，通过解析积分推导更新后的每Gaussian参数。这同时解决了位姿精修和运动模糊问题。散焦模糊补偿通过修改Gaussian的协方差来模拟散焦效应，使用可学习的每图像散焦参数。颜色不一致性处理通过可学习的每图像颜色变换矩阵解决由环境光、阴影或相机白平衡设置变化引起的颜色差异。所有机制与3DGS概率公式无缝集成，保持高效训练和渲染。其设计目的不是单独增加一个模块，而是用上述表示、监督或渲染约束消除原流程在目标任务中的具体失配。提出将运动模糊建模为相机位姿上的Gaussian分布，统一解决位姿精修和运动模糊校正。同时提出散焦模糊补偿机制和颜色不一致性处理方法。所有方案与3DGS无缝集成，保持训练效率和渲染速度。在Scannet++和Deblur-NeRF基准上达到SOTA。适用时还应保留论文给出的边界：运动模糊模型假设线性运动轨迹，不适用于非线性运动模糊；每图像额外参数增加存储；散焦模糊模型简化了真实光学模糊；颜色变换假设全局变换无法处理局部光照变化。
% METHOD: 2404.11358
\paragraph{DeblurGS（相机运动模糊轨迹优化）（2024.04）\cite{2404_11358}}
DeblurGS 关注的是另一种比普通“训练图像有点糊”更难的情形：运动模糊不仅让图像监督变模糊，还会让 SfM 位姿本身变差。原始 3DGS 默认每张图对应一个清晰、固定的相机位姿；如果输入是在曝光时间内移动相机得到的模糊图，一张训练图实际对应一段连续相机轨迹。直接用单个位姿训练 3DGS，会把曝光积分造成的拖影写进球谐颜色和高斯形状；更糟的是，错误 SfM 位姿还会在早期 densification 中诱导高斯长到错误位置。

DeblurGS 的改变是联合优化清晰 3DGS 和每张模糊图像的潜在 6DoF 相机运动。对一张模糊图，方法在曝光时间内采样多个 sub-frame pose，从这些虚拟姿态渲染若干清晰图像，再按曝光积分平均成模拟模糊图，与真实模糊观测做重建损失。相机运动用 SE(3) 轨迹参数化，并加入时间平滑约束，使曝光期间的运动不至于变成任意抖动。核心训练策略是 Gaussian Densification Annealing：在相机轨迹还不准的早期，先抑制或延迟 clone/split，避免错误位姿驱动高斯在错误深度处增殖；等轨迹优化更稳定后，再逐步恢复密度控制，让模型补充真实结构。相对 Deblurring 3DGS 的协方差/位置偏移模拟模糊，DeblurGS 更强调“模糊首先破坏了位姿与 densification 时机”；相对 GeMS，它仍依赖优化相机运动和清晰场景，而不是先用学习式 SfM/VGGSfM 替代 COLMAP。当前只确认官方项目仓库 \url{https://github.com/taekkii/deblurgs}，尚未确认完整可复现实装代码。

\textbf{实现细节：} DeblurGS（2024.04） 进一步处理了由严重模糊导致的 SfM 位姿不准问题，估计每张模糊观测的 6-DoF（六自由度，即三维旋转加三维平移）相机运动来合成对应的模糊渲染，并提出高斯致密化退火（annealing，从高温度逐步降温的策略，此处指从抑制致密化逐步放开）策略防止训练早期生成不准确的高斯。
\textbf{方法信息：} 已有可跟进的公开链接，链接为 taekkii/deblurgs (36 stars)，但尚未确认完整代码。入选依据：企业:否；代码:否；星标:否；引用:否。


% METHOD: 2405.20224
\paragraph{EvaGaussians（事件流辅助模糊重建）（2024.05）}
EvaGaussians~\cite{2405_20224} 用事件相机的高时间分辨率恢复模糊输入中的细节。初始化阶段，EDI 模型从模糊图像和事件流估计潜在清晰图像，再用 COLMAP 获取初始位姿和点云。

真正的增量在 bundle adjustment 阶段。对每张模糊图像，EDI 会给出曝光时间内若干潜在清晰时刻及其初始姿态；EvaGaussians 不把这些姿态当固定真值，而是给每个时刻加入可学习位姿 offset，与高斯位置、尺度、旋转、不透明度和颜色一起优化。blur reconstruction loss 将这些时刻的 3DGS 渲染结果平均，要求平均图匹配输入模糊图；event reconstruction loss 则随机选取两个时刻，比较渲染图的对数亮度变化和事件流积分出的亮度变化，从而利用事件的高时间分辨率监督曝光轨迹。

模糊输入还会导致几何深度不稳，因此论文加入两类事件辅助深度正则。单目深度约束把事件恢复的潜在清晰图送入深度估计网络，用作渲染深度的软监督；立体深度约束则利用曝光期间多时刻姿态形成短基线几何，约束跨时刻深度一致。训练上采用 progressive/coarse-to-fine 策略，先低分辨率稳定结构，再逐步升到全分辨率。相对 E2GS 主要把事件用于去模糊损失，EvaGaussians 更强调联合优化曝光期间相机轨迹和 3DGS 几何。\textbf{方法信息：} 引用很多，Semantic Scholar 约 61 次。
\textbf{设计动机与边界：} 相对原始 3DGS 以独立高斯和图像重建误差为主的设定，EvaGaussians 把问题落在以下可实现的改动上：方法包含三个关键组件：事件辅助初始化、联合优化和事件辅助深度正则化。初始化阶段使用Event-based Double Integral(EDI)模型从模糊图像和事件流估计潜在清晰图像，再通过COLMAP获取初始相机位姿和点云。联合优化阶段显式建模运动模糊形成过程，将模糊图像建模为曝光时间内瞬时图像的平均，通过事件流计算任意时刻的潜在图像。引入模糊重建损失和事件重建损失联合优化3DGS参数和相机轨迹。提出两种事件辅助深度正则化策略：事件辅助单目深度约束和事件辅助立体深度约束，解决模糊图像带来的几何模糊问题。最终实现从模糊输入中恢复精细纹理细节。其设计目的不是单独增加一个模块，而是用上述表示、监督或渲染约束消除原流程在目标任务中的具体失配。提出事件流辅助的Gaussian Splatting方法EvaGaussians，利用事件相机的高时间分辨率和高动态范围特性，显式建模运动模糊图像的形成过程，联合优化3DGS参数和曝光时间内相机运动轨迹，从模糊图像中恢复高保真度新视角合成结果。在恢复精细纹理细节方面显著超越现有去模糊渲染方法。适用时还应保留论文给出的边界：需要事件相机硬件，限制了通用性；事件相机和传统相机的同步校准增加系统复杂度；极端运动模糊下事件流可能溢出；不适用于静态模糊（如散焦模糊）。
% METHOD: 2406.14978
\paragraph{E2GS（事件增强 Gaussian Splatting）（2024.06）}
E2GS~\cite{2406_14978} 是最早把事件数据引入 Gaussian Splatting 的方法之一。它用事件流和模糊图像之间的 EDI 关系恢复曝光期间多个清晰时刻，并估计对应相机姿态和初始点云。

它的流程很直接。对每张模糊 RGB 图像，先把曝光时间均分为 $N$ 个时间戳，并把事件流分到相邻事件 bin 中。EDI 模型把模糊图像解释为这些时刻清晰图像的平均，同时用事件积分从一个时刻推到另一个时刻。由此得到的去模糊图像序列被送入 COLMAP，用于估计曝光期间的多个相机姿态和初始点云；真实数据中 README 也说明需要先用 EDI 去模糊，再为每张图估计多个曝光内位姿。

训练 3DGS 时有两类损失。blur rendering loss 在每个曝光内姿态渲染一张图，然后把 $N$ 张渲染图平均，要求平均结果匹配输入模糊图；event rendering loss 随机选择两个渲染时刻，把它们转成对数灰度亮度，并让亮度差与对应事件积分一致。前者保证颜色仍来自 RGB 模糊帧，后者利用事件流告诉模型“曝光期间亮度如何变化”。这样训练出的高斯表示趋向于清晰场景，而不是把模糊纹理写进 SH 颜色。相对 DiET-GS，E2GS 不依赖扩散先验，核心是用 EDI 和事件亮度变化把 3DGS 的训练监督从单张模糊图扩展到曝光时间轴。\textbf{方法信息：} 代码：\url{https://github.com/deguchihiroyuki/e2gs}。
\textbf{设计动机与边界：} 相对原始 3DGS 以独立高斯和图像重建误差为主的设定，E2GS 把问题落在以下可实现的改动上：E2GS是首个将事件相机数据融入3DGS的方法。方法包含以下核心设计：1）预处理：利用事件数据与模糊图像的对应关系进行预处理。使用Event-based Double Integral (EDI)模型，假设模糊图像是曝光时间内多张清晰图像的平均，利用事件数据估计曝光时间内N个时间戳的清晰图像。2）相机姿态和初始点云估计：通过事件辅助的去模糊估计N个等间距相机姿态，从模糊图像中恢复清晰帧来生成初始点云用于Gaussian Splatting初始化。3）训练损失：设计两种考虑模糊的损失函数：a）基于事件数据的去模糊损失，利用事件流在曝光时间内重建清晰图像并与渲染图像对比；b）基于模糊图像的直接损失，将渲染的多帧图像平均与输入模糊图像对比。4）高斯表示：标准3DGS表示，但训练过程中考虑了由事件数据提供的时序信息。整体流程使系统能够从模糊输入中重建清晰场景。其设计目的不是单独增加一个模块，而是用上述表示、监督或渲染约束消除原流程在目标任务中的具体失配。提出首个将事件数据引入Gaussian Splatting的方法E2GS，有效利用模糊图像和事件数据，显著改善图像去模糊并产生高质量新视角合成。训练速度比E2NeRF快60倍，渲染速度达140 FPS（快3500倍）。在合成和真实数据集上验证了有效性。适用时还应保留论文给出的边界：需要同时获取事件相机数据和RGB图像。对事件相机和RGB相机的同步精度要求较高。在极端低光条件下的表现可能有限。
% METHOD: 2401.00834
\paragraph{Deblurring 3DGS（协方差偏移与模糊形成）（2024.01）}
\textbf{锚点细化：} Deblurring 3D Gaussian Splatting：针对标准 3DGS 直接用模糊图像训练时无法建模邻域像素干扰的局限，用``操纵 3D 高斯几何属性正向建模模糊过程''替代图像空间反卷积，核心思想是不在图像空间反卷积去模糊，而是直接放大高斯几何属性正向重建模糊过程——大尺寸高斯覆盖更宽图像区域恰好对应邻域像素对中心像素的干扰，正是模糊的物理根源；引入小型 MLP 以高斯位置/四元数/缩放/视角方向为输入，对散焦模糊输出缩放因子、放大协方差、模拟散焦扩散，对运动模糊输出 $M$ 组时刻偏移量、构造多组高斯集合、渲染后取平均，推理时无需激活 MLP、直接用原始高斯渲染清晰图像、完全保留实时性。
Deblurring 3D Gaussian Splatting~\cite{2401_00834} 是 3DGS 去模糊的早期锚点，解决的是原始 3DGS 只能拟合模糊训练图像、推理时也输出模糊外观的问题。它不把模糊当作后处理卷积，而是把模糊写进高斯本身：训练时临时改变高斯的协方差、位置和旋转来模拟模糊成像，推理时只保留原始高斯参数，因此渲染速度仍接近 3DGS。

散焦模糊分支来自薄透镜直觉：离焦区域在成像面上形成更大的光斑。方法用一个小 MLP 读取每个高斯的属性，预测位置、旋转和尺度的偏移，尤其通过尺度偏移放大训练时投影到图像平面的高斯协方差，让它能匹配局部模糊斑；清晰区域的高斯偏移很小，模糊区域的高斯被选择性放大，因此能处理空间变化散焦，而不是整张图套同一个高斯核。相机运动模糊分支则在曝光期间采样多个相机时刻，分别渲染清晰图像后求平均，用平均图和输入模糊图做重建损失。论文还加入额外点补充和深度感知剪枝，缓解模糊输入下 SfM 点云过稀导致的远处细节缺失。

相对后续 Robust Gaussian Splatting 的统一鲁棒成像模型，Deblurring 3DGS 更像“把 NeRF 去模糊里的模糊核思想移植到 rasterization-based 3DGS”的第一步；它说明只要在训练阶段扩展高斯协方差和曝光采样，就可以从模糊图像中恢复清晰可实时渲染的表示。代码材料需要区分：DiET-GS 仓库对应另一篇后续工作，不是 Deblurring 3DGS；该方法官方实现为 benhenryL 版本。\textbf{方法信息：} 代码：\url{https://github.com/benhenryL/Deblurring-3D-Gaussian-Splatting}；引用很多，Semantic Scholar 约 95 次；GitHub 关注度较高，约 338 stars。



% METHOD: 2503.24210
\paragraph{DiET-GS（事件流与扩散先验去模糊）（2025.03）}
DiET-GS~\cite{2503_24210} 是事件辅助去模糊路线的后续增强版。E2GS 和 EvaGaussians 已经证明事件流能提供曝光期间的高时间分辨率约束，但事件只记录亮度变化，颜色和细小纹理仍可能不稳定。DiET-GS 的增量是把 EDI 事件约束和扩散先验拆成两阶段：第一阶段先让 3DGS 学到可信颜色、结构和基本清晰度，第二阶段再用生成先验补边缘和细节。

第一阶段的输入是模糊 RGB 图像和同步事件流。方法先用 EDI 从模糊图和事件中恢复初始潜在清晰图像，并在曝光时间内采样多个时刻的清晰估计，送入 SfM 得到初始相机和点云。训练时，每张模糊图对应一段近似相机轨迹；3DGS 在轨迹上的多个姿态渲染清晰瞬时图像，平均后得到模拟模糊图，用 L1 与 D-SSIM 组成的 blur reconstruction loss 对齐真实模糊输入。事件重建损失则随机取两个时刻的渲染图，经过亮度转换和可学习相机响应函数后比较对数亮度变化，使它与事件流积分出的亮度变化一致。论文还把 EDI 作为 RGB 空间和 brightness domain 的联合约束：RGB 约束保颜色，亮度域约束保边缘，EDI regularization 稳定没有事件触发的区域。

第二阶段不再重新求几何，而是对第一阶段渲染出的清晰结果使用预训练扩散模型做细节增强，形成 DiET-GS++。这个阶段的作用不是替代事件约束，而是在已有几何和颜色基础上恢复更锐利的边缘。公开 README 也按 Stage 1/Stage 2 给出训练和渲染入口，说明实现上确实是分阶段管线。相对 E2GS，DiET-GS 的主要差异是 EDI 约束更细，并额外引入扩散先验作为最后的视觉细节修复器。\textbf{方法信息：} 代码：\url{https://github.com/diet-gs/diet-gs}。
\textbf{设计动机与边界：} 相对原始 3DGS 以独立高斯和图像重建误差为主的设定，DiET-GS 把问题落在以下可实现的改动上：DiET-GS包含两阶段：(1)第一阶段：利用EDI从模糊图像和事件流恢复初始清晰图像用于SfM初始化；模糊重建损失通过在近似相机轨迹上离散化模糊形成过程监督3DGS；事件重建损失利用事件流的无模糊微秒级测量，通过模拟对数亮度变化约束渲染图像的亮度变化，确保准确颜色和清晰细节；(2)第二阶段：引入预训练扩散模型作为先验，对第一阶段渲染的清晰图像进一步去噪增强，特别提升边缘细节质量。两阶段分离设计使模型先学习基础颜色和结构，再优化细节。其设计目的不是单独增加一个模块，而是用上述表示、监督或渲染约束消除原流程在目标任务中的具体失配。提出DiET-GS框架，首次结合事件双积分约束和扩散先验进行3DGS运动去模糊。事件重建损失约束亮度变化确保准确颜色，扩散先验增强边缘细节，两阶段训练策略分步优化，在合成和真实数据上显著优于现有基线。适用时还应保留论文给出的边界：需要事件相机数据；扩散先验增强增加训练时间；对极端模糊的恢复效果有限
% METHOD: 2506.04444
\paragraph{Egocentric Photoreal Reconstruction（第一人称设备成像模型）（2025.06）}
Photoreal Scene Reconstruction from an Egocentric Device~\cite{2506_04444} 针对第一人称设备。它发现普通帧率相机位姿不足以描述滚动快门 RGB 图像，IMU 中的高频运动对于真实感 3DGS 很关键。

第一部分是 VIBA。方法不再把一张 RGB 图像看成同一时刻的全局快门观测，而是把设备 6DoF 轨迹表示成可按任意时间查询的高频连续函数。对滚动快门相机，第 $h$ 行像素的曝光时间不同，因此每一行、甚至每个校正后像素，都应该从连续轨迹中查询自己的位姿。VIBA 利用设备 IMU 和视觉观测联合优化相机参数、时间偏移和 IMU 偏置，使这些 per-pixel poses 比设备默认 frame-rate VIO 位姿更适合像素级重建。

第二部分是把真实设备成像写进 Gaussian rasterization。每个像素的颜色不只是某个瞬时位姿下的 3DGS 渲染值，而是在曝光时间内对多个运动采样的场景辐照度积分，再乘以模拟增益、镜头阴影和归一化因子，并经过相机响应函数压到图像动态范围。方法还用场景深度估计运动采样间隔，保证曝光或读出时间内的重投影误差足够小。对于低光照，论文建议有意欠曝光采集，用短曝光减少模糊，再通过 HDR/响应建模恢复亮度。

相对 Deblurring 3DGS 的模糊核建模，这项工作更像完整的设备级成像校准：它关心滚动快门、镜头阴影、曝光、增益、HDR 响应和 VIBA 位姿能否一起进入 3DGS/2DGS 家族的栅格化公式。对第一人称眼镜、Project Aria 和 Quest 这类设备，它提供的是“不要把设备视频当普通针孔相机图像”的系统锚点。\textbf{方法信息：} Meta 参与了这篇论文；项目资料：\url{https://www.projectaria.com/photoreal-reconstruction/}。

\textbf{实现细节：} Photoreal Scene Reconstruction from an Egocentric Device（2025.06） 处理第一人称设备（如 Project Aria 眼镜）采集的数据，这类设备配备多传感器（RGB 相机、单目深度传感器、IMU、眼动追踪），且用户在佩戴过程中不断移动，导致采集过程光照持续变化、视角抖动。该方法通过多传感器融合（深度图辅助几何初始化、IMU 提供姿态先验）和针对设备特性的光照补偿策略，实现真实感 3DGS 重建。
\textbf{实现补充：} 实现上还要核对输入、输出、参数、特征、初始化、训练损失、正则、渲染或推理接口，以及优化中的剪枝、分裂、采样、编码或解码步骤；这些抓手决定该方法到底改了 3DGS 管线的哪一环。
% METHOD: 2508.14682
\paragraph{GeMS（极端运动模糊下的 3DGS）（2025.08）}
GeMS~\cite{2508_14682} 针对极端运动模糊输入，不依赖清晰辅助图像。它的出发点是：Deblurring 3DGS、BAD-Gaussians 等方法虽然把模糊形成写进训练，但很多设置仍默认有清晰图像或可靠 COLMAP 初始化；在所有输入都严重模糊时，传统特征匹配和点云生成本身就可能失败。GeMS 因此先替换初始化，再优化模糊形成过程。

第一步用 VGGSfM 替代 COLMAP。VGGSfM 依赖学习式 2D point tracking、全局姿态估计和可微 bundle adjustment，比手工 keypoint matching 更能在模糊纹理上给出可用相机和点云。第二步用 3DGS-MCMC 训练场景：高斯不再只靠固定阈值 clone/split/prune，而被看作场景分布的样本，通过随机梯度 Langevin 动力学移动到高似然区域，减少模糊点云初始化带来的局部空洞。第三步是联合轨迹-几何优化：一张模糊图像对应曝光期间的连续相机运动，GeMS 用 Bézier 曲线表示这段轨迹，在多个虚拟相机时刻渲染清晰图，再按曝光积分形成模拟模糊图，与输入模糊图做物理一致的重建损失。GeMS-E 则在有事件流时先用 EDI 从事件和模糊帧恢复更清晰的图像，再送入 GeMS 管线改善初始化。相对 Deblurring 3DGS 的训练时模糊核/协方差偏移，GeMS 的重点是极端模糊下“先让 SfM 活下来，再让轨迹和高斯一起解释曝光”。\textbf{方法信息：} 当前未确认官方代码仓库；同名 \url{https://github.com/UNITES-Lab/GEM} 属于另一篇 cryo-EM 3DGS 工作，不能作为该方法代码。
\textbf{设计动机与边界：} 相对原始 3DGS 以独立高斯和图像重建误差为主的设定，GeMS 把问题落在以下可实现的改动上：GeMS是一个处理极端运动模糊的3DGS框架。核心组件：(1) VGGSfM初始化：替代COLMAP，使用深度学习端到端SfM框架VGGSfM进行姿态估计和点云生成。其学习的2D点跟踪和端到端优化提供模糊鲁棒的初始化。(2) 3DGS-MCMC建模：将高斯视为场景分布的样本，使用马尔可夫链蒙特卡洛(MCMC)采样自适应致密化和精炼几何，避免启发式剪枝策略。(3) 联合轨迹-几何优化：使用贝塞尔曲线建模曝光期间连续相机运动，联合优化相机轨迹和高斯参数，配合物理损失对齐重建几何与实际模糊形成过程。(4) GeMS-E变体：集成事件数据，使用事件驱动双重积分(EDI)去模糊恢复更清晰图像，进一步改善姿态估计和点云生成。这是首个处理极端运动模糊的3DGS框架。其设计目的不是单独增加一个模块，而是用上述表示、监督或渲染约束消除原流程在目标任务中的具体失配。首个直接从极端模糊图像进行3DGS重建的框架，无需清晰图像辅助。结合深度学习SfM和概率高斯建模，并可选集成事件相机进一步提升质量。适用时还应保留论文给出的边界：VGGSfM对某些极端模糊模式仍有局限；贝塞尔曲线对复杂运动轨迹建模能力有限；需要事件相机硬件支持GeMS-E。
% METHOD: 2602.09816
\paragraph{CompSplat（压缩视频感知训练）（2026.02）}
CompSplat~\cite{2602_09816} 处理长压缩视频。它继承 LongSplat/CF-3DGS 这类长视频或无位姿重建管线，但指出真实视频还有一个常被忽略的变量：H.264、H.265、VVC 等有损压缩会让不同帧的细节保留程度不同，低质量帧会破坏特征匹配、位姿估计和几何更新。如果仍按原始 3DGS 那样让每帧等权参与 densification、pruning 和 photometric loss，模型容易把块效应、纹理丢失和帧间不一致写进高斯几何。

CompSplat 在训练中显式估计帧级可靠性。它从编码器元数据中读取每帧 QP 和 bitrate：较低 QP 通常表示量化更轻，高 bitrate 则说明该帧需要更多信息保存复杂内容；两者线性归一后合成 frame confidence，并用 EMA 平滑到时间序列上。这个置信度接到密度控制中：高置信帧允许更积极的高斯更新和保留，低置信帧提高剪除噪声/过大高斯的倾向，避免被压缩模糊区域诱导出漂浮物。第二个模块是 quality gap-aware masking，当相邻帧质量差距大时，关键点匹配和像素监督都更不可靠，方法用关键点驱动的 mask 调制 photometric loss，使训练更依赖稳定区域。它还加入 scale-based pruning，把因压缩纹理缺失而扩散成大尺度的低可信高斯剔除。相对 Naka-GS 处理低光照成像退化，CompSplat 处理的是视频编码退化；它的改变不在高斯表示本身，而在“哪一帧、哪一区域、哪类高斯应该更可信”这个训练调度层。\textbf{方法信息：} Adobe 参与了这篇论文；仓库：\url{https://github.com/hojunking/CompSplat}，当前 README 写明代码将发布，GitHub 约 1 star。
\textbf{设计动机与边界：} 相对原始 3DGS 以独立高斯和图像重建误差为主的设定，CompSplat 把问题落在以下可实现的改动上：CompSplat包含两个核心组件：1)质量引导密度控制：通过量化参数置信度分数和比特率置信度分数的线性组合计算帧置信度，调控高斯致密化和剪枝——高质量帧获得更多高斯更新，低质量帧减少噪声高斯引入；2)质量差距感知掩码：通过关键点驱动的像素掩码调制光度监督，缓解相邻帧间的质量差距对特征匹配和位姿估计的影响。框架集成到无位姿3DGS管线中，在训练过程中持续更新帧可靠性因子。其设计目的不是单独增加一个模块，而是用上述表示、监督或渲染约束消除原流程在目标任务中的具体失配。提出首个压缩感知3DGS训练框架，通过估计帧级压缩可靠性因子（结合量化参数和比特率）和自适应剪枝策略，显式建模帧间压缩特性以缓解长压缩视频中的几何误差累积和帧间不一致性。适用时还应保留论文给出的边界：帧可靠性估计依赖编解码器元数据；对无损压缩视频无额外收益；长序列优化计算量较大。
% METHOD: 2604.11142
\paragraph{Naka-GS（低光照仿生校正与点预处理）（2026.04）}
Naka-GS~\cite{2604_11142} 处理低光照 3DGS。它不是简单把低光照图像先增强再跑 3DGS，而是把低光照错误拆成两个阶段：图像层面有暗部不可见、亮区偏色和边缘结构误差；几何层面，增强后的图像送入前馈重建模型时会产生噪声点云、局部过密点和漂浮结构。Naka-GS 因此同时修正输入图像和初始化点云。

第一阶段使用 Naka--Rushton 视网膜响应函数做物理先验增强，使暗部响应按非线性曲线抬升但在高亮区域逐渐饱和。直接 Naka 增强仍会留下亮区色偏和边界伪结构，因此方法构造 18 通道双分支输入：原始低光照图、Naka 增强图、二者残差，以及它们各自标准化后的版本。U-Net 风格的 chroma-correction 网络不直接重绘 RGB，而是预测单通道乘性校正和三通道加性校正；同时把 Naka 增强图分解成低频和高频，主要校正低频亮度/色度，把高频残差直接保留，以免纹理被过度平滑。mask-guided supervision 进一步强调亮区和边缘区域，因为这些正是低光照增强后最容易出错的位置。第二阶段把校正图像送入 VGGT 等前馈多视图重建模型，得到相机和密集点云。第三阶段的 PPM 在进入 3DGS 前清理点云：先做坐标对齐，再用 voxel pooling 压缩局部冗余点，最后按距离自适应阈值渐进剪枝，保留代表性结构并减少噪声先验。相对 Robust Gaussian Splatting 的每图像颜色矩阵，Naka-GS 更像低光照专用的“图像校正 + 点云清洗 + 3DGS 优化”管线。\textbf{方法信息：} 代码：\url{https://github.com/runyuzhu/naka-gs}；GitHub 约 3 stars。
\textbf{设计动机与边界：} 相对原始 3DGS 以独立高斯和图像重建误差为主的设定，Naka-GS 把问题落在以下可实现的改动上：三阶段流程：(1)Naka引导的色度校正：基于Naka-Rushton视网膜响应函数进行物理先验预增强，构建18通道双分支输入(原始低光照图、Naka增强图及其残差)，U-Net编码器-解码器预测乘性和加性校正图。将增强图像解耦为低频和高频分量，主要对低频分量校正而直接保留高频残差，避免过度平滑。掩码引导监督强化边缘和亮区的优化；(2)前馈多视角重建：将校正后图像输入多视角重建模型估计相机参数和密集点云先验；(3)点预处理模块(PPM)：坐标对齐、体素池化和距离自适应渐进剪枝去除噪声冗余点，保留代表性结构，改善高斯初始化质量。其设计目的不是单独增加一个模块，而是用上述表示、监督或渲染约束消除原流程在目标任务中的具体失配。提出NAKA-GS，仿生启发的低光照3DGS框架。设计Naka引导的色度校正网络，结合物理先验低光照增强、双分支输入建模、频率解耦校正和掩码引导优化。引入轻量级点预处理模块(PPM)通过坐标对齐、体素池化和距离自适应渐进剪枝去除噪声和冗余点。在NTIRE 3DRR挑战赛中大幅超越基线方法。适用时还应保留论文给出的边界：Naka-Rushton参数需要手动设定，对极端低光照(接近全黑)的效果有限，依赖前馈多视角重建模型的质量。
\textbf{尾项实现澄清：} NAKA-GS 的问题是低照度图像同时污染颜色、边缘和后续高斯初始化。它先以 Naka-Rushton 响应先验构造原图、增强图和残差组成的双分支输入，由 U-Net 预测加性/乘性色度校正，并以频率分解保留高频、以掩码监督抑制亮区伪影；随后将校正视图送入前馈多视图重建得到点云，再以坐标对齐、体素池化和距离自适应渐进剪枝形成干净初始化。损失层面可确认掩码引导校正监督与重建监督，论文包未给出完整权重。

\section{主动视觉与规划}

主动重建方法把 3DGS 从离线重建变成机器人采集决策的一部分。原始 3DGS 可以快速渲染、显式存几何，但它没有告诉机器人下一步去哪看。这里的关键是设计视角效用：基于互信息、参数协方差、对象不确定性、可渲染性、置信度图或未观测体素。
% METHOD: 2405.10142
\paragraph{GS-Planner（四旋翼主动高保真重建）（2024.05）}
GS-Planner~\cite{2405_10142} 是较早把 3DGS 放入机器人规划闭环的系统，可以作为这一节的系统锚点。它的问题定义比 NBV 指标更完整：四旋翼不仅要知道哪个视角对重建有用，还要能安全飞到那里，并且在在线任务时间内持续更新地图。

它以 3DGS/SplaTAM 式在线地图为基础，把未知体素也接入 splatting 过程，使渲染时能暴露哪些区域尚未被观察。视角评分同时考虑几何完整度和纹理质量：几何上，系统估计新视角能覆盖多少未知或低质量区域；纹理上，通过渲染质量和观测历史判断高斯是否还需要更好的视角。规划器采样候选位置和朝向后，选择信息增益高且路径代价可接受的目标。安全性则由高斯位置、密度和障碍约束给出，轨迹优化时加入避障和动力学可执行性，让输出不是抽象相机位姿，而是四旋翼能执行的轨迹。后续 GauSS-MI、ActiveGS 等工作通常只替换其中的地图置信度或视角效用定义。
\textbf{实现细节：} UNG-GS（2025.03） 在 3DGS 管道中引入显式空间不确定性场（SUF）来量化几何不确定性，实现不确定性感知的深度渲染和法线精修。Reprojection-based DoF Separation（2024.12） 提出将高斯位置自由度分离为图像平面平行 DoF（高置信度）和射线方向 DoF（低置信度），分别施加不同约束以抑制几何伪影。DET-GS（2025.08） 提出统一的深度和边缘感知正则化框架，结合层次化几何深度监督、边缘感知深度正则化和 RGB 引导的边缘保留全变差损失。DIP-GS（2025.08） 利用深度图像先验（Deep Image Prior，一种无需预训练的CNN结构本身即可作为自然图像先验的现象）的内部结构模式进行粗到细的 3DGS 优化，无需任何预训练模型。Few-shot Depth Aware GS（2024.10） 使用单目深度预测作为先验配合尺度不变深度损失，并采用低阶球谐函数建模颜色以避免过拟合。
\textbf{方法信息：} 引用很多，Semantic Scholar 约 51 次。
% METHOD: 2412.17769
\paragraph{ActiveGS（混合高斯和体素地图）（2024.12）}
ActiveGS~\cite{2412_17769} 对 GS-Planner 的一个关键问题做了更清晰的系统拆分：3DGS 很适合存高保真表面，却不适合表达未知/自由空间；体素地图适合导航和探索，却不适合最终渲染。因此它并行维护两张地图：Gaussian Surfel 地图负责重建质量，粗体素地图负责占用概率、未知区域和碰撞避障。

在高斯地图中，每个 Gaussian Surfel 除了位置、旋转、二维尺度、颜色和不透明度，还带有置信度 $k_i$。系统可以像渲染 RGB、深度、normal 一样渲染 confidence map，用这张图找到已经被建模但仍欠重建的表面。增量 mapping 时，RGB-D 图像先更新体素占用；随后在当前视角渲染颜色、深度和 opacity，并用低 opacity、颜色误差大或观测深度在当前估计前方等条件生成 densification mask，从这些像素反投影生成新的高斯。规划时，ActiveGS 同时考虑两类目标：体素地图中的未知区域用于探索，低置信高斯周围的 ROI 用于精细检查。公开实现中，\texttt{main.py} 组装 mapper、simulator 和 planner，\texttt{config/planner/confidence.yaml} 中的 confidence planner 使用候选视角数、ROI 采样数、探索权重和路径长度权重控制这一闭环。相对 GS-Planner，ActiveGS 的贡献在于把“未知空间探索”和“已见表面补拍”拆成两种明确的地图信号。
\textbf{锚点细化：} ActiveGS：针对已有主动重建系统要么用 NeRF 渲染慢、难在线更新、要么用体素/TSDF 缺乏高保真外观且依赖耗时射线投射评估信息增益的局限，用``高斯泼溅地图 + 概率体素地图混合表示 + 高斯置信度建模''替代单一地图，核心思想是把``空间理解''（体素地图建模占用概率，用于碰撞检测和路径规划）与``外观重建''（GS 地图基于 Gaussian Surfel 高保真渲染）分离、各取所长，GS 地图每基元附加置信度（距离加权余弦相似度 + 方向离散度），视角效用综合探索效用（可见未探索体素比例）和利用效用（渲染置信度图负均值）、配合 $A^*$ 路径代价。
	extbf{实现细节：} ActiveGS 同时维护两张地图：以 Gaussian Surfel 表示可渲染的高保真表面，并维护体素占据图来标记自由、未知和碰撞空间。每个 surfel 有由观测次数与多视角一致性估计的置信度，渲染置信度图后可定位尚未重建充分的区域。规划器把低置信度区域和体素未知空间合并为候选目标，再在体素图中进行避障路径规划；执行观测后增量更新高斯、置信度和占据概率。这里的核心损失/约束是重建一致性与置信度更新，而非只最小化离线新视图误差。
\textbf{实现细节：} ActiveGS（2024.12） 。Baseline：此前主动重建系统要么使用 NeRF 等隐式表示（渲染速度慢、难以在线更新），要么使用纯体素或 TSDF 地图（缺乏高保真外观），且多数系统依赖耗时的射线投射来评估候选视角的信息增益，无法满足机器人实时规划需求。改了什么：ActiveGS 提出混合地图表示——将高斯泼溅地图（Gaussian Surfel，高保真重建）与粗粒度概率体素地图（空间占用建模与路径规划）融合——并引入高斯基元的置信度（confidence）建模，使系统能在同一框架中同时完成探索（探索未知区域）和利用（精修重建不足表面）两类任务。实现原理：系统以迭代式"建图--规划"循环运行，直到达到预定任务时间。% （1）混合地图表示：体素地图 $V$ 以均匀划分的体素建模占用概率，用于碰撞检测和路径规划；GS 地图 $G$ 基于 Gaussian Surfel 表示，每个基元 $g_i$ 由位置 $\mu_i$、旋转 $q_i$、缩放 $s_i$、颜色 $c_i$、不透明度 $o_i$ 和置信度 $k_i$ 定义。% （2）增量建图：收到新 RGB-D 观测后，先渲染当前视角的颜色图、深度图和不透明度图，然后通过二值掩码 $B(p)=(O(p)<0.5)\wedge(\mathrm{avg}|I(p)-I^{*}(p)|>0.5)\wedge(D(p)-D^{*}(p)>\delta D^{*}(p))$ 识别需要增加新高斯的区域（不透明度低、颜色偏差大或出现新几何），将对应像素反投影为新的 Gaussian Surfel。训练损失为 $L = w_c L_c + w_d L_d + w_n L_n$，其中 $L_c$、$L_d$ 为颜色和深度的 $L_1$ 损失，$L_n$ 包含法向量余弦距离与全变差平滑项。每 5 步执行可见性检查，删除所有历史视角均不可见的高斯。% （3）置信度建模：每个高斯的置信度 $k_i = \sigma_i \exp(\rho_i)$ 由两部分组成——距离加权的余弦相似度 $\sigma_i = \sum_{j \in S(g_i)}(1 - d_{ij}/d_{\mathrm{far}})\,\hat{v}_{ij}^{\top}n_i$ 衡量观测视角与高斯法向的对齐程度（近距和正面观测权重更高），方向离散度 $\rho_i = 1 - \|\bar{v}_i\|$ 衡量观测方向的多样性（$\rho_i \to 0$ 表示观测角度趋同）。
\textbf{方法信息：} 代码：\url{https://github.com/dmar-bonn/active-gs}；GitHub 关注度高，约 150 stars；引用很多，Semantic Scholar 约 42 次。
% METHOD: 2503.07819
\paragraph{POp-GS（最优实验设计视角选择）（2025.03）}
POp-GS~\cite{2503_07819} 从最优实验设计角度重写 3DGS 的 NBV 问题。它不把候选视角看成“会看到多少未知区域”，而是看成一次新的测量：这次测量加入训练后，能把高斯参数的后验协方差压缩多少。这样视角选择就对应到参数估计中的信息增益。

具体流程是先从最大似然角度连接 3DGS 的重建损失和参数梯度，用 Gauss--Newton/Hessian 近似为每个高斯的参数构造协方差。为了在线计算，方法提供两种近似：对角近似忽略参数间相关性，计算便宜；块对角近似保留同一高斯内部参数的耦合，更接近真实不确定性。候选视角的得分采用 T-Optimality 或 D-Optimality：前者看协方差矩阵迹的减少，后者看 log determinant 的减少。相对 GauSS-MI 的可靠概率，POp-GS 更偏“参数估计理论”；相对 OUGS，它不依赖对象 mask，而是对全体高斯参数估计不确定性。
\textbf{实现细节：} POp-GS（2025.03） 从最优实验设计角度重新构建信息量化问题，提出 T-最优性和 D-最优性两种准则。FisherRF 系列方法（包括 Next Best Sense（2024.10） ）利用 Fisher 信息矩阵量化不确定性，并将其扩展到同时选择视角和触觉姿态。SA-ResGS（2026.01） 提出自我增强残差 3DGS，通过三角化训练视角和外推渲染视角生成增强点云来估计场景覆盖度，并提出首个针对 3DGS 的残差学习策略。
\textbf{方法信息：} Amazon 参与了这篇论文。
\textbf{设计动机与边界：} 相对原始 3DGS 以独立高斯和图像重建误差为主的设定，POp-GS 把问题落在以下可实现的改动上：方法分为三个部分：(1) 协方差矩阵近似：从最大似然角度出发，将3DGS的L1损失与参数梯度联系起来，构建每个椭球参数的协方差矩阵作为不确定性度量。提出对角近似（忽略参数间相关性，计算快）和块对角近似（考虑同一高斯内参数相关性，更准确但计算更昂贵）。(2) 信息增益计算：基于最优实验设计理论，定义T-Optimality（协方差矩阵的迹）和D-Optimality（协方差矩阵行列式的对数）作为信息增益准则。给定候选视角，计算该视角能减少的参数不确定性。(3) NBV选择：遍历候选视角，选择信息增益最大的视角作为下一个最佳视角。其设计目的不是单独增加一个模块，而是用上述表示、监督或渲染约束消除原流程在目标任务中的具体失配。首次将最优实验设计理论应用于3DGS的不确定性量化和信息增益估计，提出T-Optimality和D-Optimality两种准则，以及块对角协方差近似来捕捉参数相关性，为3DGS SLAM中的NBV选择提供理论框架。适用时还应保留论文给出的边界：块对角近似计算成本较高；候选视角的遍历搜索限制了实时应用；对动态场景的支持有限。
% METHOD: 2504.21067
\paragraph{GauSS-MI（香农互信息主动重建）（2025.04）}
GauSS-MI~\cite{2504_21067} 关注的是主动 3DGS 中最基础的问题：机器人不能只按体素覆盖率去探索，因为一个区域“看过”并不代表其中的高斯已经能稳定渲染。它把地图状态从确定性的高斯集合扩展成概率地图，为每个高斯维护可靠概率 $P(r^i)$，含义是该高斯当前用于渲染是否可信。每次新观测到来后，方法根据图像覆盖、深度覆盖、渲染颜色误差和深度误差构造逆传感器模型，再用二元贝叶斯滤波器更新高斯的对数几率。

视角评分不再直接数未知体素，而是计算候选视角和当前高斯可靠性之间的 Shannon mutual information。直观地说，一个视角如果能看到许多“可见但不可靠”的高斯，并且预期观测会显著改变这些高斯的可靠概率，它的信息增益就高；如果它只重复观察已经可靠的区域，得分就低。代码层面，该方法接入 ROS/MonoGS 风格的主动重建系统，候选视角由运动规划器生成，GauSS-MI 负责给这些候选视角排序。相对 GS-Planner 这类完整机器人系统，GauSS-MI 的核心贡献更集中在“如何用 3DGS 自身的不确定性定义 NBV 指标”。
\textbf{实现细节：} GauSS-MI（2025.04）。Baseline：此前 NBV 选择中的信息增益通常依赖体积渲染中的射线-体素相交计数或启发式的覆盖率指标，缺乏对高斯基元自身可靠性的概率建模，导致无法在 3DGS 的泼溅渲染框架下严格量化"当前地图还剩多少不确定性可被新视角消除"。改了什么：直觉上来说，GauSS-MI 的核心思路是为每个高斯赋予一个"可靠度"评分，然后利用信息论工具计算"从某个新视角观察后，整体不确定性能降低多少"。具体而言，GauSS-MI 为每个 3D 高斯引入二值随机变量 $r^{[i]}$ 表示"该高斯对渲染是否可靠"，构建概率模型 $P(r^{[i]})\in(0,1)$，并在此基础上推导出 Shannon 互信息的闭式表达式：Gaussian Splatting Shannon Mutual Information（GauSS-MI）。这使得信息增益的计算可完全在 3DGS 光栅化管线中高效完成。实现原理：系统分两步运行。
\textbf{实现补充：} 实现上还要核对输入、输出、参数、特征、初始化、训练损失、正则、渲染或推理接口，以及优化中的剪枝、分裂、采样、编码或解码步骤；这些抓手决定该方法到底改了 3DGS 管线的哪一环。
\textbf{方法信息：} 代码：\url{https://github.com/johannaxie/gauss-mi}。
\textbf{设计动机与边界：} 相对原始 3DGS 以独立高斯和图像重建误差为主的设定，GauSS-MI 把问题落在以下可实现的改动上：GauSS-MI方法包含以下核心组件：1）3DGS概率模型：为每个3D高斯定义可靠概率P(r\^{}i)，表示该高斯用于渲染的可靠性，使用对数几率表示并通过二元贝叶斯滤波器在每次新观测后更新概率；2）逆传感器模型：定义观测Z\_k下高斯可靠性的条件概率P(r\^{}i|Z\_k)，结合渲染误差和观测覆盖率；3）香农互信息准则GauSS-MI：计算新视角相对于当前3DGS地图的互信息，作为视角选择的信息增益指标；4）系统集成：将GauSS-MI与视图和运动规划器集成，构建完整的主动重建系统，机器人自主选择信息量最大的下一个最佳视角进行观测。初始假设无先验信息（P\_0=0.5），通过增量式贝叶斯更新逐步建立环境认知。其设计目的不是单独增加一个模块，而是用上述表示、监督或渲染约束消除原流程在目标任务中的具体失配。提出首个基于3DGS的视觉不确定性量化概率模型和实时视角选择准则GauSS-MI。利用香农互信息实时评估新视角的视觉信息量，结合视图和运动规划器构建主动重建系统，直接优化重建模型的视觉质量而非几何覆盖率。适用时还应保留论文给出的边界：计算互信息需要遍历所有高斯，在超大规模场景中可能存在效率瓶颈；对动态场景的适应性未充分讨论；依赖于准确的相机位姿估计。
% METHOD: 2511.09397
\paragraph{OUGS（对象感知不确定性）（2025.11）}
OUGS~\cite{2511_09397} 把主动重建从“把整个场景看完整”收缩到“把目标物体看清楚”。普通 scene-level uncertainty 容易被背景纹理、墙面或大面积未观测区域主导，机器人会把视角预算花在背景上；如果任务是重建桌上的目标物体，这种策略反而偏离目标。

OUGS 的做法是先给每个高斯的显式参数建立不确定性，包括位置、尺度、旋转、不透明度和 SH 颜色。训练时用梯度平方的指数滑动平均近似对角 Fisher 信息矩阵，由此得到每个参数的协方差；渲染候选视角时，再用可微 rasterizer 的雅可比把 3D 参数协方差传播到像素空间，得到颜色或外观不确定性图。最后，语义前景 mask 只保留目标区域的不确定性，候选视角按照 masked uncertainty 聚合得分。这样 OUGS 和 GauSS-MI 一样使用 3DGS 内部状态做 NBV，但它把不确定性限定到对象区域，更适合目标中心采集。
\textbf{实现细节：} 水下与海洋场景作为散射介质的极端形式，需对波长依赖衰减与焦散纹理进行物理建模。后向散射方面，RUSplatting 与 UW-3DGS 处理后向散射与颜色畸变；DualPhys-GS 与 Underwater Rethink 进一步对物理模型做反思性重构。RecGS 则专门处理水下焦散。与上节 OceanSplat 形成互补——前者侧重三目立体框架，后者侧重物理模型细化。
\textbf{实现补充：} 进一步拆开来看，输入侧对应图像、点云、位姿、语义提示或已训练高斯资产；输出侧对应新的高斯参数、渲染图像、编辑结果、地图状态或评测指标。表示层需要说明位置、尺度、旋转、不透明度、特征、材质或 MLP 解码器怎样改；训练层需要说明损失、正则、采样、剪枝、分裂或压缩何时发生；推理层则看渲染、查询、规划或系统接口如何接到原始 3DGS / Scaffold-GS / 4DGS 管线上。
\textbf{方法信息：} 代码：\url{https://github.com/Gatsby0916/OUGS}。
\textbf{设计动机与边界：} 相对原始 3DGS 以独立高斯和图像重建误差为主的设定，OUGS 把问题落在以下可实现的改动上：核心创新是将每个高斯的参数向量视为随机变量，初始化高斯先验，通过不同iable渲染管线的雅可比矩阵将3D参数协方差传播到像素空间，得到像素级视觉不确定性。可将总不确定性分解为几何和外观分量。在此物理基础之上，无缝集成语义分割掩码，生成目标感知不确定性分数，有效将目标从环境中分离。基于此不确定性选择最能提升目标保真度的视角。其设计目的不是单独增加一个模块，而是用上述表示、监督或渲染约束消除原流程在目标任务中的具体失配。首次从3D高斯原语的显式物理参数（位置、尺度、旋转）推导不确定性，通过渲染雅可比矩阵传播参数协方差，结合语义分割掩码生成目标感知不确定性分数，实现高效的目标中心主动视角选择。适用时还应保留论文给出的边界：依赖语义分割质量；不确定性分解假设几何与外观参数独立。
% METHOD: 2601.03024
\paragraph{SA-ResGS（自增强残差与 NBV）（2026.01）}
SA-ResGS~\cite{2601_03024} 的出发点是：稀疏主动重建里，视角选择和 3DGS 优化不是两个独立问题。宽基线候选视角可能提供新几何，但如果当前高斯在这些区域没有足够约束，普通渲染损失会给出很弱或很噪的梯度；反过来，如果只按当前不确定性选视角，又容易被早期错误地图误导。

它先从训练视图和外推视图之间做三角化，生成 self-augmented points，用这些点补充当前 SfM/GS 覆盖不到的空间线索。候选视角评分时，方法把场景写入哈希表征，比较候选视角和已选视角在几何上的差异，倾向选择能看到新结构的位置。训练侧则把普通 3DGS 损失改成 residual learning：渲染时强调欠监督高斯和可见稳定高斯之间的残差，通过不确定性过滤、dropout 采样和困难样本式采样，把原本被平均掉的弱梯度放大。相对 GauSS-MI，SA-ResGS 不只是给视角打分，还同时改变了稀疏视角下高斯被监督的方式。
\textbf{实现细节：} POp-GS（2025.03） 从最优实验设计角度重新构建信息量化问题，提出 T-最优性和 D-最优性两种准则。FisherRF 系列方法（包括 Next Best Sense（2024.10） ）利用 Fisher 信息矩阵量化不确定性，并将其扩展到同时选择视角和触觉姿态。SA-ResGS（2026.01） 提出自我增强残差 3DGS，通过三角化训练视角和外推渲染视角生成增强点云来估计场景覆盖度，并提出首个针对 3DGS 的残差学习策略。
\textbf{方法信息：} Huawei 参与了这篇论文。
\textbf{设计动机与边界：} 相对原始 3DGS 以独立高斯和图像重建误差为主的设定，SA-ResGS 把问题落在以下可实现的改动上：SA-ResGS框架包含三个核心组件：(1)物理引导的视图选择策略：通过三角化训练视图和光栅化外推视图生成SA-Points，利用基于哈希的场景表征编码，度量候选视图与已选视图的相似性，选择最具几何差异性的视图实现均匀场景覆盖；(2)不确定性感知残差监督：针对稀疏和宽基线视图导致的欠监督Gaussian问题，首次提出针对3DGS的残差学习策略，结合不确定性驱动过滤、dropout采样和困难负样本挖掘启发的采样方式，增强高不确定性Gaussian的梯度流动和训练稳定性；(3)隐式去偏：通过约束视图选择和残差监督的协同作用，隐式地减少不确定性量化中的偏差，缓解宽基线探索和稀疏视图歧义在NBV规划中的冲突效应。整体框架将视图选择与不确定性估计解耦，减少对内部学习动态的强依赖。其设计目的不是单独增加一个模块，而是用上述表示、监督或渲染约束消除原流程在目标任务中的具体失配。提出自增强残差3DGS框架(SA-ResGS)，通过三角化生成自增强点云(SA-Points)估计场景覆盖度，并首次引入针对3DGS的残差学习策略，结合不确定性驱动过滤和dropout采样增强弱监督Gaussian的梯度流动，在主动视图选择中实现更稳定的不确定性量化和更高质量的场景重建。适用时还应保留论文给出的边界：方法主要针对主动视图选择场景，需要预训练一定数量的初始视图后才能生成SA-Points；在极端稀疏视图场景下残差学习的效果可能受限。
% METHOD: 2601.07484
\paragraph{R3-RECON（基于可渲染性的主动重建）（2026.01）}
R3-RECON~\cite{2601_07484} 走了另一条路线：它不要求在线查询一个正在优化的 3DGS 或 NeRF，也不依赖当前高斯参数的梯度状态，而是估计“从某个方向看某个空间位置时，未来辐射场能不能可靠渲染”。这个量被称为 renderability，可理解为采集历史对未来视角的可外推程度。

方法维护轻量体素统计地图。每个体素记录历史观察方向和历史颜色/辐射观测的分散度：查询方向与历史方向差距越大，观察偏差越高；同一体素历史观测越不一致，观察噪声越高。二者结合得到闭式可渲染性分数，用来快速评估候选视角。由于分数来自观测统计而非 3DGS 当前优化状态，它可以在毫秒级查询，也能扩展到全景候选视点评估。相对 ActiveGS 这类直接在高斯地图上渲染置信度的做法，R3-RECON 更像一个独立的采集质量预估器，可为后续 3DGS 重建提前安排视角。
\textbf{实现细节：} $\mathbb{R}^{}3$-RECON（2026.01） 提出了无需辐射场的主动重建框架，从轻量体素地图诱导隐式的姿态条件渲染性场，在毫秒级即可查询候选视角的渲染性分数。ActiveInitSplat（2025.03） 专注于图像选择对高斯初始化的影响，基于密度和占据准则确保训练图像来自多样化视角。
\textbf{技术实现补充：} $\mathbb{R}^3$-RECON（2026.01） 提出了无需辐射场的主动重建框架，从轻量体素地图诱导隐式的姿态条件渲染性场，在毫秒级即可查询候选视角的渲染性分数。ActiveInitSplat（2025.03） 专注于图像选择对高斯初始化的影响，基于密度和占据准则确保训练图像来自多样化视角。
\textbf{方法信息：} 代码：\url{https://github.com/jkff00/r3recon}。
\textbf{设计动机与边界：} 相对原始 3DGS 以独立高斯和图像重建误差为主的设定，R3-RECON 把问题落在以下可实现的改动上：核心方法：(1)辐射表示中的可渲染性分析——从球谐(SH)拟合行为角度分析基元级可渲染性，通过观察偏差（查询方向与历史观察方向的角偏差）和观察噪声（历史辐射观察的分散度）两个因子表征拟合质量。利用球面几何分析外推区域中误差放大的影响；(2)观察统计驱动的可渲染性——维护轻量级体素统计地图，在线聚合每个体素的方向观察统计，推导出闭式标量可渲染性分数。可渲染性与图像空间重建误差强相关，自然引导NBV选择；(3)全景扩展——估计全向(360度)视图效用加速候选评估，通过一次性评估完整视点的全方位价值避免逐方向评估的冗余；(4)主动重建流程——将可渲染性评分集成到探索框架中，支持RGB-D相机在线增量采集。方法完全解耦视图效用与瞬态优化状态。其设计目的不是单独增加一个模块，而是用上述表示、监督或渲染约束消除原流程在目标任务中的具体失配。提出R3-RECON，从轻量级体素地图导出隐式、位姿条件的可渲染性场，通过在线观察统计计算标量可渲染性分数，毫秒级查询候选视点，无需梯度或辐射场训练，实现比基于3DGS内部状态的NBV方法更均匀的新视角质量。适用时还应保留论文给出的边界：可渲染性模型假设Lambertian或弱非Lambertian场景；全景评估在复杂几何环境中可能不够精确。

\section{鲁棒重建与干扰物处理}

开放场景中，训练图像可能包含行人、车辆、风吹树叶、反光变化或临时遮挡。原始 3DGS 会把这些不一致内容吸收到静态高斯里，或者 densification 时为干扰物长出额外高斯。鲁棒方法需要先发现“哪些像素不属于静态场景”，再避免它们影响几何和颜色。
% METHOD: 2406.20055
\paragraph{SpotlessSplats（语义特征干扰物检测）（2024.06）}
SpotlessSplats~\cite{2406_20055} 是干扰物鲁棒 3DGS 的代表性锚点。传统鲁棒 NeRF 往往用颜色残差发现 transient pixels，但 3DGS 有自适应 densification：局部残差大时，优化器可能继续 clone 或 split 高斯，颜色残差的方差被训练过程放大，直接按 RGB 残差筛像素容易误判。SpotlessSplats 的核心改变，是把干扰物检测从 RGB 空间移到语义特征空间，让 mask 由更稳定的高层语义而不是逐像素颜色差来决定。

具体流程先对每张训练图预计算 Stable Diffusion 特征图，再从这些特征中估计 inlier/outlier mask。SLS-agg 对特征图做空间连通的层次聚类，把图像划成若干语义一致的组件，并根据组件与静态场景的兼容性估计内点概率；SLS-mlp 则训练轻量 MLP，把扩散特征映射为 inlier probability，并用高置信残差标签和 Lipschitz 正则交替优化 MLP 与 3DGS。得到的 mask 进入鲁棒重建损失，被判为 distractor 的像素不再驱动颜色、opacity 和 densification，因此不会继续诱导模型为临时物体长出高斯。方法还提出 utilization-based pruning，用高斯被实际使用的程度删除冗余 splats；公开实现中以 \texttt{--ubp} 开关接入。相对后面的 PDF-GS，SpotlessSplats 更像“先用语义特征找干扰物，再训练干净场景”的路线；它也成为后续鲁棒 3DGS 方法比较 mask 来源和 densification 耦合时常用的参照。
\textbf{实现细节：} 核心创新是在丰富的语义特征空间而非RGB空间中检测异常值。使用Stable Diffusion提取的特征图作为预计算步骤。方法一（空间聚类）：对特征图进行无监督的层次聚类，将图像过分割为固定数量的空间连通组件。通过计算每个聚类的内点概率来识别干扰物聚类。方法二（学习方法）：训练一个小MLP以无监督方式预测可能对应干扰物的特征子空间。自适应鲁棒损失：基于特征空间的干扰物掩码替代手工设计的鲁棒核函数。稀疏化策略：与鲁棒优化兼容，通过正则化减少冗余高斯。该策略是通用的，即使在无干扰物的数据集上也能减少2-4倍高斯数量而不损失质量。关键解决了3DGS自适应密集化过程中颜色残差方差大导致传统基于残差的干扰物检测失效的问题。 首次将文本到图像扩散模型的语义特征用于3DGS干扰物检测，替代基于颜色残差的方法，解决了3DGS自适应密集化过程引入颜色残差方差干扰干扰物检测的问题。附带稀疏化策略可在无干扰物数据集上也减少2-4倍高斯数量。 依赖预训练扩散模型的质量和特征提取。聚类数量和超参数需要调整。对大面积干扰物的处理仍有挑战。
\textbf{设计动机与边界：} 相对原始 3DGS 以独立高斯和图像重建误差为主的设定，SpotlessSplats 把问题落在以下可实现的改动上：核心创新是在丰富的语义特征空间而非RGB空间中检测异常值。使用Stable Diffusion提取的特征图作为预计算步骤。方法一（空间聚类）：对特征图进行无监督的层次聚类，将图像过分割为固定数量的空间连通组件。通过计算每个聚类的内点概率来识别干扰物聚类。方法二（学习方法）：训练一个小MLP以无监督方式预测可能对应干扰物的特征子空间。自适应鲁棒损失：基于特征空间的干扰物掩码替代手工设计的鲁棒核函数。稀疏化策略：与鲁棒优化兼容，通过正则化减少冗余高斯。该策略是通用的，即使在无干扰物的数据集上也能减少2-4倍高斯数量而不损失质量。关键解决了3DGS自适应密集化过程中颜色残差方差大导致传统基于残差的干扰物检测失效的问题。其设计目的不是单独增加一个模块，而是用上述表示、监督或渲染约束消除原流程在目标任务中的具体失配。首次将文本到图像扩散模型的语义特征用于3DGS干扰物检测，替代基于颜色残差的方法，解决了3DGS自适应密集化过程引入颜色残差方差干扰干扰物检测的问题。附带稀疏化策略可在无干扰物数据集上也减少2-4倍高斯数量。适用时还应保留论文给出的边界：依赖预训练扩散模型的质量和特征提取。聚类数量和超参数需要调整。对大面积干扰物的处理仍有挑战。
\textbf{方法信息：} Google / DeepMind 参与了这篇论文；代码：\url{https://github.com/lilygoli/SpotLessSplats}；GitHub 关注度高，约 247 stars；引用很多，Semantic Scholar 约 68 次。



% METHOD: 2505.15737
\paragraph{RUSplatting（水下稀疏鲁棒重建）（2025.05）}
RUSplatting~\cite{2505_15737} 面向稀疏水下场景。这里的干扰不是行人或车辆，而是水体介质本身：不同颜色通道衰减速度不同，后向散射会给远处物体叠加雾化颜色，少量视角又会让 COLMAP 点云和 3DGS 初始化都不稳定。若直接套用原始 3DGS，模型往往把颜色偏移、散射和真实纹理一起写进球谐颜色，得到的高斯场在新视角下容易发灰、发蓝或边缘漂移。

RUSplatting 的第一项改变，是在 3DGS 颜色预测外接一个水下成像分解模块。五层 MLP 以深度和视角编码为输入，分别为 R/G/B 通道预测直射透射率、后向散射透射率、衰减系数和背景光偏置，再把高斯渲染出的直接颜色与介质项组合成最终像素。这样，颜色校正不再只是后处理，而是参与训练过程。第二项改变是中间帧插值：方法用 RIFE 生成补充视角，把这些帧用于改善 SfM 点和训练覆盖；自适应帧权重为插值帧学习不确定性，避免低质量插值破坏真实观测。第三项改变是边缘感知平滑损失，用 TV 与 Sobel 约束水下噪声，同时保留物体轮廓。相对 Naka-GS 处理低光照曝光不足，RUSplatting 处理的是水体吸收与散射造成的成像退化。
\textbf{实现细节：} 水下与海洋场景作为散射介质的极端形式，需对波长依赖衰减与焦散纹理进行物理建模。后向散射方面，RUSplatting 与 UW-3DGS 处理后向散射与颜色畸变；DualPhys-GS 与 Underwater Rethink 进一步对物理模型做反思性重构。RecGS 则专门处理水下焦散。与上节 OceanSplat 形成互补——前者侧重三目立体框架，后者侧重物理模型细化。
\textbf{方法信息：} 代码：\url{https://github.com/theflash987/RUSplatting}。



% METHOD: 2604.12580
\paragraph{PDF-GS（渐进式干扰物过滤）（2026.04）}
PDF-GS~\cite{2604_12580} 利用 3DGS 的自过滤现象：多视角不一致区域在优化中往往难以稳定解释，训练到一定阶段后，模型渲染与输入图像之间的差异会暴露行人、车辆或临时遮挡。它没有额外依赖扩散语义特征，而是把“3DGS 自己解释不了的区域”变成 mask 线索，再通过多阶段训练逐步过滤干扰物。

流程分为 progressive filtering phases 和 final reconstruction phase。第一个阶段先训练一个粗 3DGS，随后每一轮都用上一阶段模型渲染训练视图，与原图在颜色和特征空间比较得到 discrepancy map；高差异区域被视为 distractor mask。关键设计是每个过滤阶段不沿用上一轮已经被污染的高斯，而是从 SfM 点重新初始化新的高斯参数，再用 mask 后的图像和结构保持损失训练，颜色更新则更稀疏，以避免早期错误 mask 把真实背景也抹掉。最后的重建阶段使用过滤后的干净视图恢复细节。相对 SpotlessSplats，PDF-GS 的优势是 pipeline 更贴近标准 3DGS 训练，主要增加阶段调度、差异图生成和重新初始化；代价是结果依赖自过滤是否能在早期把干扰物和静态背景拉开。
\textbf{实现细节：} PDF-GS（2026.04） 放大了 3DGS 固有的自过滤特性，通过渐进式多阶段优化逐步移除干扰物：先利用差异线索逐步过滤，再从净化后的高斯表示恢复精细细节。与 DeSplat 的显式干扰物建模不同，PDF-GS 依赖 3DGS 自身的过滤能力来渐进剔除瞬态元素。
\textbf{方法信息：} 代码：\url{https://github.com/kangrnin/PDF-GS}。
\textbf{设计动机与边界：} 相对原始 3DGS 以独立高斯和图像重建误差为主的设定，PDF-GS 把问题落在以下可实现的改动上：PDF-GS基于一个关键观察：3DGS天然倾向于抑制多视角不一致内容（自过滤现象）。框架分为两个阶段：（1）渐进过滤阶段（phase 1到K-1）：从初始高斯参数集θ\^{}(1)开始，每个阶段通过前阶段渲染输出与训练图像的差异定位干扰物区域，生成掩码后重新从SfM点初始化高斯参数集θ\^{}(k)并用掩码图像优化。重新初始化而非复用前阶段高斯，防止误差累积（如颜色漂移、几何退化）。使用加权SSIM和L1损失组合；（2）重建阶段（phase K）：使用纯净的训练图像恢复被过滤阶段衰减的精细细节。整个过程轻量级，可直接集成到现有3DGS框架中。其设计目的不是单独增加一个模块，而是用上述表示、监督或渲染约束消除原流程在目标任务中的具体失配。提出PDF-GS框架，利用3DGS固有的自过滤特性（倾向于抑制多视角不一致区域），通过渐进式多阶段优化逐步去除干扰物。每个过滤阶段通过差异图识别干扰物区域并重新初始化高斯参数进行优化，最后通过重建阶段恢复精细细节。无需架构修改或额外推理开销。适用时还应保留论文给出的边界：渐进式优化增加了训练时间（多个阶段）；依赖SfM点初始化质量；极端密集干扰物场景可能需要更多过滤阶段。
% METHOD: 2604.21631
\paragraph{DualSplat（失败到先验的瞬态物体处理）（2026.04）}
DualSplat~\cite{2604_21631} 提出 failure-to-prior 范式，解决瞬态物体检测和干净重建之间的循环依赖：没有干净场景，残差和语义线索不够可靠；没有瞬态 mask，3DGS 又会把临时物体写进静态场景。它的做法不是反复过滤图像，而是先让一轮保守重建“失败得可见”，再把这些失败痕迹转换为第二轮重建的先验。

第一阶段只做保守 3DGS 重建。由于瞬态物体在不同视角中可见性不一致，它们会在渲染结果里表现为不完整碎片、背景泄漏和局部残差。DualSplat 同时使用三类线索生成对象级伪 mask：渲染图与输入图的 photometric residual，DINO 特征中的语义不匹配，以及 SAM2 给出的实例边界；SAM2 的 object masks 用来把像素残差聚合到实例层面，再通过自适应阈值确定哪些实例更像 transient。第二阶段重新训练干净 3DGS，先把这些伪 mask 当作硬先验屏蔽瞬态区域，再用轻量 MLP 在线细化 mask，让监督逐渐从外部先验转向当前模型的自一致性。相对 PDF-GS，它更重视对象级 mask 和多模态线索；相对 SpotlessSplats，它不把语义特征作为唯一判断，而是把第一轮重建失败本身也纳入证据。
\textbf{实现细节：} DualSplat（2026.04） 提出"失败到先验"框架：将首轮保守训练中的重建失败（瞬态物体通常表现为不完整碎片）转化为第二轮重建的显式先验，结合光度残差、特征不匹配和 SAM2 实例边界构建对象级伪掩码。该方法的核心洞察是瞬态物体仅出现在部分视角中，在保守初始训练中往往表现为不完整碎片，这些失败本身恰好可以转化为干扰物的显式先验；一个轻量级 MLP 在线精修伪掩码，逐渐从先验监督转向自一致性约束。
\textbf{方法信息：} 代码：\url{https://github.com/Lans1ot/DualSplat}。
\textbf{设计动机与边界：} 相对原始 3DGS 以独立高斯和图像重建误差为主的设定，DualSplat 把问题落在以下可实现的改动上：DualSplat的核心是两阶段Failure-to-Prior流程：(1)第一阶段（失败暴露）：进行保守的3DGS重建，瞬态物体因跨视角不一致可见性而产生不完整碎片；结合三种线索构建伪掩码——光度残差（渲染图与输入的差异）、特征不匹配（使用DINO特征检测语义不一致）、SAM2实例边界（获取对象级分割）；通过局部掩码策略促进碎片检测；(2)第二阶段（先验指导重建）：以伪掩码为硬先验掩蔽瞬态区域进行干净的3DGS优化；引入轻量级MLP在线细化掩码，从先验监督逐渐过渡到自一致性监督，随着场景几何稳定化自动调整。关键创新是解耦瞬态发现与干净场景优化。其设计目的不是单独增加一个模块，而是用上述表示、监督或渲染约束消除原流程在目标任务中的具体失配。提出创新的'失败到先验'范式DualSplat，将首轮重建失败转化为显式先验指导二次重建，打破瞬态检测与场景重建的循环依赖，在瞬态密集场景中显著优于现有方法。适用时还应保留论文给出的边界：两阶段流程增加总训练时间；第一阶段保守训练的质量影响伪掩码精度；对非常小的瞬态物体可能检测不到。
% METHOD: 2606.29496
\paragraph{Entropy Mask Rectification（熵引导免干扰重建）（2026.06）}
Rectifying Mask via Entropy for Distractor-Free 3DGS~\cite{2606_29496} 也称 RefineSplat，处理的是比普通 transient removal 更难的歧义干扰物：行人、树叶或临时物体可能与静态背景颜色相近，语义也不总是清楚，单靠光度残差或一次性 SAM mask 容易把静态结构误删，或把干扰物留在高斯场里。它的改变，是把 mask 置信度本身纳入建模，用 entropy 判断哪些实例区域处在“不确定、容易误判”的状态，再把这个不确定性接到 density control。

流程先用冻结的实例分割模型得到每张训练图中的候选 instance masks，并用 DINOv2 等语义特征训练小 MLP 预测实例身份。局部 mask 分支计算每个实例的平均 entropy：如果某个实例在不同视图、语义特征或预测中表现出高不确定性，就更可能是 ambiguous distractor。全局分支再利用 entropy 与语义特征捕获细粒度、边界不完整或分割错误的干扰区域。最终得到的 mask 不是固定阈值结果，而是由局部实例不确定性和全局语义不确定性共同修正。

它还把 entropy 用进高斯增长。原始 3DGS densification 主要看位置梯度，干扰物区域如果颜色残差大，反而会诱导模型在那里继续 clone/split 高斯。RefineSplat 改成 entropy-aware positional gradients：高不确定区域的梯度被重新加权，让 densification 更少追随 transient artifact，同时通过合并重叠高斯把静态结构对齐到实例级别。相对 SpotlessSplats 的扩散语义特征异常检测，它更强调“mask 错误怎样反馈到高斯密度控制”；相对 PDF-GS 的渐进过滤，它不只是重启训练，而是在训练中持续修正 mask 与 densification 的耦合。
\textbf{实现细节：} Rectifying Mask via Entropy for Distractor-Free 3DGS in Ambiguous Scenarios。Rectifying Mask via Entropy for Distractor-Free 3DGS in Ambiguous Scenarios（2026.06） 延续 3DGS 在"Distractor-robust Gaussian reconstruction"上的基本设定，主要改变是：基于熵的自适应 transient mask 与 density control。3DGS 的显式高斯表示和可微渲染能力保持为基础。其设计动机是：针对颜色或语义歧义下难以区分 transient 元素和静态场景的问题，提出系统化 mask 构建和密度控制。具体而言，结合熵信号和实例 mask 捕获模糊干扰物，并用 entropy-aware positional gradients 调整 Gaussian density control，使高斯更好对齐真实静态结构。适用边界是：依赖熵估计和实例 mask 质量，对长期静止的动态物体或语义错误仍可能失效。
\textbf{技术实现补充：} Rectifying Mask via Entropy for Distractor-Free 3DGS in Ambiguous Scenarios。Rectifying Mask via Entropy for Distractor-Free 3DGS in Ambiguous Scenarios（2026.06） 延续 3DGS 在"Distractor-robust Gaussian reconstruction"上的基本设定，主要改变是：基于熵的自适应 transient mask 与 density control。3DGS 的显式高斯表示和可微渲染能力保持为基础。其设计动机是：针对颜色或语义歧义下难以区分 transient 元素和静态场景的问题，提出系统化 mask 构建和密度控制。具体而言，结合熵信号和实例 mask 捕获模糊干扰物，并用 entropy-aware positional gradients 调整 Gaussian density control，使高斯更好对齐真实静态结构。适用边界是：依赖熵估计和实例 mask 质量，对长期静止的动态物体或语义错误仍可能失效。
\textbf{方法信息：} Samsung 参与了这篇论文。
\textbf{设计动机与边界：} 相对原始 3DGS 以独立高斯和图像重建误差为主的设定，Rectifying Mask via Entropy for Distractor-Free 3DGS in Ambiguous Scenarios 把问题落在以下可实现的改动上：结合熵信号和实例 mask 捕获模糊干扰物，并用 entropy-aware positional gradients 调整 Gaussian density control，使高斯更好对齐真实静态结构。其设计目的不是单独增加一个模块，而是用上述表示、监督或渲染约束消除原流程在目标任务中的具体失配。针对颜色或语义歧义下难以区分 transient 元素和静态场景的问题，提出系统化 mask 构建和密度控制。适用时还应保留论文给出的边界：依赖熵估计和实例 mask 质量，对长期静止的动态物体或语义错误仍可能失效。

\section{深度、立体、超分与补充先验}

本节收纳与稀疏鲁棒相关但任务更专门的方法：有的用免微调扩散模型修复外推视角，有的研究低分辨率输入的 3DGS 超分，有的用 2D 基础模型和立体先验增强初始化。这些方法的外部信息不一定充分，但都补全了“输入不完美时怎样给 3DGS 可靠先验”的图谱。多光谱 3DGS 方向已有项目页和仓库占位，但完整代码尚未释放，适合放入后续持续关注清单，而不是作为本版正文可信方法展开。
% METHOD: 2311.13398
\paragraph{Depth-Regularized Optimization for 3D Gaussian Splatting in Few-Shot Images（稀疏视角、鲁棒性与主动观测）（2023.11）}
Depth-Regularized Optimization for 3D Gaussian Splatting in Few-Shot Images~\cite{2311_13398} 关注的问题是：在稀疏视角、少样本或鲁棒条件下重建 3DGS Depth-Regularized Optimization for 3D Gaussian Splatting in Few-Shot Images 是宽召回补全中发现的论文，主要根据标题、摘要和 arXiv 类别判定其与 3DGS 的关系。
相对原始 3DGS，Depth-Regularized Optimization for 3D Gaussian Splatting in Few-Shot Images 的变化不是只替换损失函数，而是把表示、初始化、约束或网络预测路径重新组织到任务流程中：该工作被召回是因为论文文本包含 Gaussian/3D Gaussian 相关表述；当前分类先确定其任务域和与 3DGS 的相关性，后续证据流水线再补代码、企业和引用信息。
\textbf{实现细节：} 具体实现可以按输入、核心模块和渲染/优化路径来读。输入端首先取得任务所需的图像、相机、文本、时间序列、语义或几何线索，并把这些线索转换成可以约束高斯中心、尺度、旋转、不透明度、颜色或附加特征的条件。核心结构中，方法使用的主要机制是：从摘要层面看，Depth-Regularized Optimization for 3D Gaussian Splatting in Few-Shot Images 围绕“在稀疏视角、少样本或鲁棒条件下重建 3DGS”展开。若其正文或代码后续可读，应进一步核对它具体改动的是表示、初始化、渲染、优化、损失、动态建模、压缩或应用接口中的哪一环。 训练或推理时，这些模块通常先给出初始化或中间特征，再通过可微 Gaussian rasterization、重建损失、感知/语义约束、几何正则或任务损失回传到高斯参数；如果方法包含前馈网络，则网络输出的是高斯属性、变形场、压缩码、语义特征或控制变量，而不是只输出最终 2D 图像。渲染路径仍以显式高斯投影、深度排序和 alpha 合成为基础，只是在采样、属性预测、局部编辑、动态变形、压缩编码或下游任务接口处加入额外约束。这样写入本章时，它与前后方法的关系是：相对原始 3DGS 和本章已有锚点方法，它可以理解为 3DGS extends\_or\_uses 宽召回分类阶段的关系判断，后续需要按论文方法部分复核。
\textbf{证据与边界：} 本轮召回分类尚未逐项核对实验表格；后续证据补全阶段需要读取论文源码/PDF或项目页确认数据集、指标和对比方法。 该条来自宽召回批次；源码下载受 arXiv e-print 网络限制，部分论文暂时只能依据标题、摘要和类别分类，需要在入可信书前做人工/LLM精读复核。 本书只采用本地 paper pack、数据库分类、源码/引用登记和论文摘要中可核验的信息；没有本地证据支撑的速度、排行榜或开源完整性不在正文中强化。\textbf{方法信息：} 代码：\url{https://github.com/robot0321/DepthRegularizedGS}；引用信号：242 citations (7.7/month, domain-relative cutoff 2.2/month, semantic\_scholar)；代码关注：robot0321/depthregularizedgs has 275 stars (8.7/month, domain-relative cutoff 11.2/month)。
\textbf{设计动机与边界：} 相对原始 3DGS 以独立高斯和图像重建误差为主的设定，Depth-Regularized Optimization for 3D Gaussian Splatting in Few-Shot Images 把问题落在以下可实现的改动上：从摘要层面看，Depth-Regularized Optimization for 3D Gaussian Splatting in Few-Shot Images 围绕“在稀疏视角、少样本或鲁棒条件下重建 3DGS”展开。若其正文或代码后续可读，应进一步核对它具体改动的是表示、初始化、渲染、优化、损失、动态建模、压缩或应用接口中的哪一环。其设计目的不是单独增加一个模块，而是用上述表示、监督或渲染约束消除原流程在目标任务中的具体失配。该工作被召回是因为论文文本包含 Gaussian/3D Gaussian 相关表述；当前分类先确定其任务域和与 3DGS 的相关性，后续证据流水线再补代码、企业和引用信息。适用时还应保留论文给出的边界：该条来自宽召回批次；源码下载受 arXiv e-print 网络限制，部分论文暂时只能依据标题、摘要和类别分类，需要在入可信书前做人工/LLM精读复核。
% METHOD: 2404.10318
\paragraph{SRGS（3DGS 超分辨率统一框架）（2024.04）}
SRGS~\cite{2404_10318} 面向低分辨率多视角输入。原始 3DGS 假设训练图像本身有足够分辨率；如果输入是下采样图像，直接在高分辨率空间优化会缺少纹理监督，单独对训练图做 2D 超分又容易生成视角不一致的高频细节。SRGS 的核心作用，是把 3DGS 超分辨率拆成两个可复用模块：外部 2D SR 先验负责补高频，内部跨视角一致性负责约束这些高频必须能被同一个高斯场解释。

在训练上，SRGS 先用固定的 2D SR 模型（如 SwinIR、StableSR 或 Real-ESRGAN 类模型）把每个低分辨率训练视图变成 pseudo-HR target，再让 HR 3DGS 渲染去拟合这些目标；这是 prior term。与此同时，系统把 HR 渲染下采样回输入分辨率，与原始 LR 图像计算一致性损失；这是 regularization term，用来防止 2D SR 先验在不同视图中注入互相冲突的细节。统一写法可以理解为 $L=\lambda_e L_{\mathrm{prior}} + (1-\lambda_e)L_{\mathrm{reg}}$。后续 3DGS SR 方法大多沿三条线改它：换更强的 2D SR 先验，增强跨视角几何正则，或者对两类损失做选择性加权和 mask。相对 FreeFix 的外推视角修复，SRGS 关注的是“训练图分辨率不足”这一输入退化。
\textbf{实现细节：} 针对此前 3DGS 超分方法多采用单图像超分范式、直接对渲染低分辨率图像进行 2D 上采样，忽略了高斯表征中固有的多视角一致性约束、导致跨视角纹理不一致和棋盘伪影的局限，SRGS（2024.04） 提出将 3DGS 超分辨率统一为先验注入和跨视角正则化的模块化框架。MVGSR（2025.12） 首次引入核线约束的多视角注意力机制到 3DGS 超分辨率，通过基于相机姿态的辅助视角选择适配任意组织的多视角数据集。IE-SRGS（2025.11） 联合利用外部 2D 超分辨率先验和内部 3DGS 特征，通过掩码引导融合策略协同利用两者互补优势。Close-up-GS（2025.03） 通过渐进式自训练和 3D 感知生成模型增强近距视角的细节渲染。
\textbf{方法信息：} 代码：\url{https://github.com/xiangfeng66/sr-gs}。



% METHOD: 2410.11080
\paragraph{Depth-Aware Few-shot 3DGS（深度先验与低阶球谐）（2024.10）}
Few-shot Novel View Synthesis using Depth Aware 3D Gaussian Splatting~\cite{2410_11080} 是少视角 3DGS 中一条轻量深度先验路线。原始 3DGS 在 5 个左右视角下容易只拟合训练图像颜色，而没有恢复正确 3D 形状；同时，高阶球谐给了模型过强的视角相关颜色容量，容易把错误几何伪装成训练视角上看起来正确的纹理。

它的做法有三处。第一，用 GLPN 等预训练单目深度模型为训练图生成 depth prior，并改造 rasterizer 额外渲染深度图；训练时用尺度不变深度损失约束渲染深度与单目深度的 Pearson 相关性，而不是强行匹配绝对尺度。第二，把球谐阶数降到 1，减少少视角条件下颜色分支的过拟合能力，让模型更依赖几何而不是视角相关颜色补偿。第三，取消原始 3DGS 中周期性删除低透明度高斯的策略，因为在少视角下这些低 opacity splats 可能仍是稀疏点云扩展所需的几何种子。相对 FSGS 通过虚拟视图扩充监督，这个方法更像是对标准 3DGS 训练环节做最小改动：增加深度渲染损失、降低颜色容量、放缓剪枝。
\textbf{方法信息：} 代码：\url{https://github.com/raja-kumar/depth-aware-3dgs}。
\textbf{设计动机与边界：} 相对原始 3DGS 以独立高斯和图像重建误差为主的设定，Few-shot Novel View Synthesis using Depth Aware 3D Gaussian Splatting 把问题落在以下可实现的改动上：方法包含三个关键改进：（1）深度先验：使用预训练单目深度估计网络获取深度先验，引入尺度不变深度损失（Pearson相关性）约束渲染深度与估计深度的一致性，解决少样本下3D形状不准确的问题。（2）低阶球谐函数：原始3DGS使用高阶球谐函数（最高4阶）表示颜色以恢复高频细节，但在少样本条件下容易过拟合训练视角。改用低阶球谐函数减少参数量。（3）保留高斯策略：观察到原始方法周期性移除低透明度高斯会导致非常稀疏的点云，因此在少样本设置中保留所有高斯以维持更好的重建。其设计目的不是单独增加一个模块，而是用上述表示、监督或渲染约束消除原流程在目标任务中的具体失配。提出深度感知的少样本3DGS方法，通过单目深度先验和低阶球谐函数避免过拟合，在PSNR上提升10.5%，SSIM提升6%，LPIPS提升14.1%。适用时还应保留论文给出的边界：依赖单目深度估计质量；对大尺度场景的效果未验证；低阶球谐函数可能丢失视角相关效果
% METHOD: 2510.14270
\paragraph{GauSSmart（2D 基础模型与几何滤波）（2025.10）}
GauSSmart~\cite{2510_14270} 结合 2D foundation models 和几何滤波增强 3DGS。它针对的不是单一深度先验，而是标准 COLMAP 初始化常见的两个问题：点云里有远离真实表面的伪点，且稀疏区域没有足够高斯支撑。若直接进入 3DGS densification，伪点会带来噪声，高斯增长也可能继续偏向已有纹理强的区域。

方法先做 convex-guided outlier removal：构建点云凸包，并按点到凸包的距离过滤离群点。与 kNN 半径过滤不同，它不要求局部邻域密度均匀，因此更适合稀疏且密度不均的 COLMAP 点云。第二步用相机聚类选择少量代表视角运行 SAM/SAM2，避免对所有图像做分割；再把 3D 点投影到这些视图的 segment masks 中，为点分配稳定的 segment id。对于点数不足的 segment，系统按原始点的协方差结构采样新点，执行 segment-aware densification，补足欠表示区域。第三步在训练时加入 DINOv3 embedding alignment loss，把渲染结果的高层语义特征拉向输入图像的 DINO 特征。相对 CoMapGS 的共可见性约束，GauSSmart 更偏向“用 2D 基础模型清理初始化、补点、再监督语义一致性”的路线。
\textbf{实现细节：} Depth Estimation Based on 3DGS（2024.09） 利用 3DGS 的可微渲染来辅助单目深度估计的自监督训练。Enhancing Sparse-view（2025.01） 通过 DINO-ViT 特征实现多视角语义一致性，并提出局部深度正则化约束深度值。Multimodal-Prior-Guided（2026.03） 融合光度渲染残差、语义先验和几何先验来驱动分层高斯致密化。GauSSmart（2025.10） 桥接 2D 基础模型和 3DGS，利用 DINO 语义特征引导高斯致密化和精修。 % ======================================================================。
\textbf{方法信息：} 代码：\url{https://github.com/alevalve/gaussmart}。
\textbf{设计动机与边界：} 相对原始 3DGS 以独立高斯和图像重建误差为主的设定，GauSSmart 把问题落在以下可实现的改动上：GauSSmart包含三个核心组件。第一，凸包离群值移除：针对COLMAP点云中的伪点，提出基于凸包的距离阈值滤波方法，独立评估每个点与凸包的距离进行移除，不依赖局部邻域，避免了传统kNN方法在密度不均场景下的局限性。第二，分割感知点云稠密化：利用SAM生成的分割掩码，通过协方差驱动的采样策略在欠表示区域增加点密度。具体地，在每个分割掩码内，根据原始点的协方差结构采样新点，同时考虑局部几何和分割面积。还提出相机聚类和图像选择策略，通过k-means在相机位置上聚类，选择紧凑的几何感知视角子集以降低分割计算量。第三，DINOv3嵌入对齐训练：引入基于DINOv3特征的轻量级嵌入损失，确保渲染的高斯外观与分割级别语义一致，提升空间保真度和对象连贯性。其设计目的不是单独增加一个模块，而是用上述表示、监督或渲染约束消除原流程在目标任务中的具体失配。提出GauSSmart混合框架，通过凸包离群值移除、基于SAM分割的感知点云稠密化、以及基于DINOv3嵌入的语义对齐训练损失，有效桥接2D基础模型和3DGS重建，在稀疏覆盖区域提升覆盖度并保留精细结构细节。适用时还应保留论文给出的边界：对2D基础模型质量有依赖；相机聚类策略需要调参；在极大规模场景上的扩展性有待验证。
% METHOD: 2601.20857
\paragraph{FreeFix（免微调图像扩散外推修复）（2026.01）}
FreeFix~\cite{2601_20857} 处理的是稀疏或部分观测 3DGS 在外推视角上的 artifacts。GSFixer 这类方法可以用扩散先验修复外推视图，但通常要为任务微调扩散模型；FreeFix 的改变，是保留预训练图像扩散模型的泛化能力，只在 3DGS 侧做交错式更新。它也不直接换成视频扩散模型，因为长序列视频模型成本高，且时间下采样会破坏逐像素 guidance。

流程先训练一个普通 3DGS，再沿连续外推轨迹逐视角处理。对于当前外推相机，系统用当前高斯场渲染图像和额外属性，生成 per-pixel confidence mask；低置信区域交给 SDXL 或 Flux 这类图像扩散模型修复，高置信区域尽量保持原渲染，避免扩散模型改坏已可靠内容。修复后的图像不是孤立输出，而是加入 fixed image set，作为下一轮 3DGS 优化的监督；优化后的 3DGS 再用于下一外推视角。这样，前面已经修复的视图会影响后续视图，形成 2D 修复和 3D 场景更新的交错闭环。论文还提到在 3D refinement 中使用 affine transform 抑制扩散图像的颜色风格漂移，使高斯场不必追随每张修复图的局部色调变化。相对 GSFixer，FreeFix 的重点不是训练一个更适配 3D artifacts 的扩散模型，而是设计“免微调扩散 + 置信 mask + 3D 迭代回灌”的工作流。
\textbf{设计动机与边界：} 相对原始 3DGS 以独立高斯和图像重建误差为主的设定，FreeFix 把问题落在以下可实现的改动上：FreeFix的流程为：给定已训练的3DGS模型，采样外推视图并渲染2D图像，使用图像扩散模型(IDM)进行精炼，将精炼结果整合回3DGS。核心技术包括：(1) 2D-3D交错精炼策略：交替进行2D图像精炼和3DGS模型更新，通过迭代逐步提升外推视图质量。选择IDM而非VDM（视频扩散模型）是因为VDM计算复杂度高（注意力二次复杂度）、输出长度受限、3D VAE的时间下采样阻碍逐像素引导。(2) 逐像素置信度掩码：通过分析3DGS渲染过程中的不确定性，生成逐像素掩码识别需要精炼的区域，将扩散模型引导集中在不确定区域，避免过度精炼已高质量区域。方法保留了扩散模型的泛化能力，同时逼近微调方法的保真度。其设计目的不是单独增加一个模块，而是用上述表示、监督或渲染约束消除原流程在目标任务中的具体失配。提出FreeFix，在无需微调扩散模型的前提下，利用预训练图像扩散模型增强3DGS外推视图渲染。核心创新包括：(1) 交错式2D-3D精炼策略，证明图像扩散模型可在无需视频扩散模型的情况下实现一致性精炼；(2) 逐像素置信度掩码识别不确定区域进行针对性改进，在多数据集上达到或超越微调方法的性能。适用时还应保留论文给出的边界：交错迭代增加了总精炼时间；精炼效果受限于预训练IDM的能力；对极端外推角度的改善有限；需要多次迭代才能收敛。
\textbf{方法信息：} Huawei 参与了这篇论文；代码：\url{https://github.com/hyzhou404/FreeFix}。



% METHOD: 2606.22091
\paragraph{ACEsplat（RGB 与位姿下的无外部几何初始化）（2026.06）}
ACEsplat~\cite{2606_22091} 关心现场捕获部署：用户有 RGB 图像和相机位姿，例如来自移动端 SLAM 或已有定位系统，但不希望先跑 SfM 点云、监督深度或大型前馈模型。原始 3DGS 在这种情况下缺一个可用点云；Alternatives to SfM Initialization 讨论了随机初始化和 NeRF 引导初始化，ACEsplat 则给出一条更直接的内部几何路线：先在当前场景自监督学习 scene coordinate regression，再用这个场景坐标先验初始化高斯。

第一阶段沿用 ACE 风格的 SCR。给定多张 RGB 和已知内外参，网络预测每个像素对应的 3D scene coordinate，通过多视图重投影误差自监督训练；无效深度或误差过大的预测被过滤，训练过程中阈值随进度调整。这个 SCR 网络的输出不是最终渲染结果，而是一个 scene-specific internal geometry prior：它给每个像素提供粗场景坐标，也提供与几何一致的中间特征。第二阶段冻结 SCR backbone，用 dense feature upsampling adapter 把低分辨率 SCR 特征提升到图像分辨率，再由轻量 Gaussian initialization head 预测高斯的非中心属性，例如不透明度、协方差和颜色。中心由 SCR scene coordinates 提供，属性由高斯初始化头补齐，随后进入普通 per-scene 3DGS 优化。

相对 InstantSplat，ACEsplat 不依赖 DUSt3R/MASt3R 直接给出点图和相机，而是在已知 pose 条件下快速训练一个当前场景的 SCR 几何先验；相对纯随机初始化，它把高斯中心和特征放回像素对齐的几何框架中。论文报告 SCR mapping 加后续 3DGS 优化可在单 GPU 上约 15--25 分钟完成，用于机器人和混合现实这类需要现场快速建场的场景。
\textbf{实现细节：} ACEsplat。ACEsplat（2026.06） 延续 3DGS 在"无外部几何初始化"上的基本设定，主要改变是：ACEsplat：自监督scene coordinate regression构建内部几何先验，再进行加速3D Gaussian scene regression。3DGS 的显式高斯表示和可微渲染能力保持为基础。其设计动机是：提出无需预计算SfM或深度的两阶段快速3DGS优化流程，提升现场捕获部署可靠性。具体而言，第一阶段SCR模块从RGB和位姿自监督学习场景坐标并建立内部几何先验；第二阶段用该先验加速高斯场优化。适用边界是：仍需已知相机位姿；SCR先验失败会影响后续高斯优化。
\textbf{方法信息：} ByteDance 参与了这篇论文。



% METHOD: 2606.30545
\paragraph{StereoGS（虚拟双目先验）（2026.06）}
StereoGS~\cite{2606_30545} 针对单目深度先验在少视角 3DGS 中的尺度歧义。Depth-Aware 3DGS、DNGaussian 等方法可以用单目深度约束渲染深度，但单目深度常只有相对尺度，不同视图之间不一定一致；在 3-view 或 6-view 训练中，这会让高斯贴到错误深度或形成 floaters。StereoGS 的改变，是把“单张图的深度先验”换成“虚拟双目对 + foundation stereo”的绝对尺度约束。

它包含三个训练期模块。第一是 consistency-aware dense initialization：用 MVSAnywhere 这类 zero-shot 多视图深度模型为每个训练视角估计深度，并通过跨视图重投影过滤异常值，融合成比 SfM 更密、更完整的初始点云。第二是 stereo depth regularization：训练过程中选取一个真实左视图和由当前 3DGS 渲染出的右视图，构造 virtual stereo pair；FoundationStereo 输出左视图 disparity/depth，系统再用遮挡、背景和异常 disparity mask 得到有效区域，只在可靠像素上对渲染深度施加约束。这样监督既利用了真实图像的纹理，又把当前高斯几何拉向双目一致的绝对尺度。第三是 gradient-aware opacity decay：不再周期性粗暴 reset opacity，而是根据高斯相对 opacity gradient 动态衰减冗余基元，减少 sparse-view 过拟合和噪声高斯。

相对 MVPGS 或 Binocular3DGS 的 MVS/双目思想，StereoGS 的关键在于把 dense initialization、online stereo depth regularization 和 opacity decay 放成一套训练期几何约束；推理阶段不需要额外 stereo 网络，仍按普通 3DGS 渲染。需要注意同名仓库很多，本书对应的是 ECCV 2026 项目页和 \texttt{StringerYwh00/StereoGS}，不是 BMVC 2024 的 \texttt{sadrasafa/StereoGS}。
\textbf{实现细节：} StereoGS。StereoGS（2026.06） 延续 3DGS 在"Stereo priors for sparse-view Gaussian Splatting"上的基本设定，主要改变是：构造虚拟 stereo pairs 并使用 stereo depth regularization 建立双目一致性。3DGS 的显式高斯表示和可微渲染能力保持为基础。其设计动机是：在优化中构造虚拟双目对，引入更可靠的 binocular consistency，提高稀疏输入下的几何约束。适用边界是：立体先验质量和虚拟视角构造会影响效果；极少视角和弱纹理区域仍可能不稳。
\textbf{方法信息：} 代码：\url{https://github.com/StringerYwh00/StereoGS}；GitHub 约 6 stars。
\textbf{设计动机与边界：} 相对原始 3DGS 以独立高斯和图像重建误差为主的设定，StereoGS 把问题落在以下可实现的改动上：相比 monocular depth priors，StereoGS 的 stereo depth regularization 提供跨视图尺度一致的深度监督，减少 sparse-view 3DGS 的 defective geometry。其设计目的不是单独增加一个模块，而是用上述表示、监督或渲染约束消除原流程在目标任务中的具体失配。在优化中构造虚拟双目对，引入更可靠的 binocular consistency，提高稀疏输入下的几何约束。适用时还应保留论文给出的边界：立体先验质量和虚拟视角构造会影响效果；极少视角和弱纹理区域仍可能不稳。
% METHOD: 2606.31556
\paragraph{AugSplat（辐射场先验增强 3DGS）（2026.06）}
AugSplat~\cite{2606_31556} 是“用 NeRF/Radiance Field 给 3DGS 补训练视角”的路线。少视角下，NeRF 类连续体表示通常比显式高斯更能在弱几何区域内保持平滑，而 3DGS 的优势是最终实时渲染。AugSplat 不把 radiance field 当作最终表示，也不直接从 radiance field 导出高斯；它把 radiance field 当成合成监督源，先让它生成附近新视角，再用这些视角稳定 3DGS 的早期优化。

流程从稀疏、已标定图像开始。代码中先用 VGGT 准备 COLMAP 风格的相机、点云和深度，再用 Nerfstudio 训练一个 depth-nerfacto radiance-field ensemble。ensemble 在附近视角渲染合成图像，同时给出每个像素的不确定性：如果多个 radiance field 成员对同一像素颜色方差大，说明这个合成监督不可靠，后续损失权重就降低。然后进入 gsplat/3DGS 训练，合成视角与真实输入共同监督高斯。

AugSplat 给了两种接法。Staged AugSplat 先用合成视角做短 warm-up，让粗高斯先学到更完整的几何和低频外观，然后切换回真实图像监督做最终细化；Dual AugSplat 则每步同时采真实图和合成图，合成损失权重按指数衰减，早期帮助几何稳定，后期避免 NeRF artifact 主导外观。相对 MVPGS/StereoGS 的几何先验，AugSplat 借的是 radiance field 的连续渲染能力；相对 RadSplat 的表示级初始化，它更像数据增强，把额外视角作为训练监督，而不是改变高斯参数化。
\textbf{实现细节：} AugSplat。AugSplat（2026.06） 延续 3DGS 在"Radiance-field priors for sparse-view Gaussian Splatting"上的基本设定，主要改变是：Radiance Field-informed Gaussian Splatting，用 NeRF/RF 信息增强稀疏视图 3DGS 初始化和优化。3DGS 的显式高斯表示和可微渲染能力保持为基础。其设计动机是：结合 NeRF 类方法在少视角鲁棒重建上的优势和 Gaussian Splatting 的实时渲染优势。具体而言，在稀疏输入下，利用 radiance field 提供更稳健的几何或外观线索，缓解不完整、噪声点云初始化导致的 Gaussian 优化失败。适用边界是：可能引入额外 radiance field 训练成本；稀疏极限和复杂遮挡场景仍需要验证。
\textbf{技术实现补充：} AugSplat。AugSplat（2026.06） 延续 3DGS 在"Radiance-field priors for sparse-view Gaussian Splatting"上的基本设定，主要改变是：Radiance Field-informed Gaussian Splatting，用 NeRF/RF 信息增强稀疏视图 3DGS 初始化和优化。3DGS 的显式高斯表示和可微渲染能力保持为基础。其设计动机是：结合 NeRF 类方法在少视角鲁棒重建上的优势和 Gaussian Splatting 的实时渲染优势。具体而言，在稀疏输入下，利用 radiance field 提供更稳健的几何或外观线索，缓解不完整、噪声点云初始化导致的 Gaussian 优化失败。适用边界是：可能引入额外 radiance field 训练成本；稀疏极限和复杂遮挡场景仍需要验证。
\textbf{方法信息：} 代码：\url{https://github.com/llazzaroni/AugSplat}；GitHub 约 6 stars。

\begin{takeaway}
本章方法可以按“补什么信号”来记。FSGS、CoherentGS、MVPGS、NexusGS、Sparse2DGS、StereoGS 补几何；DropGaussian、Co-Adaptation、DropoutGS、SE-GS 补结构正则；GSFixer、3DGS-Enhancer、FreeFix、ArtifactWorld、FlowR、Splatent 补生成式新视角先验；InstantSplat、LongSplat、SPFSplat、FlowMap 和 InstantSfM 补位姿与初始化；Deblurring 3DGS、E2GS、EvaGaussians、GeMS、Naka-GS、CompSplat 补真实成像过程；GauSS-MI、GS-Planner、ActiveGS 等让 3DGS 主动决定下一视角；SpotlessSplats、PDF-GS、DualSplat 则把动态干扰物排除在静态场景之外。稀疏和鲁棒 3DGS 的核心不再只是“多加正则”，而是识别当前缺少的是几何、位姿、外观、成像、语义还是采集策略。
\end{takeaway}
\textbf{设计动机与边界：} 相对原始 3DGS 以独立高斯和图像重建误差为主的设定，AugSplat 把问题落在以下可实现的改动上：在稀疏输入下，利用 radiance field 提供更稳健的几何或外观线索，缓解不完整、噪声点云初始化导致的 Gaussian 优化失败。其设计目的不是单独增加一个模块，而是用上述表示、监督或渲染约束消除原流程在目标任务中的具体失配。结合 NeRF 类方法在少视角鲁棒重建上的优势和 Gaussian Splatting 的实时渲染优势。适用时还应保留论文给出的边界：可能引入额外 radiance field 训练成本；稀疏极限和复杂遮挡场景仍需要验证。
